\input{preamble}

% OK, start here.
%
\begin{document}

\title{Groupes de Chow des espaces}

\maketitle

\phantomsection
\label{section-phantom}


\tableofcontents


\section{Introduction}
\label{section-introduction}

\noindent
Dans ce chapitre, nous étudions d'abord les groupes de Chow des espaces
algébriques. Après les avoir définis, nous définissons les classes de Chern des
fibrés vectoriels comme des opérateurs sur ces groupes de Chow. La stratégie
est exactement la même que dans le cas des schémas. Nous invitons donc le
lecteur à consulter l'introduction du chapitre correspondant pour les schémas
(Homologie de Chow, section \ref{chow-section-introduction}).

\medskip\noindent
Quelques articles sur le sujet : \cite{edidin-graham} et \cite{kresch_cycle}.



\section{Cadre}
\label{section-setup}

\noindent
Nous fixons d'abord la catégorie des espaces algébriques que nous considérerons.
Il faut garder à l'esprit, tout au long de ce chapitre, que
« décent $+$ localement noethérien » équivaut à « quasi-séparé $+$ localement
noethérien » d'après Espaces décents, lemme
\ref{decent-spaces-lemma-locally-Noetherian-decent-quasi-separated}.

\begin{situation}
\label{situation-setup}
Ici, $S$ est un schéma et $B$ un espace algébrique sur $S$. Nous supposons
que $B$ est quasi-séparé, localement noethérien et universellement caténaire
(Espaces décents, définition
\ref{decent-spaces-definition-universally-catenary}). De plus, nous supposons
donnée une fonction de dimension $\delta : |B| \longrightarrow
\mathbf{Z}$. On dit que $X/B$ est {\it bon} si $X$ est un espace algébrique
sur $B$ dont le morphisme de structure $f : X \to B$ est quasi-séparé et
localement de type fini. Dans ce cas, nous définissons
$$
\delta = \delta_{X/B} : |X| \longrightarrow \mathbf{Z}
$$
comme l'application qui à $x$ associe $\delta(f(x))$ plus le degré de
transcendance de $x/f(x)$ (Morphismes d'espaces, définition
\ref{spaces-morphisms-definition-dimension-fibre}). Il s'agit d'une fonction
de dimension d'après Compléments sur les morphismes d'espaces, lemme
\ref{spaces-more-morphisms-lemma-universally-catenary-dimension-function}.
\end{situation}

\noindent
Un cas particulier est celui où $S = B$ est un schéma et $(S, \delta)$ vérifie
les hypothèses d'Homologie de Chow, situation \ref{chow-situation-setup}.
Ainsi, $B$ peut être le spectre d'un corps (Homologie de Chow, exemple
\ref{chow-example-field}) ou $B = \Spec(\mathbf{Z})$ (Homologie de Chow, exemple
\ref{chow-example-domain-dimension-1}).

\medskip\noindent
De nombreux lemmes, propositions, théorèmes et définitions concernant les
espaces algébriques sont plus faciles à établir dans le cadre de la situation
\ref{situation-setup} car les espaces algébriques avec lesquels nous
travaillons sont quasi-séparés (et donc a fortiori décents) et localement
noethériens. Nous le signalerons au fil du chapitre par des remarques comme
les suivantes.

\begin{remark}
\label{remark-sober}
Dans la situation \ref{situation-setup}, si $X/B$ est bon, alors $|X|$ est un
espace topologique sobre. Voir Propriétés des espaces, lemme
\ref{spaces-properties-lemma-quasi-separated-sober} ou Espaces décents,
proposition \ref{decent-spaces-proposition-reasonable-sober}. Nous
l'utiliserons sans autre mention pour choisir des points génériques de
sous-ensembles fermés irréductibles de $|X|$.
\end{remark}

\begin{remark}
\label{remark-integral}
Dans la situation \ref{situation-setup}, si $X/B$ est bon, alors $X$ est
intègre (Espaces sur les corps, définition
\ref{spaces-over-fields-definition-integral-algebraic-space}) si et seulement
si $X$ est réduit et $|X|$ est irréductible. De plus, pour tout point $\xi \in
|X|$, il existe un sous-espace fermé intègre unique $Z \subset X$ tel que
$\xi$ soit le point générique du sous-ensemble fermé $|Z| \subset |X|$ ; voir
Espaces sur les corps, lemme
\ref{spaces-over-fields-lemma-decent-irreducible-closed}.
\end{remark}

\noindent
Si $B$ est jacobsonien et que $\delta$ envoie les points fermés sur zéro, alors
$\delta$ est la fonction qui envoie un point à la dimension de sa fermeture.

\begin{lemma}
\label{lemma-delta-is-dimension}
Dans la situation \ref{situation-setup}, supposons que $B$ soit jacobsonien et que
$\delta(b) = 0$ pour chaque point fermé $b$ de $|B|$. Soit $X/B$ bon. Si
$Z \subset X$ est un sous-espace fermé intègre de point générique $\xi \in
|Z|$, alors les entiers suivants sont égaux :
\begin{enumerate}
\item $\delta(\xi) = \delta_{X/B}(\xi)$,
\item $\dim(|Z|)$,
\item $\text{codim}(\{z\}, |Z|)$ pour $z \in |Z|$ fermé,
\item la dimension de l'anneau local de $Z$ à $z$ pour $z \in |Z|$ fermé, et
\item $\dim(\mathcal{O}_{Z, \overline{z}})$ pour $z \in |Z|$ fermé.
\end{enumerate}
\end{lemma}

\begin{proof}
Soient $X$, $Z$, $\xi$ comme dans le lemme. Puisque $X$ est localement de type
fini sur $B$, l'espace $X$ est jacobsonien ; voir Espaces décents, lemme
\ref{decent-spaces-lemma-Jacobson-universally-Jacobson}. Ainsi
$X_{\text{ft-pts}} \subset |X|$ est l'ensemble des points fermés d'après
Espaces décents, lemme \ref{decent-spaces-lemma-decent-Jacobson-ft-pts}. Étant donné
une chaîne $T_0 \supset \ldots \supset T_e$ de sous-ensembles fermés
irréductibles de $|Z|$, l'ensemble $T_e \cap X_{\text{ft-pts}}$ est non vide
d'après Morphismes d'espaces, lemme
\ref{spaces-morphisms-lemma-enough-finite-type-points}. Ainsi on peut toujours
supposer qu'une telle chaîne se termine par $T_e = \{z\}$ pour un point fermé
$z \in |Z|$. Il s'ensuit que $\dim(Z) = \sup_z \text{codim}(\{z\}, |Z|)$, où
$z$ court sur les points fermés de $|Z|$. Nous avons $\text{codim}(\{z\}, Z) =
\delta(\xi) - \delta(z)$ d'après Topologie, lemme
\ref{topology-lemma-dimension-function-catenary}. Par Morphismes d'espaces,
lemme \ref{spaces-morphisms-lemma-finite-type-points-morphism}, l'image de $z$
est un point de type fini de $B$, c'est-à-dire un point fermé de $|B|$. Par
Morphismes d'espaces, lemme
\ref{spaces-morphisms-lemma-jacobson-finite-type-points}, le degré de
transcendance de $z/b$ est $0$. Par hypothèse, nous concluons que
$\delta(z) = \delta(b) = 0$. On obtient donc l'égalité
$$
\dim(|Z|) = \text{codim}(\{z\}, Z) = \delta(\xi)
$$
pour tout point fermé $z \in |Z|$. Enfin, $\text{codim}(\{z\}, Z)$ est égal à
la dimension de l'anneau local de $Z$ en $z$ d'après Espaces décents, lemme
\ref{decent-spaces-lemma-codimension-local-ring}, laquelle est à son tour égale
à $\dim(\mathcal{O}_{Z, \overline{z}})$ d'après Propriétés des espaces, lemme
\ref{spaces-properties-lemma-dimension-local-ring}.
\end{proof}

\noindent
Dans la situation du lemme précédent, la valeur de $\delta$ au point
générique d'un sous-ensemble fermé irréductible est la dimension de ce
sous-ensemble. Cela motive la définition suivante.

\begin{definition}
\label{definition-delta-dimension}
Dans la situation \ref{situation-setup}, pour tout $X/B$ bon et tout
sous-ensemble fermé irréductible $T \subset |X|$, on définit
$$
\dim_\delta(T) = \delta(\xi)
$$
où $\xi \in T$ est le point générique de $T$. Nous appelons ce nombre la {\it
$\delta$-dimension de $T$}. Si $T \subset |X|$ est un sous-ensemble fermé,
alors nous définissons $\dim_\delta(T)$ comme la borne supérieure des
$\delta$-dimensions des composantes irréductibles de $T$. Si $Z$ est un sous-espace fermé
de $X$, alors nous définissons $\dim_\delta(Z) = \dim_\delta(|Z|)$.
\end{definition}

\noindent
Bien sûr, cela signifie simplement que $\dim_\delta(T) = \sup \{\delta(t) \mid
t \in T\}$.







\section{Cycles}
\label{section-cycles}

\noindent
Cette section est l'analogue d'Homologie de Chow, section
\ref{chow-section-cycles}.

\medskip\noindent
Puisque nous ne supposons pas nos espaces quasi-compacts, il faut être un peu
prudent dans la définition des cycles. Nous devons autoriser les sommes infinies,
car une fonction rationnelle peut, par exemple, avoir une infinité de pôles.
En revanche, si $X$ est quasi-compact, alors un cycle est
une somme finie comme d'habitude.

\begin{definition}
\label{definition-cycles}
Dans la situation \ref{situation-setup}, soit $X/B$ bon. Soit $k
\in \mathbf{Z}$.
\begin{enumerate}
\item Un {\it cycle sur $X$} est une somme formelle
$$
\alpha = \sum n_Z [Z]
$$
où la somme s'étend sur les sous-espaces fermés intègres $Z \subset X$,
chaque $n_Z \in \mathbf{Z}$ et $\{|Z|; n_Z \not = 0\}$ est une famille
localement finie de sous-ensembles de $|X|$ (Topologie, définition
\ref{topology-definition-locally-finite}).
\item Un {\it $k$-cycle} sur $X$ est un cycle
$$
\alpha = \sum n_Z [Z]
$$
où $n_Z \not = 0 \Rightarrow \dim_\delta(Z) = k$.
\item Le groupe abélien de tous les $k$-cycles sur $X$ est noté $Z_k(X)$.
\end{enumerate}
\end{definition}

\noindent
En d'autres termes, un $k$-cycle sur $X$ est une combinaison formelle
$\mathbf{Z}$-linéaire localement finie de sous-espaces fermés intègres (remarque
\ref{remark-integral}) de $\delta$-dimension $k$. L'addition des $k$-cycles
$\alpha = \sum n_Z[Z]$ et $\beta = \sum m_Z[Z]$ est donné par
$$
\alpha + \beta = \sum (n_Z + m_Z)[Z],
$$
c'est-à-dire en additionnant les coefficients.




\section{Multiplicités}
\label{section-multiplicities}

\noindent
Cette section rassemble quelques résultats simples sur les longueurs et les
multiplicités.

\begin{lemma}
\label{lemma-length}
Soit $S$ un schéma et soit $X$ un espace algébrique sur $S$. Soit $\mathcal{F}$
un $\mathcal{O}_X$-module quasi-cohérent. Soit $x \in |X|$. Soit $d \in
\{0, 1, 2, \ldots, \infty\}$. Les assertions suivantes sont équivalentes :
\begin{enumerate}
\item
$\text{length}_{\mathcal{O}_{X, \overline{x}}} \mathcal{F}_{\overline{x}} = d$
\item il existe un morphisme \'etale $U \to X$, avec $U$ un schéma, et un
point $u \in U$ d'image $x$, tels que $\text{length}_{\mathcal{O}_{U, u}}
(\mathcal{F}|_U)_u = d$
\item pour tout morphisme \'etale $U \to X$, avec $U$ un schéma, et tout
point $u \in U$ d'image $x$, on a $\text{length}_{\mathcal{O}_{U, u}}
(\mathcal{F}|_U)_u = d$
\end{enumerate}
\end{lemma}

\begin{proof}
Soient $U \to X$ et $u \in U$ comme en (2) ou (3). On sait alors que
$\mathcal{O}_{X, \overline{x}}$ est l'hensélisation stricte de
$\mathcal{O}_{U, u}$ et que
$$
\mathcal{F}_{\overline{x}} =
(\mathcal{F}|_U)_u \otimes_{\mathcal{O}_{U, u}} \mathcal{O}_{X, \overline{x}}
$$
Voir Propriétés des espaces, lemmes
\ref{spaces-properties-lemma-describe-etale-local-ring} et
\ref{spaces-properties-lemma-stalk-quasi-coherent}. Il s'ensuit, par Algèbre,
lemme \ref{algebra-lemma-pullback-module}, par la platitude de
$\mathcal{O}_{U, u} \to \mathcal{O}_{X, \overline{x}}$ et par le fait que
$\mathcal{O}_{X, \overline{x}}/\mathfrak m_u\mathcal{O}_{X,
\overline{x}}$ est le corps résiduel de $\mathcal{O}_{X, \overline{x}}$, que
les longueurs sont égales. Ces propriétés des hensélisations strictes se
trouvent dans Compléments d'algèbre, lemme
\ref{more-algebra-lemma-dumb-properties-henselization}.
\end{proof}

\begin{definition}
\label{definition-length-at-x}
Soit $S$ un schéma et soit $X$ un espace algébrique sur $S$. Soit $\mathcal{F}$
un $\mathcal{O}_X$-module quasi-cohérent. Soit $x \in |X|$. Soit $d \in
\{0, 1, 2, \ldots, \infty\}$. Nous disons que {\it $\mathcal{F}$ est de longueur
$d$ en $x$} si les conditions équivalentes du lemme \ref{lemma-length} sont
satisfaites.
\end{definition}

\begin{lemma}
\label{lemma-length-closed-immersion}
Soit $S$ un schéma. Soit $i : Y \to X$ une immersion fermée d'espaces
algébriques sur $S$. Soit $\mathcal{G}$ un $\mathcal{O}_Y$-module
quasi-cohérent. Soit $y \in |Y|$ d'image $x \in |X|$. Soit $d \in \{0,
1, 2, \ldots, \infty\}$. Les assertions suivantes sont équivalentes :
\begin{enumerate}
\item $\mathcal{G}$ a une longueur $d$ à $y$, et
\item $i_*\mathcal{G}$ a une longueur $d$ à $x$.
\end{enumerate}
\end{lemma}

\begin{proof}
Choisissons un morphisme \'etale $f : U \to X$, où $U$ est un schéma, et un
point $u \in U$ d'image $x$. Posons $V = Y \times_X U$. Notons $g : V \to Y$
et $j : V \to U$ les projections. Alors $j : V \to U$ est une immersion fermée
et il existe un unique point $v \in V$ d'images $y \in |Y|$ et $u \in U$
(voir Propriétés des espaces, lemme
\ref{spaces-properties-lemma-points-cartesian}, et Espaces, lemme
\ref{spaces-lemma-base-change-immersions}). Nous avons $j_*(\mathcal{G}|_V) =
(i_*\mathcal{G})|_U$ en tant que modules sur le schéma $V$, et $j_*$ est
l'image directe « usuelle » des modules pour le morphisme de schémas $j$ ;
voir la discussion qui entoure Cohomologie des espaces, équation
(\ref{spaces-cohomology-equation-representable-higher-direct-image}). On se
ramène ainsi au cas des schémas : si $i : Y \to X$ est une immersion fermée de
schémas, alors
$$
(i_*\mathcal{G})_x = \mathcal{G}_y
$$
en tant que modules sur $\mathcal{O}_{X, x}$, la structure de module au membre
de droite étant donnée par la surjection $i_y^\sharp : \mathcal{O}_{X, x} \to
\mathcal{O}_{Y, y}$. L'égalité résulte donc d'Algèbre, lemme
\ref{algebra-lemma-length-independent}.
\end{proof}

\begin{lemma}
\label{lemma-length-finite}
Soit $S$ un schéma et soit $X$ un espace algébrique localement noethérien sur
$S$. Soit $\mathcal{F}$ un $\mathcal{O}_X$-module cohérent. Soit $x \in |X|$.
Les assertions suivantes sont équivalentes :
\begin{enumerate}
\item il existe un morphisme \'etale $U \to X$, avec $U$ un schéma, et un point
$u \in U$ d'image $x$, tel que $u$ soit un point générique d'une composante
irréductible de $\text{Supp}(\mathcal{F}|_U)$,
\item pour tout morphisme \'etale $U \to X$, avec $U$ un schéma, et tout point
$u \in U$ d'image $x$, le point $u$ est générique dans une composante irréductible
de $\text{Supp}(\mathcal{F}|_U)$,
\item la longueur de $\mathcal{F}$ à $x$ est finie et non nulle.
\end{enumerate}
Si $X$ est décent (de manière équivalente, quasi-séparé), ces assertions sont
également équivalentes à
\begin{enumerate}
\item[(4)] $x$ est un point générique d'une composante irréductible de
$\text{Supp}(\mathcal{F})$.
\end{enumerate}
\end{lemma}

\begin{proof}
Supposons que $f : U \to X$ soit un morphisme \'etale, avec $U$ un schéma, et
que $u \in U$ ait pour image $x$. Alors $\mathcal{F}|_U = f^*\mathcal{F}$ est
un $\mathcal{O}_U$-module cohérent sur le schéma localement noethérien $U$ ; en
particulier, $(\mathcal{F}|_U)_u$ est un $\mathcal{O}_{U, u}$-module fini ; voir
Cohomologie des espaces, lemme
\ref{spaces-cohomology-lemma-coherent-Noetherian} et Cohomologie des schémas,
lemme \ref{coherent-lemma-coherent-Noetherian}. Rappelons que le support de
$\mathcal{F}|_U$ est un sous-ensemble fermé de $U$ (Morphismes, lemme
\ref{morphisms-lemma-support-finite-type}) et que le support de
$(\mathcal{F}|_U)_u$ est l'image réciproque du support de $\mathcal{F}|_U$ par le
morphisme $\Spec(\mathcal{O}_{U, u}) \to U$. Ainsi $u$ est un point générique
d'une composante irréductible de $\text{Supp}(\mathcal{F}|_U)$ si et seulement
si le support de $(\mathcal{F}|_U)_u$ est égal à l'idéal maximal de
$\mathcal{O}_{U, u}$. L'équivalence de (1), (2) et (3) résulte alors d'Algèbre,
lemme \ref{algebra-lemma-support-point}.

\medskip\noindent
Si $X$ est décent, choisissons un morphisme \'etale $f : U \to X$ et un point
$u \in U$ d'image $x$. L'image réciproque du support de $\mathcal{F}$ est le
support de $\mathcal{F}|_U$ ; voir Morphismes d'espaces, lemme
\ref{spaces-morphisms-lemma-support-finite-type}. De plus, les spécialisations
$x' \leadsto x$ dans $|X|$ se relèvent en des spécialisations $u' \leadsto u$
dans $U$ et toute spécialisation non triviale $u' \leadsto u$ dans $U$
correspond à une spécialisation non triviale $f(u') \leadsto f(u)$ dans $|X|$,
voir Espaces décents, lemmes \ref{decent-spaces-lemma-decent-specialization}
et \ref{decent-spaces-lemma-decent-no-specializations-map-to-same-point}. En
utilisant le fait que $|X|$ et $U$ sont des espaces topologiques sobres
(Espaces décents, proposition \ref{decent-spaces-proposition-reasonable-sober}
et Schémas, lemme \ref{schemes-lemma-scheme-sober}), nous concluons que $x$
est un point générique du support de $\mathcal{F}$ si et seulement si $u$ est
un point générique du support de $\mathcal{F}|_U$. Nous concluons que (4) est
équivalent à (1).

\medskip\noindent
L'assertion entre parenthèses résulte d'Espaces décents, lemme
\ref{decent-spaces-lemma-locally-Noetherian-decent-quasi-separated}.
\end{proof}

\begin{lemma}
\label{lemma-point-of-max-dimension}
Dans la situation \ref{situation-setup}, soit $X/B$ bon. Soit $T
\subset |X|$ un sous-ensemble fermé et $t \in T$. Si $\dim_\delta(T) \leq k$
et $\delta(t) = k$, alors $t$ est un point générique d'une composante
irréductible de $T$.
\end{lemma}

\begin{proof}
Nous savons que $t$ appartient à une composante irréductible $T' \subset T$.
Soit $t' \in T'$ le point générique. Alors $k \geq \delta(t') \geq \delta(t)$.
Puisque $\delta$ est une fonction de dimension, nous voyons que $t = t'$.
\end{proof}



\section{Cycle associé à un sous-espace fermé}
\label{section-cycle-of-closed-subscheme}

\noindent
Cette section est l'analogue d'Homologie de Chow, section
\ref{chow-section-cycle-of-closed-subscheme}.

\begin{remark}
\label{remark-irreducible-component}
Dans la situation \ref{situation-setup}, soit $X/B$ bon. Soit $Y
\subset X$ un sous-espace fermé. D'après les remarques \ref{remark-sober} et
\ref{remark-integral}, il existe des bijections (c'est-à-dire des
correspondances $1$-à-$1$) entre les ensembles suivants :
\begin{enumerate}
\item les composantes irréductibles $T$ de $|Y|$,
\item les points génériques des composantes irréductibles de $|Y|$, et
\item les sous-espaces fermés intègres $Z \subset Y$ tels que $|Z|$ soit une
composante irréductible de $|Y|$.
\end{enumerate}
Dans ce chapitre, nous appelons tout $Z$ comme en (3) une {\it composante
irréductible de $Y$}, et nous appelons $\xi \in |Z|$ son {\it point générique}.
\end{remark}

\begin{definition}
\label{definition-cycle-associated-to-closed-subscheme}
Dans la situation \ref{situation-setup}, soit $X/B$ bon. Soit $Y
\subset X$ un sous-espace fermé.
\begin{enumerate}
\item Pour une composante irréductible $Z \subset Y$ de point générique $\xi$,
la longueur de $\mathcal{O}_Y$ en $\xi$ (définition
\ref{definition-length-at-x}) est appelée la {\it multiplicité de $Z$ dans
$Y$}. Par le lemme \ref{lemma-length-finite} appliqué à $\mathcal{O}_Y$ sur
$Y$, c'est un entier positif.
\item Supposons $\dim_\delta(Y) \leq k$. Le {\it $k$-cycle associé à $Y$} est
$$
[Y]_k = \sum m_{Z, Y}[Z]
$$
où la somme porte sur les composantes irréductibles $Z$ de $Y$ de
$\delta$-dimension $k$ et $m_{Z, Y}$ est la multiplicité de $Z$ dans $Y$. Il
s'agit d'un $k$-cycle d'après Espaces sur les corps, lemme
\ref{spaces-over-fields-lemma-components-locally-finite}.
\end{enumerate}
\end{definition}

\noindent
Il importe de noter que nous ne définissons $[Y]_k$ que si la
$\delta$-dimension de $Y$ ne dépasse pas $k$. En d’autres termes, par
convention, l'écriture $[Y]_k$ implique donc que $\dim_\delta(Y) \leq
k$.







\section{Cycle associé à un faisceau cohérent}
\label{section-cycle-of-coherent-sheaf}

\noindent
Cette section est l'analogue d'Homologie de Chow, section
\ref{chow-section-cycle-of-coherent-sheaf}.

\begin{definition}
\label{definition-cycle-associated-to-coherent-sheaf}
Dans la situation \ref{situation-setup}, soit $X/B$ bon. Soient
$\mathcal{F}$ un $\mathcal{O}_X$-module cohérent.
\begin{enumerate}
\item Pour un sous-espace fermé intègre $Z \subset X$ de point générique
$\xi$ tel que $|Z|$ soit une composante irréductible de
$\text{Supp}(\mathcal{F})$, la longueur de $\mathcal{F}$ en $\xi$ (définition
\ref{definition-length-at-x}) est appelée la {\it multiplicité de $Z$ dans
$\mathcal{F}$}. D'après le lemme \ref{lemma-length-finite}, il s'agit d'un
entier positif.
\item Supposons $\dim_\delta(\text{Supp}(\mathcal{F})) \leq k$. Le {\it
$k$-cycle associé à $\mathcal{F}$} est
$$
[\mathcal{F}]_k = \sum m_{Z, \mathcal{F}}[Z]
$$
où la somme porte sur les sous-espaces fermés intègres $Z \subset X$
correspondant aux composantes irréductibles de $\text{Supp}(\mathcal{F})$ de
$\delta$-dimension $k$, et $m_{Z, \mathcal{F}}$ est la multiplicité de $Z$
dans $\mathcal{F}$. Il s'agit d'un $k$-cycle d'après Espaces sur les corps, lemme
\ref{spaces-over-fields-lemma-components-locally-finite}.
\end{enumerate}
\end{definition}

\noindent
Il importe de noter que nous ne définissons $[\mathcal{F}]_k$ que si
$\mathcal{F}$ est cohérent et que la $\delta$-dimension de
$\text{Supp}(\mathcal{F})$ ne dépasse pas $k$. En d’autres termes, par
convention, l'écriture $[\mathcal{F}]_k$ implique donc que
$\mathcal{F}$ est cohérent sur $X$ et $\dim_\delta(\text{Supp}(\mathcal{F}))
\leq k$.

\begin{lemma}
\label{lemma-reformulate-coeff-coherent}
Dans la situation \ref{situation-setup}, soit $X/B$ bon. Soit
$\mathcal{F}$ un $\mathcal{O}_X$-module cohérent avec
$\dim_\delta(\text{Supp}(\mathcal{F})) \leq k$. Soit $Z$ un sous-espace fermé
intègre de $X$ tel que $\dim_\delta(Z) = k$. Soit $\xi \in |Z|$ son point
générique. Alors le coefficient de $Z$ dans $[\mathcal{F}]_k$ est la
longueur de $\mathcal{F}$ en $\xi$.
\end{lemma}

\begin{proof}
Observons que $|Z|$ est une composante irréductible de
$\text{Supp}(\mathcal{F})$ si et seulement si $\xi \in
\text{Supp}(\mathcal{F})$ ; voir le lemme \ref{lemma-point-of-max-dimension}.
De plus, la longueur de $\mathcal{F}$ en $\xi$ est nulle si $\xi \not \in
\text{Supp}(\mathcal{F})$. Le résultat découle alors de la définition
\ref{definition-cycle-associated-to-coherent-sheaf}.
\end{proof}

\begin{lemma}
\label{lemma-cycle-closed-coherent}
Dans la situation \ref{situation-setup}, soit $X/B$ bon. Soit $Y
\subset X$ un sous-espace fermé. Si $\dim_\delta(Y) \leq k$, alors $[Y]_k =
[i_*\mathcal{O}_Y]_k$ où $i : Y \to X$ est le morphisme d'inclusion.
\end{lemma}

\begin{proof}
Soit $Z$ un sous-espace fermé intègre de $X$ tel que $\dim_\delta(Z) = k$. Si $Z
\not \subset Y$, le coefficient de $Z$ est nul dans $[Y]_k$ comme dans
$[i_*\mathcal{O}_Y]_k$. Si $Z \subset Y$, alors le point générique de $Z$ peut
être considéré comme un point $y \in |Y|$ d'image $x \in |X|$. Alors le
coefficient de $Z$ dans $[Y]_k$ est la longueur de $\mathcal{O}_Y$ en $y$ et
le coefficient de $Z$ dans $[i_*\mathcal{O}_Y]_k$ est la longueur de
$i_*\mathcal{O}_Y$ en $x$. L'égalité des coefficients résulte donc du lemme
\ref{lemma-length-closed-immersion}.
\end{proof}

\begin{lemma}
\label{lemma-additivity-sheaf-cycle}
Dans la situation \ref{situation-setup}, soit $X/B$ bon. Soit $0 \to
\mathcal{F} \to \mathcal{G} \to \mathcal{H} \to 0$ une suite exacte courte
de $\mathcal{O}_X$-modules cohérents. Supposons que la $\delta$-dimension de
chacun des supports de $\mathcal{F}$, $\mathcal{G}$ et $\mathcal{H}$ soit
$\leq k$. Alors
$[\mathcal{G}]_k = [\mathcal{F}]_k + [\mathcal{H}]_k$.
\end{lemma}

\begin{proof}
Soit $Z$ un sous-espace fermé intègre de $X$ tel que $\dim_\delta(Z) = k$. Il
suffit de montrer que les coefficients de $Z$ dans $[\mathcal{G}]_k$,
$[\mathcal{F}]_k$ et $[\mathcal{H}]_k$ satisfont l'additivité correspondante.
Par le lemme \ref{lemma-reformulate-coeff-coherent}, il suffit de montrer
$$
\text{the length of }\mathcal{G}\text{ at }x =
\text{the length of }\mathcal{F}\text{ at }x +
\text{the length of }\mathcal{H}\text{ at }x
$$
pour tout $x \in |X|$. D'après la définition \ref{definition-length-at-x},
cela résulte immédiatement de l'additivité des longueurs ; voir Algèbre, lemme
\ref{algebra-lemma-length-additive}.
\end{proof}





\section{Préliminaires à l'image directe propre}
\label{section-preparation-pushforward}

\noindent
Cette section est l'analogue de l'Homologie de Chow, section
\ref{chow-section-preparation-pushforward}.

\begin{lemma}
\label{lemma-proper-image}
Dans la situation \ref{situation-setup}, soient $X,Y/B$ bons et soit $f : X
\to Y$ un morphisme au-dessus de $B$. Si $Z \subset X$ est un sous-espace fermé
intègre, alors il existe un unique sous-espace fermé intègre $Z' \subset Y$
tel qu'il existe un diagramme commutatif
$$
\xymatrix{
Z \ar[r] \ar[d] & X \ar[d]^f \\
Z' \ar[r] & Y
}
$$
avec $Z \to Z'$ dominant. Si $f$ est propre, alors $Z \to Z'$ est propre et
surjectif.
\end{lemma}

\begin{proof}
Soit $\xi \in |Z|$ le point générique. Soit $Z' \subset Y$ le sous-espace
fermé intègre dont le point générique est $\xi' = f(\xi)$ ; voir la remarque
\ref{remark-integral}. Puisque $\xi \in |f^{-1}(Z')| = |f|^{-1}(|Z'|)$ d'après
Propriétés des espaces, lemme \ref{spaces-properties-lemma-points-cartesian},
et puisque $Z$ est le sous-espace réduit tel que $|Z| =
\overline{\{\xi\}}$, nous voyons que $Z \subset f^{-1}(Z')$ en tant que
sous-espaces fermés de $X$ (voir Propriétés des espaces, lemme
\ref{spaces-properties-lemma-map-into-reduction}). On obtient ainsi notre
morphisme $Z \to Z'$. Ce morphisme est dominant puisque le point générique de
$Z$ s'envoie sur le point générique de $Z'$. L'unicité de $Z'$ est claire. Si
$f$ est propre, alors $Z \to Y$ est propre en tant que
composition de morphismes propres (Morphismes d'espaces, lemmes
\ref{spaces-morphisms-lemma-base-change-proper} et
\ref{spaces-morphisms-lemma-closed-immersion-proper}). Il en résulte que $Z \to
Z'$ est propre d'après Morphismes d'espaces, lemme
\ref{spaces-morphisms-lemma-universally-closed-permanence}. La surjectivité en
résulte, puisque l'image d'un morphisme propre est fermée.
\end{proof}

\begin{remark}
\label{remark-residue-field}
Dans la situation \ref{situation-setup}, soit $X/B$ bon. Tout $x \in
|X|$ peut être représenté par un monomorphisme (unique) $\Spec(k) \to X$, où
$k$ est un corps ; voir Espaces décents, lemme
\ref{decent-spaces-lemma-decent-points-monomorphism}. Alors $k$ est le {\it
corps résiduel de} $x$ et se note $\kappa(x)$. Rappelons que $X$ possède un
sous-schéma ouvert dense $U \subset X$ (Propriétés des espaces, proposition
\ref{spaces-properties-proposition-locally-quasi-separated-open-dense-scheme}).
Si $x \in U$, alors $\kappa(x)$ coïncide avec le corps résiduel de $x$
dans le schéma $U$. Voir Espaces décents, section
\ref{decent-spaces-section-residue-fields-henselian-local-rings}.
\end{remark}

\begin{remark}
\label{remark-function-field}
Dans la situation \ref{situation-setup}, soit $X/B$ bon. Supposons que
$X$ soit intègre. Dans ce cas, le {\it corps des fonctions} $R(X)$ de $X$ est
défini et est égal au corps résiduel de $X$ en son point générique. Voir
Espaces sur les corps, définition
\ref{spaces-over-fields-definition-function-field}. En combinant cela avec la
remarque \ref{remark-integral}, nous voyons que, pour tout $x \in X$, le
corps résiduel $\kappa(x)$ est le corps des fonctions de l'unique sous-espace fermé
intègre unique $Z \subset X$ dont le point générique est $x$.
\end{remark}

\begin{lemma}
\label{lemma-equal-dimension}
Dans la situation \ref{situation-setup}, soient $X, Y/B$ bons et soit $f : X
\to Y$ un morphisme au-dessus de $B$. Supposons $X$ et $Y$ intègres et
$\dim_\delta(X) = \dim_\delta(Y)$. Alors, ou bien $f$ se factorise par un
sous-espace fermé strict de $Y$, ou bien $f$ est dominant et l'extension de
corps de fonctions $R(X) / R(Y)$ est finie.
\end{lemma}

\begin{proof}
Par le lemme \ref{lemma-proper-image}, il existe un sous-espace fermé intègre
unique $Z \subset Y$ tel que $f$ se factorise par un morphisme dominant $X
\to Z$. Alors $Z = Y$ si et seulement si $\dim_\delta(Z) = \dim_\delta(Y)$.
D'autre part, par notre construction des fonctions de dimension (voir la
situation \ref{situation-setup}), nous avons $\dim_\delta(X) =
\dim_\delta(Z) + r$, où $r$ est le degré de transcendance de l'extension
$R(X)/R(Z)$. Le lemme résulte alors d'Espaces sur les corps, lemme
\ref{spaces-over-fields-lemma-finite-degree}.
\end{proof}

\begin{lemma}
\label{lemma-quasi-compact-locally-finite}
Dans la situation \ref{situation-setup}, soient $X, Y/B$ bons. Soit $f :
X \to Y$ un morphisme sur $B$. Supposons que $f$ soit quasi-compact et que
$\{T_i\}_{i \in I}$ soit une famille localement finie de parties fermées de
$|X|$. Alors $\{\overline{|f|(T_i)}\}_{i \in I}$ est une famille localement
finie de parties fermées de $|Y|$.
\end{lemma}

\begin{proof}
Soit $V \subset |Y|$ une partie ouverte quasi-compacte. Alors $|f|^{-1}(V)
\subset |X|$ est quasi-compact par Morphismes d'espaces, Lemme
\ref{spaces-morphisms-lemma-quasi-compact-is-quasi-compact}. L’ensemble $\{i
\in I : T_i \cap |f|^{-1}(V) \not = \emptyset \}$ est donc fini par un simple
argument topologique que nous omettons. Or cet ensemble est aussi égal à
$$
\{i \in I : |f|(T_i) \cap V \not = \emptyset \} =
\{i \in I : \overline{|f|(T_i)} \cap V \not = \emptyset \}
$$
ce qui démontre le lemme.
\end{proof}










\section{Image directe propre}
\label{section-proper-pushforward}

\noindent
Cette section est l'analogue de l'Homologie de Chow, section
\ref{chow-section-proper-pushforward}.

\begin{definition}
\label{definition-proper-pushforward}
Dans la situation \ref{situation-setup}, soient $X, Y/B$ bons. Soit $f : X
\to Y$ un morphisme au-dessus de $B$. Supposons $f$ propre.
\begin{enumerate}
\item Soit $Z \subset X$ un sous-espace fermé intègre avec $\dim_\delta(Z) =
k$. Soit $Z' \subset Y$ l'image de $Z$ comme dans le lemme
\ref{lemma-proper-image}. Nous définissons
$$
f_*[Z] =
\left\{
\begin{matrix}
0 & \text{if} & \dim_\delta(Z')< k, \\
\deg(Z/Z') [Z'] & \text{if} & \dim_\delta(Z') = k.
\end{matrix}
\right.
$$
Le degré de $Z$ sur $Z'$ est défini et fini si $\dim_\delta(Z') =
\dim_\delta(Z)$ d'après le lemme \ref{lemma-equal-dimension} et Espaces sur
les corps, définition \ref{spaces-over-fields-definition-degree}.
\item Soit $\alpha = \sum n_Z [Z]$ un $k$-cycle sur $X$. L'{\it image directe}
de $\alpha$ est la somme
$$
f_* \alpha = \sum n_Z f_*[Z]
$$
où chaque $f_*[Z]$ est défini comme ci-dessus. Cette somme est localement finie
d'après le lemme \ref{lemma-quasi-compact-locally-finite} ci-dessus.
\end{enumerate}
\end{definition}

\noindent
Par définition, l'image directe propre des cycles
$$
f_* : Z_k(X) \longrightarrow Z_k(Y)
$$
est un homomorphisme de groupes abéliens. Elle fait de $X \mapsto Z_k(X)$ un
foncteur covariant sur la catégorie dont les objets sont les espaces
algébriques bons au-dessus de $B$ et dont les morphismes sont les morphismes
propres au-dessus de $B$.

\begin{lemma}
\label{lemma-compose-pushforward}
Dans la situation \ref{situation-setup}, soient $X, Y, Z/B$ bons. Soient $f
: X \to Y$ et $g : Y \to Z$ des morphismes propres au-dessus de $B$. Alors $g_* \circ
f_* = (g \circ f)_*$ comme applications $Z_k(X) \to Z_k(Z)$.
\end{lemma}

\begin{proof}
Soit $W \subset X$ un sous-espace fermé intègre de dimension $k$. Considérons
les sous-espaces fermés intègres $W' \subset Y$ et $W'' \subset Z$ obtenus en
appliquant le lemme \ref{lemma-proper-image} à $f$ et $W$, puis à $g$ et $W'$.
Alors $W \to W'$ et $W' \to W''$ sont surjectifs et propres. Il faut montrer
que $g_*(f_*[W]) = (f \circ g)_*[W]$. Si $\dim_\delta(W'') <
k$, alors les deux côtés sont nuls. Si $\dim_\delta(W'') = k$, alors nous
voyons que $W \to W'$ et $W' \to W''$ satisfont tous deux aux hypothèses du
lemme \ref{lemma-equal-dimension}. Ainsi
$$
g_*(f_*[W]) = \deg(W/W')\deg(W'/W'')[W''],
\quad
(f \circ g)_*[W] = \deg(W/W'')[W''].
$$
La conclusion résulte alors d'Espaces sur les corps, lemme
\ref{spaces-over-fields-lemma-degree-composition}.
\end{proof}

\begin{lemma}
\label{lemma-cycle-push-sheaf}
Dans la situation \ref{situation-setup}, soit $f : X \to Y$ un morphisme
propre d'espaces algébriques bons au-dessus de $B$.
\begin{enumerate}
\item Soit $Z \subset X$ un sous-espace fermé avec $\dim_\delta(Z) \leq k$.
Alors
$$
f_*[Z]_k = [f_*{\mathcal O}_Z]_k.
$$
\item Soit $\mathcal{F}$ un faisceau cohérent sur $X$ tel que
$\dim_\delta(\text{Supp}(\mathcal{F})) \leq k$. Alors
$$
f_*[\mathcal{F}]_k = [f_*{\mathcal F}]_k.
$$
\end{enumerate}
Remarquons que l'énoncé a un sens puisque $f_*\mathcal{F}$ et $f_*\mathcal{O}_Z$
sont des $\mathcal{O}_Y$-modules cohérents d'après Cohomologie des espaces, lemme
\ref{spaces-cohomology-lemma-proper-pushforward-coherent}.
\end{lemma}

\begin{proof}
La partie (1) découle de (2) et du lemme \ref{lemma-cycle-closed-coherent}.
Soit $\mathcal{F}$ un faisceau cohérent sur $X$. Supposons que
$\dim_\delta(\text{Supp}(\mathcal{F})) \leq k$. D'après Cohomologie des espaces,
lemme \ref{spaces-cohomology-lemma-coherent-support-closed}, il existe une
immersion fermée $i : Z \to X$ et un $\mathcal{O}_Z$-module cohérent
$\mathcal{G}$ tel que $i_*\mathcal{G} \cong \mathcal{F}$ et tel que le support
de $\mathcal{F}$ soit $Z$. Soit $Z' \subset Y$ l'image schématique de
$f|_Z : Z \to Y$ ; voir Morphismes d'espaces, définition
\ref{spaces-morphisms-definition-scheme-theoretic-image}. Considérons le
diagramme commutatif
$$
\xymatrix{
Z \ar[r]_i \ar[d]_{f|_Z} &
X \ar[d]^f \\
Z' \ar[r]^{i'} & Y
}
$$
d'espaces algébriques au-dessus de $B$. Remarquons que $f|_Z$ est surjectif
(cela résulte de Morphismes d'espaces, lemme
\ref{spaces-morphisms-lemma-quasi-compact-scheme-theoretic-image}, et du fait
que $|f|$ est fermé) et propre (cela résulte de Morphismes d'espaces, lemmes
\ref{spaces-morphisms-lemma-base-change-proper},
\ref{spaces-morphisms-lemma-closed-immersion-proper}, et
\ref{spaces-morphisms-lemma-universally-closed-permanence}). En parcourant le
diagramme par les deux chemins, nous obtenons $f_*\mathcal{F} =
f_*i_*\mathcal{G} = i'_*(f|_Z)_*\mathcal{G}$. Supposons le résultat établi pour
les immersions fermées et pour $f|_Z$. Alors
$$
f_*[\mathcal{F}]_k = f_*i_*[\mathcal{G}]_k
= (i')_*(f|_Z)_*[\mathcal{G}]_k =
(i')_*[(f|_Z)_*\mathcal{G}]_k =
[(i')_*(f|_Z)_*\mathcal{G}]_k = [f_*\mathcal{F}]_k
$$
comme voulu. Le cas d'une immersion fermée découle du lemme
\ref{lemma-length-closed-immersion} et des définitions. Nous nous sommes donc
réduits au cas où $\dim_\delta(X) \leq k$ et $f : X \to Y$ est propre et
surjectif.

\medskip\noindent
Supposons que $\dim_\delta(X) \leq k$ et que $f : X \to Y$ soit propre et
surjectif. Pour toute composante irréductible $Z \subset Y$ de point
générique $\eta$, il existe un point $\xi \in X$ tel que $f(\xi) = \eta$. Ainsi
$\delta(\eta) \leq \delta(\xi) \leq k$. On voit donc que dans les expressions
$$
f_*[\mathcal{F}]_k = \sum n_Z[Z],
\quad
\text{et}
\quad
[f_*\mathcal{F}]_k = \sum m_Z[Z].
$$
dès que $n_Z \not = 0$, ou que $m_Z \not = 0$, le sous-espace fermé
intègre $Z$ est en fait une composante irréductible de $Y$ de
$\delta$-dimension $k$ (voir le lemme \ref{lemma-point-of-max-dimension}).
Choisissons un tel sous-espace fermé intègre $Z \subset Y$ et notons $\eta$ son
point générique. Remarquons que, pour tout $\xi \in X$ tel que $f(\xi) = \eta$,
nous avons
$\delta(\xi) \geq k$ et donc $\xi$ est également un point générique d'une
composante irréductible de $X$ de $\delta$-dimension $k$ (voir le lemme
\ref{lemma-point-of-max-dimension}). D'après Espaces sur les corps, lemme
\ref{spaces-over-fields-lemma-finite-in-codim-1}, il existe un sous-espace
ouvert $\eta \in V \subset Y$ tel que $f^{-1}(V) \to V$ soit fini. Puisque
$\eta$ est un point générique d'une composante irréductible de $|Y|$, nous
pouvons supposer que $V$ est un schéma affine ; voir Propriétés des espaces,
proposition
\ref{spaces-properties-proposition-locally-quasi-separated-open-dense-scheme}.
En remplaçant $Y$ par $V$ et $X$ par $f^{-1}(V)$, nous nous réduisons au cas où
$Y$ est affine et où $f$ est fini. En particulier, $X$ et $Y$ sont des schémas,
et nous sommes ramenés au résultat correspondant pour les schémas ; voir
Homologie de Chow, lemme \ref{chow-lemma-cycle-push-sheaf} (appliqué avec $S = Y$).
\end{proof}













\section{Préliminaires à l'image réciproque plate}
\label{section-preparation-flat-pullback}

\noindent
Cette section est l'analogue de l'Homologie de Chow, section
\ref{chow-section-preparation-flat-pullback}.

\medskip\noindent
Rappelons qu'un morphisme d'espaces algébriques est dit de dimension relative
$r$ si, localement pour la topologie \'etale sur la source et sur le but, on
obtient un morphisme de schémas de dimension relative $r$. La définition
précise est équivalente, quoique légèrement différente en réalité ; voir
Morphismes d'espaces, définition
\ref{spaces-morphisms-definition-relative-dimension}.

\begin{lemma}
\label{lemma-flat-inverse-image-dimension}
Dans la situation \ref{situation-setup}, soient $X, Y/B$ bons. Soit $f :
X \to Y$ un morphisme au-dessus de $B$. Supposons que $f$ soit plat de
dimension relative $r$. Pour toute partie fermée $T \subset |Y|$, nous avons
$$
\dim_\delta(|f|^{-1}(T)) = \dim_\delta(T) + r.
$$
pourvu que $|f|^{-1}(T)$ ne soit pas vide. Si $Z \subset Y$ est un
sous-schéma fermé intègre et que $Z' \subset f^{-1}(Z)$ est une composante
irréductible, alors $Z'$ domine $Z$ et $\dim_\delta(Z') = \dim_\delta(Z) + r$.
\end{lemma}

\begin{proof}
Puisque la $\delta$-dimension d'une partie fermée est la borne supérieure des
$\delta$-dimensions de ses composantes irréductibles, il suffit de démontrer
la dernière assertion. Nous pouvons remplacer $Y$ par le sous-schéma fermé intègre
$Z$ et $X$ par $f^{-1}(Z) = Z \times_Y X$. Nous pouvons donc supposer que $Z =
Y$ est intègre et que $f$ est un morphisme plat de dimension relative $r$.
Puisque $Y$ est localement noethérien, le morphisme $f$, qui est localement de
type fini, est en fait localement de présentation finie. Le lemme
\ref{spaces-morphisms-lemma-fppf-open} de Morphismes d'espaces s'applique donc,
et $f$ est ouvert. Soit $\xi \in X$ un point générique d'une composante
irréductible de $X$. Comme $f$ est ouvert, $f(\xi)$ est le point
générique $\eta$ de $Z = Y$. Ainsi $Z'$ domine $Z = Y$. Enfin, la proposition
\ref{spaces-properties-proposition-locally-quasi-separated-open-dense-scheme}
de Propriétés des espaces montre que $\xi$ et $\eta$ appartiennent aux lieux
schématiques de $X$ et de $Y$.
Puisque $\xi$ est un point générique de $X$, nous voyons que $\mathcal{O}_{X,
\xi} = \mathcal{O}_{X_\eta, \xi}$ n'a qu'un seul idéal premier et a donc la
dimension $0$ (on peut utiliser les anneaux locaux usuels, car $\xi$ et
$\eta$ appartiennent aux lieux schématiques de $X$ et de $Y$). Le lemme
\ref{spaces-morphisms-lemma-dimension-fibre-at-a-point} de Morphismes d'espaces
(et la définition des morphismes d'une dimension relative donnée) montre alors que
le degré de transcendance de $\kappa(\xi)$ sur $\kappa(\eta)$ est $r$. En
d'autres termes, $\delta(\xi) = \delta(\eta) + r$, comme voulu.
\end{proof}

\noindent
Voici le lemme que nous utiliserons pour montrer que l'image réciproque plate
d'une famille localement finie de sous-schémas fermés est localement finie.

\begin{lemma}
\label{lemma-inverse-image-locally-finite}
Dans la situation \ref{situation-setup}, soient $X, Y/B$ bons. Soit $f :
X \to Y$ un morphisme au-dessus de $B$. Supposons que $\{T_i\}_{i \in I}$ soit
une famille localement finie de parties fermées de $|Y|$. Alors
$\{|f|^{-1}(T_i)\}_{i \in I}$ est une famille localement finie de
parties fermées de $X$.
\end{lemma}

\begin{proof}
Soit $U \subset |X|$ une partie ouverte quasi-compacte. Puisque l'image
$|f|(U) \subset |Y|$ est une partie quasi-compacte, il existe une partie $V
\subset |Y|$ ouverte quasi-compacte telle que $|f|(U) \subset V$. Remarquons que
$$
\{i \in I : |f|^{-1}(T_i) \cap U \not = \emptyset \}
\subset
\{i \in I : T_i \cap V \not = \emptyset \}.
$$
Le membre de droite étant fini par hypothèse, on conclut.
\end{proof}

















\section{Image réciproque plate}
\label{section-flat-pullback}

\noindent
Cette section est l'analogue de l'Homologie de Chow, section
\ref{chow-section-flat-pullback}.

\medskip\noindent
Soit $S$ un schéma et soit $f : X \to Y$ un morphisme d'espaces algébriques
au-dessus de $S$. Soit $Z \subset Y$ un sous-espace fermé. Dans ce chapitre, nous
utiliserons parfois la terminologie {\it image réciproque schématique}
pour l'image réciproque $f^{-1}(Z)$ de $Z$ construite dans Morphismes des
espaces, définition
\ref{spaces-morphisms-definition-inverse-image-closed-subspace}. L'image
réciproque schématique est le produit fibré
$$
\xymatrix{
f^{-1}(Z) \ar[r] \ar[d] & X \ar[d] \\
Z \ar[r] & Y
}
$$
Si $\mathcal{I} \subset \mathcal{O}_Y$ est le faisceau quasi-cohérent d'idéaux
correspondant à $Z$ dans $Y$, alors $f^{-1}(\mathcal{I})\mathcal{O}_X$ est le
faisceau quasi-cohérent d'idéaux correspondant à $f^{-1}(Z)$ dans $X$.

\begin{definition}
\label{definition-flat-pullback}
Dans la situation \ref{situation-setup}, soient $X, Y/B$ bons. Soit $f :
X \to Y$ un morphisme au-dessus de $B$. Supposons que $f$ soit plat de dimension
relative $r$.
\begin{enumerate}
\item Soit $Z \subset Y$ un sous-espace fermé intègre de $\delta$-dimension
$k$. Nous définissons $f^*[Z]$ comme le $(k+r)$-cycle sur $X$ associé à
l'image réciproque schématique
$$
f^*[Z] = [f^{-1}(Z)]_{k+r}.
$$
Cela a du sens puisque $\dim_\delta(f^{-1}(Z)) = k + r$ par le lemme
\ref{lemma-flat-inverse-image-dimension}.
\item Soit $\alpha = \sum n_i [Z_i]$ un $k$-cycle sur $Y$. L'{\it image
réciproque plate de $\alpha$ par $f$} est la somme
$$
f^* \alpha = \sum n_i f^*[Z_i]
$$
où chaque $f^*[Z_i]$ est défini comme ci-dessus. La somme est localement finie
par le lemme \ref{lemma-inverse-image-locally-finite}.
\item On note $f^* : Z_k(Y) \to Z_{k + r}(X)$ l'homomorphisme de groupes abéliens
ainsi obtenu.
\end{enumerate}
\end{definition}

\noindent
Une immersion ouverte est plate. C'est un cas particulier important, quoique
trivial, de morphisme plat. Si $U \subset X$ est ouvert, l'image réciproque par
$j : U \to X$ d'un cycle est parfois appelée la {\it restriction} du cycle à
$U$. Remarquons que, dans ce cas, les applications
$$
j^* : Z_k(X) \longrightarrow Z_k(U)
$$
sont toutes {\it surjectives}. En effet, étant donné un sous-espace fermé
intègre $Z' \subset U$, nous pouvons prendre pour $Z$ l'adhérence de
$Z'$ dans $X$ et la considérer comme un sous-espace fermé réduit de $X$ (voir
Propriétés des espaces, définition
\ref{spaces-properties-definition-reduced-induced-space}). On a alors clairement
$Z \cap U = Z'$, autrement dit $j^*[Z] = [Z']$, d'où la surjectivité. En fait,
on a un peu mieux.

\begin{lemma}
\label{lemma-exact-sequence-open}
Dans la situation \ref{situation-setup}, soit $X/B$ bon. Soit $U
\subset X$ un sous-espace ouvert. Soit $Y$ le sous-espace fermé réduit de $X$
avec $|Y| = |X| \setminus |U|$ et notons $i : Y \to X$ le morphisme
d'inclusion. Pour chaque $k \in \mathbf{Z}$, la séquence
$$
\xymatrix{
Z_k(Y) \ar[r]^{i_*} & Z_k(X) \ar[r]^{j^*} & Z_k(U) \ar[r] & 0
}
$$
est une suite exacte de groupes abéliens.
\end{lemma}

\begin{proof}
Nous avons vu ci-dessus que $j^*$ est surjectif. Supposons d'abord que $X$
soit quasi-compact. Alors $Z_k(X)$ est un $\mathbf{Z}$-module libre de base
donnée par les éléments $[Z]$ où $Z \subset X$ est fermé intègre de
$\delta$-dimension $k$. Un tel élément de base s'envoie soit sur l'élément de
base $[Z \cap U]$ de $Z_k(U)$, soit sur zéro si $Z \subset Y$. Le lemme est donc
clair dans ce cas. Le cas général est similaire et la preuve est omise.
\end{proof}

\begin{lemma}
\label{lemma-etale-pullback}
Dans la situation \ref{situation-setup}, soit $f : X \to Y$ un morphisme
\'etale de bons espaces algébriques au-dessus de $B$. Si $Z \subset Y$ est un
sous-espace fermé intègre, alors $f^*[Z] = \sum [Z']$ où la somme est sur les
composantes irréductibles (remarque \ref{remark-irreducible-component}) de
$f^{-1}(Z)$.
\end{lemma}

\begin{proof}
La signification du lemme est que le coefficient de $[Z']$ est $1$. Cela
découle du fait que $f^{-1}(Z)$ est un espace algébrique réduit car il est
\'etale sur l'espace algébrique intègre $Z$.
\end{proof}

\begin{lemma}
\label{lemma-compose-flat-pullback}
Dans la situation \ref{situation-setup}, soient $X, Y, Z/B$ bons. Soient $f
: X \to Y$ et $g : Y \to Z$ des morphismes plats de dimensions relatives $r$
et $s$ au-dessus de $B$. Alors $g \circ f$ est plat de dimension relative $r + s$ et
$$
f^* \circ g^* = (g \circ f)^*
$$
comme applications $Z_k(Z) \to Z_{k + r + s}(X)$.
\end{lemma}

\begin{proof}
La composition est plate de dimension relative $r + s$ d'après Morphismes
d'espaces, lemmes
\ref{spaces-morphisms-lemma-dimension-fibre-at-a-point-additive} et
\ref{spaces-morphisms-lemma-composition-flat}. Supposons que
\begin{enumerate}
\item $A \subset Z$ est un sous-espace fermé intègre de $\delta$-dimension
$k$,
\item $A' \subset Y$ est un sous-espace fermé intègre de $\delta$-dimension
$k + s$ tel que $A' \subset g^{-1}(A)$, et
\item $A'' \subset Y$ est un sous-espace fermé intègre de $\delta$-dimension
$k + s + r$ tel que $A'' \subset f^{-1}(W')$.
\end{enumerate}
Nous devons montrer que le coefficient $n$ de $[A'']$ dans $(g \circ f)^*[A]$
est égal au coefficient $m$ de $[A'']$ dans $f^*(g^*[A])$. On peut
choisir un diagramme commutatif
$$
\xymatrix{
U \ar[d] \ar[r] & V \ar[d] \ar[r] & W \ar[d] \\
X \ar[r] & Y \ar[r] & Z
}
$$
où $U, V, W$ sont des schémas, les flèches verticales sont \'etales, et il
existe des points $u \in U$, $v \in V$, $w \in W$ tels que $u \mapsto v
\mapsto w$ et tels que $u, v, w$ s'envoient sur les points génériques de $A'',
A', A$. (Détails omis.) Nous avons alors des homomorphismes plats d'anneaux
locaux $\mathcal{O}_{W, w} \to \mathcal{O}_{V, v}$, $\mathcal{O}_{V, v} \to
\mathcal{O}_{U, u}$, et en utilisant à plusieurs reprises le lemme
\ref{lemma-length}, nous trouvons
$$
n = \text{length}_{\mathcal{O}_{U, u}}(
\mathcal{O}_{U, u}/\mathfrak m_w\mathcal{O}_{U, u})
$$
et
$$
m =
\text{length}_{\mathcal{O}_{V, v}}(
\mathcal{O}_{V, v}/\mathfrak m_w\mathcal{O}_{V, v})
\text{length}_{\mathcal{O}_{U, u}}(
\mathcal{O}_{U, u}/\mathfrak m_v\mathcal{O}_{U, u})
$$
L'égalité découle donc du lemme d'Algèbre
\ref{algebra-lemma-pullback-transitive}.
\end{proof}

\begin{lemma}
\label{lemma-pullback-coherent}
Dans la situation \ref{situation-setup}, soient $X, Y/B$ bons. Soit $f :
X \to Y$ un morphisme plat de dimension relative $r$.
\begin{enumerate}
\item Soit $Z \subset Y$ un sous-espace fermé avec $\dim_\delta(Z) \leq k$.
Alors $\dim_\delta(f^{-1}(Z)) \leq k + r$ et $[f^{-1}(Z)]_{k +
r} = f^*[Z]_k$ dans $Z_{k + r}(X)$.
\item Soit $\mathcal{F}$ un faisceau cohérent sur $Y$ avec
$\dim_\delta(\text{Supp}(\mathcal{F})) \leq k$. Nous avons alors
$\dim_\delta(\text{Supp}(f^*\mathcal{F})) \leq k + r$ et
$$
f^*[{\mathcal F}]_k = [f^*{\mathcal F}]_{k+r}
$$
dans $Z_{k + r}(X)$.
\end{enumerate}
\end{lemma}

\begin{proof}
La partie (1) découle de la partie (2) par le lemme
\ref{lemma-cycle-closed-coherent} et le fait que $f^*\mathcal{O}_Z =
\mathcal{O}_{f^{-1}(Z)}$.

\medskip\noindent
Preuve de (2). Comme $X$, $Y$ sont localement noethériens, nous pouvons
appliquer le lemme
\ref{spaces-cohomology-lemma-coherent-Noetherian} de Cohomologie des espaces :
$\mathcal{F}$ est de type fini, donc $f^*\mathcal{F}$ est de type fini (Modules
sur les sites, lemme \ref{sites-modules-lemma-local-pullback}), et par suite
$f^*\mathcal{F}$ est cohérent (de nouveau Cohomologie des espaces, lemme
\ref{spaces-cohomology-lemma-coherent-Noetherian}). Le lemme a donc un sens.
Soit $W \subset Y$ un sous-espace fermé intègre de $\delta$-dimension $k$, et
soit $W' \subset X$ un sous-espace fermé intègre de dimension $k + r$
s'envoyant dans $W$ par $f$. Nous devons montrer que le
coefficient $n$ de $[W']$ dans $f^*[{\mathcal F}]_k$ est égal au
coefficient $m$ de $[W']$ dans $[f^*{\mathcal F}]_{k+r}$. On peut choisir un
diagramme commutatif
$$
\xymatrix{
U \ar[d] \ar[r] & V \ar[d] \\
X \ar[r] & Y
}
$$
où $U, V$ sont des schémas, les flèches verticales sont \'etales, et il existe
des points $u \in U$, $v \in V$ tels que $u \mapsto v$ et tels que $u, v$
s'envoient sur les points génériques de $W', W$. (Détails omis.) Considérons le
germe $M = (\mathcal{F}|_V)_v$ comme un $\mathcal{O}_{V, v}$-module.
(Remarquons que $M$ est de longueur finie d'après nos hypothèses de dimension,
mais nous n'avons en fait pas besoin de le vérifier ; voir le lemme
\ref{lemma-length-finite}.) Nous avons $(f^*\mathcal{F}|_U)_u =
\mathcal{O}_{U, u} \otimes_{\mathcal{O}_{V, v}} M$. Ainsi on voit que
$$
n = \text{length}_{\mathcal{O}_{U, u}}
(\mathcal{O}_{U, u} \otimes_{\mathcal{O}_{V, v}} M)
\quad
\text{et}
\quad
m = \text{length}_{\mathcal{O}_{V, v}}(M)
\text{length}_{\mathcal{O}_{V, v}}(
\mathcal{O}_{U, u}/\mathfrak m_v \mathcal{O}_{U, u})
$$
Ainsi, l'égalité découle du lemme d'Algèbre
\ref{algebra-lemma-pullback-module}.
\end{proof}









\section{Image directe et image réciproque}
\label{section-push-pull}

\noindent
Cette section est l'analogue de l'Homologie de Chow, section
\ref{chow-section-push-pull}.

\medskip\noindent
Dans cette section, nous vérifions que l'image directe propre et l'image
réciproque plate sont compatibles lorsque cela a un sens. D'après le travail que
nous avons effectué ci-dessus, cela est une conséquence de la cohomologie et
du changement de base.

\begin{lemma}
\label{lemma-flat-pullback-proper-pushforward}
Dans la situation \ref{situation-setup}, soit
$$
\xymatrix{
X' \ar[r]_{g'} \ar[d]_{f'} & X \ar[d]^f \\
Y' \ar[r]^g & Y
}
$$
un diagramme cartésien de bons espaces algébriques au-dessus de $B$.
Supposons $f : X \to Y$ propre et $g : Y' \to Y$ plat de dimension
relative $r$. Alors également $f'$ est propre et $g'$ est plat de dimension
relative $r$. Pour tout $k$-cycle $\alpha$ sur $X$, nous avons
$$
g^*f_*\alpha = f'_*(g')^*\alpha
$$
dans $Z_{k + r}(Y')$.
\end{lemma}

\begin{proof}
L'assertion selon laquelle $f'$ est propre découle de Morphismes d'espaces,
lemme \ref{spaces-morphisms-lemma-base-change-proper}. Celle selon laquelle
$g'$ est plat de dimension relative $r$ découle de Morphismes d'espaces, lemmes
\ref{spaces-morphisms-lemma-dimension-fibre-after-base-change} et
\ref{spaces-morphisms-lemma-base-change-flat}. Il suffit de prouver l'égalité
des cycles lorsque $\alpha = [W]$ pour un sous-espace fermé intègre $W
\subset X$ de $\delta$-dimension $k$. Remarquons que, dans ce cas, $\alpha
= [\mathcal{O}_W]_k$ ; voir le lemme \ref{lemma-cycle-closed-coherent}. D'après les
lemmes \ref{lemma-cycle-push-sheaf} et \ref{lemma-pullback-coherent} il suffit
donc de montrer que $f'_*(g')^*\mathcal{O}_W$ est isomorphe à
$g^*f_*\mathcal{O}_W$. Cela découle de la cohomologie et du changement de
base ; voir Cohomologie des espaces, lemme
\ref{spaces-cohomology-lemma-flat-base-change-cohomology}.
\end{proof}

\begin{lemma}
\label{lemma-finite-flat}
Dans la situation \ref{situation-setup}, soient $X, Y/B$ bons. Soit $f :
X \to Y$ un morphisme fini localement libre de degré $d$ (voir Morphismes des
espaces, définition \ref{spaces-morphisms-definition-finite-locally-free}).
Alors $f$ est à la fois propre et plat de dimension relative $0$, et
$$
f_*f^*\alpha = d\alpha
$$
pour tout $\alpha \in Z_k(Y)$.
\end{lemma}

\begin{proof}
Un morphisme fini localement libre est plat et fini d'après Morphismes
d'espaces, lemme \ref{spaces-morphisms-lemma-finite-flat}, et un morphisme
fini est propre d'après Morphismes d'espaces, lemme
\ref{spaces-morphisms-lemma-finite-proper}. Nous omettons de vérifier qu'un
morphisme fini est de dimension relative $0$. La formule a donc un sens. Pour
la démontrer, soit $Z \subset Y$ un sous-schéma fermé intègre de
$\delta$-dimension $k$. Il suffit de démontrer la formule pour
$\alpha = [Z]$. Puisque le changement de base d'un morphisme fini localement
libre est fini localement libre (Morphismes des espaces, lemme
\ref{spaces-morphisms-lemma-base-change-finite-locally-free}), on voit que
$f_*f^*\mathcal{O}_Z$ est un faisceau localement libre de rang fini $d$ sur
$Z$. Il est donc clair que $f_*f^*\mathcal{O}_Z$ est de longueur $d$ au point
générique de $Z$. Par conséquent,
$$
f_*f^*[Z] = f_*f^*[\mathcal{O}_Z]_k =
[f_*f^*\mathcal{O}_Z]_k = d[Z]
$$
où nous avons utilisé les lemmes \ref{lemma-pullback-coherent} et
\ref{lemma-cycle-push-sheaf}.
\end{proof}










\section{Préliminaires aux diviseurs principaux}
\label{section-preparation-principal-divisors}

\noindent
Cette section est l'analogue de l'Homologie de Chow, section
\ref{chow-section-preparation-principal-divisors}. Certains éléments de cette
section recoupent en partie la discussion d'Espaces sur les corps, section
\ref{spaces-over-fields-section-Weil-divisors}.

\begin{lemma}
\label{lemma-divisor-delta-dimension}
Dans la situation \ref{situation-setup}, soit $X/B$ bon. Supposons que
$X$ soit intègre.
\begin{enumerate}
\item Si $Z \subset X$ est un sous-espace fermé intègre, alors les assertions
suivantes sont équivalentes :
\begin{enumerate}
\item $Z$ est un diviseur premier,
\item $|Z|$ est de codimension $1$ dans $|X|$, et
\item $\dim_\delta(Z) = \dim_\delta(X) - 1$.
\end{enumerate}
\item Si $Z$ est une composante irréductible d'un diviseur de Cartier effectif
sur $X$, alors $\dim_\delta(Z) = \dim_\delta(X) - 1$.
\end{enumerate}
\end{lemma}

\begin{proof}
La partie (1) découle de la définition d'un diviseur premier (Espaces sur les
corps, définition \ref{spaces-over-fields-definition-Weil-divisor}), d'Espaces
décents, lemme \ref{decent-spaces-lemma-codimension-local-ring}, et de la
définition d'une fonction de dimension (Topologie, définition
\ref{topology-definition-dimension-function}).

\medskip\noindent
Soit $D \subset X$ un diviseur de Cartier effectif. Soit $Z \subset D$
une composante irréductible et soit $\xi \in |Z|$ le point générique.
Choisissons un voisinage \'etale $(U, u) \to (X, \xi)$ où $U = \Spec(A)$ et où
$D \times_X U$ est défini par un non-diviseur de zéro $f \in A$ ; voir
Diviseurs sur les espaces, lemme
\ref{spaces-divisors-lemma-characterize-effective-Cartier-divisor}. Alors $u$
est le point générique de $V(f)$ d'après Espaces décents, lemme
\ref{decent-spaces-lemma-decent-generic-points}. Par conséquent,
$\mathcal{O}_{U, u}$ est de dimension $1$ d'après le théorème de l'idéal
principal de Krull (Algèbre, lemme \ref{algebra-lemma-minimal-over-1}). Ainsi,
$\xi$ est un point de codimension $1$ sur $X$ et $Z$ est un diviseur premier,
comme voulu.
\end{proof}












\section{Diviseurs principaux}
\label{section-principal-divisors}

\noindent
Cette section est l'analogue de l'Homologie de Chow, section
\ref{chow-section-principal-divisors}. La définition suivante est l'analogue
d'Espaces sur les corps, définition
\ref{spaces-over-fields-definition-principal-divisor}, dans notre cadre.

\begin{definition}
\label{definition-principal-divisor}
Dans la situation \ref{situation-setup}, soit $X/B$ bon. Supposons que
$X$ soit intègre et que $\dim_\delta(X) = n$. Soit $f \in R(X)^*$. Le
{\it diviseur principal associé à $f$} est le $(n - 1)$-cycle
$$
\text{div}(f) = \text{div}_X(f) = \sum \text{ord}_Z(f) [Z]
$$
défini dans Espaces sur les corps, définition
\ref{spaces-over-fields-definition-principal-divisor}. Cela a un sens car
les diviseurs premiers sont de $\delta$-dimension $n - 1$ par le lemme
\ref{lemma-divisor-delta-dimension}.
\end{definition}

\noindent
Dans la situation de la définition, pour $f, g \in R(X)^*$, nous avons
$$
\text{div}_X(fg) = \text{div}_X(f) + \text{div}_X(g)
$$
dans $Z_{n - 1}(X)$. Voir Espaces sur les corps, lemme
\ref{spaces-over-fields-lemma-div-additive}. Le lemme suivant nous permettra
de ramener les assertions sur les diviseurs principaux au cas des schémas.

\begin{lemma}
\label{lemma-etale-pullback-principal-divisor}
Dans la situation \ref{situation-setup}, soit $f : X \to Y$ un morphisme
\'etale de bons espaces algébriques au-dessus de $B$. Supposons que $Y$ soit intègre.
Soit $g \in R(Y)^*$. Comme cycles sur $X$ nous avons
$$
f^*(\text{div}_Y(g)) =
\sum\nolimits_{X'} (X' \to X)_*\text{div}_{X'}(g \circ f|_{X'})
$$
où la somme s'étend sur les composantes irréductibles de $X$ (remarque
\ref{remark-irreducible-component}).
\end{lemma}

\begin{proof}
L'application $|X| \to |Y|$ est ouverte. L'ensemble des composantes irréductibles
de $|X|$ est localement fini dans $|X|$. Nous concluons que $f|_{X'} : X' \to
Y$ est dominant pour toute composante irréductible $X' \subset X$. Ainsi $g
\circ f|_{X'}$ est défini (Morphismes des espaces, section
\ref{spaces-morphisms-section-rational-maps}), donc $\text{div}_{X'}(g \circ
f|_{X'})$ est défini. De plus, la somme est localement finie et on trouve que
le membre de droite est bien un cycle sur $X$. Le membre de gauche est défini par
la définition \ref{definition-flat-pullback} et le fait qu'un morphisme
\'etale est plat de dimension relative $0$.

\medskip\noindent
Puisque $f$ est \'etale, nous voyons que $\delta_X(x) = \delta_y(f(x))$ pour
tout $x \in |X|$. Ainsi, si $\dim_\delta(Y) = n$, alors $\dim_\delta(X') =
n$ pour toute composante irréductible $X'$ de $X$ (puisque les points
génériques de $X$ s'envoient sur le point générique de $Y$ ; voir ci-dessus).
Les deux membres sont donc des $(n - 1)$-cycles.

\medskip\noindent
Soit $Z \subset X$ un sous-espace fermé intègre avec $\dim_\delta(Z) = n -
1$. Pour démontrer l'égalité, il faut montrer que les coefficients de $Z$
dans les deux membres sont égaux. Soit $Z' \subset Y$ le sous-espace fermé intègre construit
dans le lemme \ref{lemma-proper-image}. Alors $\dim_\delta(Z') = n - 1$ aussi.
Soit $\xi \in |Z|$ le point générique. Alors $\xi' = f(\xi) \in |Z'|$ est le
point générique. Considérons le diagramme commutatif
$$
\xymatrix{
\Spec(\mathcal{O}_{X, \xi}^h) \ar[r] \ar[d] & X \ar[d] \\
\Spec(\mathcal{O}_{Y, \xi'}^h) \ar[r] & Y
}
$$
d'Espaces décents, remarque
\ref{decent-spaces-remark-functoriality-henselian-local-ring}. Nous devons
être légèrement prudents car les anneaux locaux noéthériens réduits
$\mathcal{O}_{X, \xi}^h$ et $\mathcal{O}_{Y, \xi'}^h$ ne sont pas
nécessairement intègres. Nous travaillons donc avec des anneaux totaux de
fractions $Q(-)$ plutôt que des corps des fractions. Par définition, pour
obtenir le coefficient de $Z'$ dans $\text{div}_Y(g)$ nous écrivons l'image de
$g$ dans $Q(\mathcal{O}_{Y, \xi'}^h)$ comme $a/b$ avec $a, b \in
\mathcal{O}_{Y, \xi'}^h$ non-diviseurs de zéro et nous prenons
$$
\text{ord}_{Z'}(g) =
\text{length}_{\mathcal{O}_{Y, \xi'}^h}
(\mathcal{O}_{Y, \xi'}^h/a \mathcal{O}_{Y, \xi'}^h) -
\text{length}_{\mathcal{O}_{Y, \xi'}^h}
(\mathcal{O}_{Y, \xi'}^h/b \mathcal{O}_{Y, \xi'}^h)
$$
Remarquons que le coefficient de $Z$ dans $f^*\text{div}_Y(G)$ est ce même
entier ; voir le lemme \ref{lemma-etale-pullback}. Supposons que $\xi \in X'$.
On peut alors considérer les applications
$$
\mathcal{O}_{Y, \xi'}^h \to
\mathcal{O}_{X, \xi}^h \to
\mathcal{O}_{X', \xi}^h
$$
La première flèche est plate et la seconde est un homomorphisme surjectif
d'anneaux locaux noethériens réduits de dimension $1$. Par conséquent, ces
deux homomorphismes envoient les non-diviseurs de zéro sur des non-diviseurs
de zéro, et le coefficient de $Z'$ dans $\text{div}_{X'}(g \circ f|_{X'})$
est
$$
\text{ord}_Z(g \circ f|_{X'}) =
\text{length}_{\mathcal{O}_{X', \xi}^h}
(\mathcal{O}_{X', \xi}^h/a \mathcal{O}_{X', \xi}^h) -
\text{length}_{\mathcal{O}_{Y, \xi'}^h}
(\mathcal{O}_{X', \xi}^h/b \mathcal{O}_{X', \xi}^h)
$$
par la même prescription que ci-dessus. Il suffit donc de montrer
$$
\text{length}_{\mathcal{O}_{Y, \xi'}^h}
(\mathcal{O}_{Y, \xi'}^h/a \mathcal{O}_{Y, \xi'}^h) =
\sum\nolimits_{\xi \in |X'|}
\text{length}_{\mathcal{O}_{X', \xi}^h}
(\mathcal{O}_{X', \xi}^h/a \mathcal{O}_{X', \xi}^h)
$$
Tout d'abord, puisque l'homomorphisme d'anneaux $\mathcal{O}_{Y, \xi'}^h \to
\mathcal{O}_{X, \xi}^h$ est plate et non ramifiée, nous avons
$$
\text{length}_{\mathcal{O}_{Y, \xi'}^h}
(\mathcal{O}_{Y, \xi'}^h/a \mathcal{O}_{Y, \xi'}^h) =
\text{length}_{\mathcal{O}_{X, \xi}^h}
(\mathcal{O}_{X, \xi}^h/a \mathcal{O}_{X, \xi}^h)
$$
par Algèbre, lemme \ref{algebra-lemma-pullback-module}. Soient $\mathfrak q_1,
\ldots, \mathfrak q_t$ les idéaux premiers non maximaux de $\mathcal{O}_{X,
\xi}^h$, et posons $R_j = \mathcal{O}_{X, \xi}^h/\mathfrak q_j$. Pour $X'$
comme ci-dessus, notons $J(X') \subset \{1, \ldots, t\}$ l'ensemble des
indices tels que $\mathfrak q_j$ corresponde à un point de $X'$, c'est-à-dire
tels que, par l'homomorphisme surjectif $\mathcal{O}_{X, \xi}^h \to
\mathcal{O}_{X', \xi}$, l'idéal premier $\mathfrak q_j$ corresponde à un idéal
premier de $\mathcal{O}_{X', \xi}$. D'après Homologie de Chow, lemme
\ref{chow-lemma-additivity-divisors-restricted}, on obtient
$$
\text{length}_{\mathcal{O}_{X, \xi}^h}
(\mathcal{O}_{X, \xi}^h/a \mathcal{O}_{X, \xi}^h) =
\sum\nolimits_j \text{length}_{R_j}(R_j/a R_j)
$$
et
$$
\text{length}_{\mathcal{O}_{X', \xi}^h}
(\mathcal{O}_{X', \xi}^h/a \mathcal{O}_{X', \xi}^h) =
\sum\nolimits_{j \in J(X')} \text{length}_{R_j}(R_j/a R_j)
$$
Le résultat du lemme en découle, car $\{1, \ldots, t\}$ est l'union
disjointe des ensembles $J(X')$ : chaque point de codimension $0$ sur $X$
appartient à un unique $X'$.
\end{proof}







\section{Diviseurs principaux et image directe}
\label{section-two-fun}

\noindent
Cette section est l'analogue de l'Homologie de Chow, section
\ref{chow-section-two-fun}.

\begin{lemma}
\label{lemma-proper-pushforward-alteration}
Dans la situation \ref{situation-setup}, soient $X, Y/B$ bons. Supposons
que $X$, $Y$ soient intègres et que $n = \dim_\delta(X) = \dim_\delta(Y)$.
Soit $p : X \to Y$ un morphisme propre dominant. Soit $f \in R(X)^*$.
Posons
$$
g = \text{Nm}_{R(X)/R(Y)}(f).
$$
Nous avons alors $p_*\text{div}(f) = \text{div}(g)$.
\end{lemma}

\begin{proof}
Nous allons le déduire du cas des schémas par localisation \'etale. Soit $Z
\subset Y$ un sous-espace fermé intègre de $\delta$-dimension $n - 1$.
Nous voulons montrer que les coefficients de $[Z]$ dans $p_*\text{div}(f)$ et
$\text{div}(g)$ sont égaux. Appliquons le lemme
\ref{spaces-over-fields-lemma-finite-in-codim-1} d'Espaces sur les corps au
morphisme $p : X \to Y$ et au point générique $\xi \in |Z|$. Nous pouvons
alors remplacer $Y$ par un sous-espace ouvert contenant $\xi$ et supposer que
$p : X \to Y$ est fini. Choisissons un voisinage \'etale $(V, v) \to (Y, \xi)$
où $V$ est un schéma affine. D'après le lemme \ref{lemma-etale-pullback}, il
suffit de démontrer l'égalité des cycles après image réciproque sur $V$. Posons
$U = V \times_Y X$ et
considérez le diagramme commutatif
$$
\xymatrix{
U \ar[r]_a \ar[d]_{p'} & X \ar[d]^p \\
V \ar[r]^b & Y
}
$$
Soient $V_j \subset V$, $j = 1, \ldots, m$, les composantes irréductibles de
$V$. Pour chaque $i$, soient $U_{j, i}$, $i = 1, \ldots, n_j$, les composantes
irréductibles de $U$ dominant $V_j$. Notons $p'_{j, i} : U_{j, i} \to V_j$ la
restriction de $p' : U \to V$. D'après le cas des schémas (Homologie de Chow,
lemme \ref{chow-lemma-proper-pushforward-alteration}), on a
$$
p'_{j, i, *}\text{div}_{U_{j, i}}(f_{j, i}) = \text{div}_{V_j}(g_{j, i})
$$
où $f_{j, i}$ est la restriction de $f$ à $U_{j, i}$ et $g_{j, i}$ est la
norme de $f_{j, i}$ pour l'extension finie $R(U_{j, i})/R(V_j)$. Nous
avons
\begin{align*}
b^* p_*\text{div}_X(f)
& =
p'_* a^* \text{div}_X(f) \\
& =
p'_*\left(\sum\nolimits_{j, i}
(U_{j, i} \to U)_*\text{div}_{U_{j, i}}(f_{j, i})\right) \\
& =
\sum\nolimits_{j, i}
(V_j \to V)_*p'_{j, i, *}\text{div}_{U_{j, i}}(f_{j, i}) \\
& =
\sum\nolimits_j
(V_j \to V)_*\left(\sum\nolimits_i \text{div}_{V_j}(g_{j, i})\right) \\
& =
\sum\nolimits_j (V_j \to V)_*\text{div}_{V_j}(\prod\nolimits_i g_{j, i})
\end{align*}
par les lemmes \ref{lemma-flat-pullback-proper-pushforward},
\ref{lemma-etale-pullback-principal-divisor} et
\ref{lemma-compose-pushforward}. Pour terminer la preuve, en utilisant à
nouveau le lemme \ref{lemma-etale-pullback-principal-divisor}, il suffit de
montrer que
$$
g \circ b|_{V_j} = \prod\nolimits_i g_{j, i}
$$
comme éléments du corps des fonctions de $V_j$. En termes de corps, il s'agit
de l'énoncé suivant : soit $L/K$ une extension finie. Soit $M/K$ une extension
finie séparable. Écrivons $M \otimes_K L = \prod M_i$. Alors, pour $t \in L$
dont les images sont $t_i \in M_i$, l'image de $\text{Norm}_{L/K}(t)$ dans $M$ est
$\prod \text{Norm}_{M_i/M}(t_i)$. Nous omettons la preuve.
\end{proof}







\section{Équivalence rationnelle}
\label{section-rational-equivalence}

\noindent
Cette section est l'analogue de l'Homologie de Chow, section
\ref{chow-section-rational-equivalence}. Dans cette section, nous définissons
l'{\it équivalence rationnelle} sur les $k$-cycles. Nous autoriserons des
sommes localement finies d'images de diviseurs principaux (par immersions
fermées). Cela conduit à des phénomènes assez étranges (voir les exemples du
chapitre sur les schémas). Toutefois, sans les autoriser, nous ne savons pas
montrer que l'intersection avec les classes de Chern de
fibrés en droites passe au quotient par l'équivalence rationnelle.

\begin{definition}
\label{definition-rational-equivalence}
Dans la situation \ref{situation-setup}, soit $X/B$ bon. Soit $k
\in \mathbf{Z}$.
\begin{enumerate}
\item Étant donné une famille localement finie $\{W_j \subset X\}$ de
sous-espaces fermés intègres telle que $\dim_\delta(W_j) = k + 1$, et un
élément $f_j \in R(W_j)^*$ quelconque, nous pouvons considérer
$$
\sum (i_j)_*\text{div}(f_j) \in Z_k(X)
$$
où $i_j : W_j \to X$ est le morphisme d'inclusion. Cela a du sens puisque le
morphisme $\coprod i_j : \coprod W_j \to X$ est propre.
\item On dit que $\alpha \in Z_k(X)$ est {\it rationnellement équivalent à
zéro} si $\alpha$ est un cycle de la forme affichée ci-dessus.
\item On dit que $\alpha, \beta \in Z_k(X)$ sont {\it rationnellement
équivalents} et on écrit $\alpha \sim_{rat} \beta$ si $\alpha - \beta$ est
rationnellement équivalent à zéro.
\item Nous définissons
$$
\CH_k(X) = Z_k(X) / \sim_{rat}
$$
comme le {\it groupe de Chow des $k$-cycles sur $X$}. On l'appelle parfois le
{\it groupe de Chow des $k$-cycles modulo l'équivalence rationnelle sur $X$}.
\end{enumerate}
\end{definition}

\noindent
Il existe de nombreuses autres relations d'équivalence intéressantes.
L'équivalence rationnelle est la plus grossière de toutes. Un lemme très
simple mais important est le suivant.

\begin{lemma}
\label{lemma-restrict-to-open}
Dans la situation \ref{situation-setup}, soit $X/B$ bon. Soit $U
\subset X$ un sous-espace ouvert. Soit $Y$ le sous-espace fermé réduit de $X$
avec $|Y| = |X| \setminus |U|$ et notons $i : Y \to X$ le morphisme
d'inclusion. Soit $k \in \mathbf{Z}$. Supposons que $\alpha, \beta \in
Z_k(X)$. Si $\alpha|_U \sim_{rat} \beta|_U$ alors il existe un cycle $\gamma
\in Z_k(Y)$ tel que
$$
\alpha \sim_{rat} \beta + i_*\gamma.
$$
Autrement dit, la séquence
$$
\xymatrix{
\CH_k(Y) \ar[r]^{i_*} & \CH_k(X) \ar[r]^{j^*} & \CH_k(U) \ar[r] & 0
}
$$
est une suite exacte de groupes abéliens.
\end{lemma}

\begin{proof}
Soit $\{W_j\}_{j \in J}$ une famille localement finie de sous-espaces
fermés intègres de $U$ de $\delta$-dimension $k + 1$, et soient $f_j \in
R(W_j)^*$ des éléments tels que $(\alpha - \beta)|_U = \sum
(i_j)_*\text{div}(f_j)$ comme dans la définition. Soit $W_j' \subset X$ le
sous-espace fermé intègre correspondant de $X$, c'est-à-dire ayant le même
point générique que $W_j$. Supposons que $V \subset X$ soit un ouvert
quasi-compact. Alors $V \cap U$ est également un ouvert quasi-compact de $U$,
car $V$ est noethérien. Ainsi, l'ensemble $\{j \in J \mid W_j \cap V \not =
\emptyset\} = \{j \in J \mid W'_j \cap V \not = \emptyset\}$ est fini puisque
$\{W_j\}$ est localement finie. Autrement dit, $\{W'_j\}$ est également
localement finie. Puisque $R(W_j) = R(W'_j)$, on voit que
$$
\alpha - \beta - \sum (i'_j)_*\text{div}(f_j)
$$
est un cycle sur $X$ dont la restriction à $U$ est nulle. Le lemme découle
alors du lemme \ref{lemma-exact-sequence-open}.
\end{proof}

\begin{remark}
\label{remark-infinite-sums-rational-equivalences}
Dans la situation \ref{situation-setup}, soit $X/B$ bon. Supposons que
soient données des familles infinies $\alpha_i, \beta_i \in Z_k(X)$,
$i \in I$, de $k$-cycles sur $X$. Supposons que les supports de $\alpha_i$ et
$\beta_i$ forment des familles localement finies de parties fermées de $X$, de
sorte que $\sum \alpha_i$ et $\sum \beta_i$ soient définies comme cycles. De
plus, supposons que $\alpha_i \sim_{rat} \beta_i$ pour tout $i$. Il n'est alors
pas clair que $\sum \alpha_i \sim_{rat} \sum \beta_i$. En effet, le problème
est que les équivalences rationnelles peuvent être données
par des familles localement finies $\{W_{i, j}, f_{i, j} \in R(W_{i,
j})^*\}_{j \in J_i}$ mais que l'union $\{W_{i, j}\}_{i \in I, j\in J_i}$ peut
ne pas être localement finie.

\medskip\noindent
En pratique, dans de nombreux cas, on dispose d'une famille localement finie de
parties fermées $\{T_i\}_{i \in I}$ de $|X|$ telle que $\alpha_i,
\beta_i$ soient supportés sur $T_i$ et que $\alpha_i \sim_{rat} \beta_i$ «
sur » $T_i$. Plus précisément, les familles $\{W_{i, j}, f_{i, j} \in R(W_{i,
j})^*\}_{j \in J_i}$ sont constituées de sous-espaces fermés intègres $W_{i,
j}$ tels que $|W_{i, j}| \subset T_i$. Dans ce cas, on a bien $\sum \alpha_i
\sim_{rat} \sum \beta_i$ sur $X$, simplement parce que la famille
$\{W_{i, j}\}_{i \in I, j\in J_i}$ est alors automatiquement localement finie.
\end{remark}











\section{Équivalence rationnelle, image directe et image réciproque}
\label{section-properties-rational-equivalence}

\noindent
Cette section est l'analogue de l'Homologie de Chow, section
\ref{chow-section-properties-rational-equivalence}. Dans cette section, nous
montrons que l'image réciproque plate et l'image directe propre sont
compatibles avec l'équivalence rationnelle.

\begin{lemma}
\label{lemma-prepare-flat-pullback-rational-equivalence}
Dans la situation \ref{situation-setup}, soient $X, Y/B$ bons. Supposons
que $Y$ soit intègre et que $\dim_\delta(Y) = k$. Soit $f : X \to Y$ un morphisme
plat de dimension relative $r$. Alors pour $g \in R(Y)^*$ nous avons
$$
f^*\text{div}_Y(g) =
\sum m_{X', X} (X' \to X)_*\text{div}_{X'}(g \circ f|_{X'})
$$
comme $(k + r - 1)$-cycles sur $X$, où la somme est indexée par les composantes
irréductibles $X'$ de $X$, et où $m_{X', X}$ est la multiplicité de $X'$ dans $X$.
\end{lemma}

\begin{proof}
Remarquons que toute composante irréductible de $X$ domine $Y$ (lemme
\ref{lemma-flat-inverse-image-dimension}) et donc la composition $g \circ
f|_{X'}$ est définie (Morphismes des espaces, section
\ref{spaces-morphisms-section-rational-maps}). Nous réduirons cela au cas des
schémas. Choisissons un schéma $V$ et un morphisme \'etale surjectif $V \to Y$.
Choisissons un schéma $U$ et un morphisme \'etale surjectif $U \to V \times_Y
X$. Considérons le diagramme $$
\xymatrix{
U \ar[r]_a \ar[d]_h & X \ar[d]^f \\
V \ar[r]^b & Y
}
$$
Puisque $a$ est surjectif et \'etale, il découle du lemme
\ref{lemma-etale-pullback} qu'il suffit de prouver l'égalité des cycles après
image réciproque par $a$. On peut utiliser le lemme
\ref{lemma-etale-pullback-principal-divisor} pour écrire
$$
b^*\text{div}_Y(g) = \sum (V' \to V)_*\text{div}_{V'}(g \circ b|_{V'})
$$
où la somme est indexée par les composantes irréductibles $V'$ de $V$. En utilisant le
lemme \ref{lemma-flat-pullback-proper-pushforward}, on trouve
$$
h^*b^*\text{div}_Y(g) =
\sum (V' \times_V U \to U)_*(h')^*\text{div}_{V'}(g \circ b|_{V'})
$$
où $h' : V' \times_V U \to V'$ est la projection. On peut appliquer le lemme
dans le cas des schémas (Homologie de Chow, lemme
\ref{chow-lemma-prepare-flat-pullback-rational-equivalence}) au morphisme $h'
: V' \times_V U \to V'$ pour voir que l'on a
$$
(h')^*\text{div}_{V'}(g \circ b|_{V'}) =
\sum
m_{U', V' \times_V U}
(U' \to V' \times_V U)_*\text{div}_{U'}(g \circ b|_{V'} \circ h'|_{U'})
$$
où la somme est indexée par les composantes irréductibles $U'$ de $V' \times_V U$.
Chaque $U'$ apparaissant dans cette somme est une composante irréductible de
$U$ et inversement chaque composante irréductible $U'$ de $U$ est une
composante irréductible de $V' \times_V U$ pour une composante irréductible
unique $V' \subset V$. Étant donné une composante irréductible $U' \subset U$,
notons $\overline{a(U')} \subset X$ l'« image » dans $X$ (lemme
\ref{lemma-proper-image}) ; c'est une composante irréductible de $X$ par
exemple d'après le lemme \ref{lemma-flat-inverse-image-dimension}. La multiplicité
$m_{U', V' \times_V U}$ est égale à la multiplicité $m_{\overline{a(U')}, X}$.
Cela découle de l'égalité $h^*a^*[Y] = b^*f^*[Y]$ (lemme
\ref{lemma-compose-flat-pullback}), des définitions et du lemme
\ref{lemma-etale-pullback}. En combinant tout ce que nous venons de dire, nous
obtenons
$$
a^*f^*\text{div}_Y(g) =
h^*b^*\text{div}_Y(g) =
\sum m_{\overline{a(U')}, X}
(U' \to U)_*\text{div}_{U'}(g \circ (f \circ a)|_{U'})
$$
À présent, analysons ce que devient le membre de droite de la formule de
l'énoncé après image réciproque par $a$. Tout d'abord, le lemme
\ref{lemma-flat-pullback-proper-pushforward} donne
$$
a^*\sum m_{X', X} (X' \to X)_*\text{div}_{X'}(g \circ f|_{X'}) =
\sum m_{X', X} (X' \times_X U \to U)_*(a')^*\text{div}_{X'}(g \circ f|_{X'})
$$
où $a' : X' \times_X U \to X'$ est la projection. Par le lemme
\ref{lemma-etale-pullback-principal-divisor}, on obtient
$$
(a')^*\text{div}_{X'}(g \circ f|_{X'}) =
\sum (U' \to X' \times_X U)_*\text{div}_{U'}(g \circ (f \circ a)|_{U'})
$$
où la somme est indexée par les composantes irréductibles $U'$ de $X' \times_X U$. Ces
$U'$ sont des composantes irréductibles de $U$ et sont en fait exactement les
composantes irréductibles de $U$ telles que $\overline{a(U')} = X'$. La comparaison
avec ce qui précède permet de conclure.
\end{proof}

\begin{lemma}
\label{lemma-flat-pullback-rational-equivalence}
Dans la situation \ref{situation-setup}, soient $X, Y/B$ bons. Soit $f :
X \to Y$ un morphisme plat de dimension relative $r$. Soit donnée une
équivalence rationnelle $\alpha \sim_{rat} \beta$ entre deux $k$-cycles sur
$Y$. Alors $f^*\alpha
\sim_{rat} f^*\beta$ comme $(k + r)$-cycles sur $X$.
\end{lemma}

\begin{proof}
Que faut-il montrer ? Supposons que l'on se donne une famille
$$
i_j : W_j \longrightarrow Y
$$
d'immersions fermées, chaque $W_j$ étant intègre de $\delta$-dimension $k +
1$, ainsi que des fonctions rationnelles $g_j \in R(W_j)^*$. Supposons en outre que la
famille $\{|i_j|(|W_j|)\}_{j \in J}$ soit localement finie dans $|Y|$. Il
faut alors montrer que
$$
f^*(\sum i_{j, *}\text{div}(g_j)) = \sum f^*i_{j, *}\text{div}(g_j)
$$
est rationnellement équivalent à zéro sur $X$. La somme de droite a un sens
d'après le lemme \ref{lemma-inverse-image-locally-finite}.

\medskip\noindent
Considérons les produits fibrés
$$
i'_j : W'_j = W_j \times_Y X \longrightarrow X.
$$
et notons $f_j : W'_j \to W_j$ la première projection. Par le lemme
\ref{lemma-flat-pullback-proper-pushforward}, nous pouvons écrire la somme
ci-dessus sous la forme
$$
\sum i'_{j, *}(f_j^*\text{div}(g_j))
$$
Par le lemme \ref{lemma-prepare-flat-pullback-rational-equivalence}, on voit
que chaque $f_j^*\text{div}(g_j)$ est rationnellement équivalent à zéro sur
$W'_j$. Par conséquent, chaque $i'_{j, *}(f_j^*\text{div}(g_j))$ est
rationnellement équivalent à zéro. Il en va de même pour la somme affichée par
la discussion dans la remarque
\ref{remark-infinite-sums-rational-equivalences}.
\end{proof}

\begin{lemma}
\label{lemma-proper-pushforward-rational-equivalence}
Dans la situation \ref{situation-setup}, soient $X, Y/B$ bons. Soit $p :
X \to Y$ un morphisme propre. Supposons que $\alpha, \beta \in Z_k(X)$ soient
rationnellement équivalents. Alors $p_*\alpha$ est rationnellement équivalent
à $p_*\beta$.
\end{lemma}

\begin{proof}
Que faut-il montrer ? Supposons que l'on se donne une famille
$$
i_j : W_j \longrightarrow X
$$
d'immersions fermées, chaque $W_j$ étant intègre de $\delta$-dimension $k +
1$, ainsi que des fonctions rationnelles $f_j \in R(W_j)^*$. Supposons en outre
que la famille $\{i_j(W_j)\}_{j \in J}$ soit localement finie dans $X$. Il faut
alors montrer que
$$
p_*\left(\sum i_{j, *}\text{div}(f_j)\right)
$$
est rationnellement équivalent à zéro sur $X$.

\medskip\noindent
Remarquons que la somme est égale à
$$
\sum p_*i_{j, *}\text{div}(f_j).
$$
Soit $W'_j \subset Y$ le sous-espace fermé intègre qui est l'image de $p
\circ i_j$ ; voir le lemme \ref{lemma-proper-image}. La famille $\{W'_j\}$ est
localement finie dans $Y$ par le lemme
\ref{lemma-quasi-compact-locally-finite}. Il suffit donc de montrer, pour un
$j$ donné, soit que $p_*i_{j, *}\text{div}(f_j) = 0$, soit que ce cycle soit égal à
$i'_{j, *}\text{div}(g_j)$ pour un certain $g_j \in R(W'_j)^*$.

\medskip\noindent
Les arguments ci-dessus nous ramènent donc au cas d'un seul sous-espace fermé
intègre $W \subset X$ de $\delta$-dimension $k + 1$. Soit $f \in
R(W)^*$. Soit $W' = p(W)$ comme ci-dessus. On obtient un diagramme commutatif
de morphismes
$$
\xymatrix{
W \ar[r]_i \ar[d]_{p'} & X \ar[d]^p \\
W' \ar[r]^{i'} & Y
}
$$
Remarquons que $p_*i_*\text{div}(f) = i'_*(p')_*\text{div}(f)$ d'après le lemme
\ref{lemma-compose-pushforward}. Comme expliqué ci-dessus, nous devons montrer
que $(p')_*\text{div}(f)$ est le diviseur d'une fonction rationnelle sur $W'$
ou zéro. Il y a trois cas à distinguer.

\medskip\noindent
Le cas $\dim_\delta(W') < k$. Dans ce cas automatiquement $(p')_*\text{div}(f)
= 0$ et il n'y a rien à prouver.

\medskip\noindent
Le cas $\dim_\delta(W') = k$. Montrons que $(p')_*\text{div}(f) = 0$ dans ce
cas. Puisque $(p')_*\text{div}(f)$ est un $k$-cycle, nous voyons que
$(p')_*\text{div}(f) = n[W']$ pour un certain $n \in \mathbf{Z}$. Afin de
prouver que $n = 0$, on peut remplacer $W'$ par un sous-espace ouvert non vide.
En particulier, nous pouvons supposer, et supposons, que $W'$ est
un schéma. Soit $\eta \in W'$ le point générique. Soit $K = \kappa(\eta) =
R(W')$ le corps des fonctions. Considérons le diagramme de changement de base
$$
\xymatrix{
W_\eta \ar[r] \ar[d]_c & W \ar[d]^{p'} \\
\Spec(K) \ar[r]^\eta & W'
}
$$
Remarquons que $c$ est propre. De plus, $|W_\eta|$ est de dimension $1$ : le
lemme \ref{decent-spaces-lemma-topology-fibre} d'Espaces décents identifie
$|W_\eta|$ à la partie de $|W|$ formée des points s'envoyant sur $\eta$ ; or,
puisque $\dim_\delta(W) = k + 1$ et $\delta(\eta) = k$, les points de
$W_\eta$ ont nécessairement une valeur $\delta$ égale à $k$ ou à $k + 1$.
Ainsi, les anneaux locaux sont de dimension $\leq 1$ d'après Espaces décents, lemme
\ref{decent-spaces-lemma-codimension-local-ring}. Par Espaces sur les corps,
lemme \ref{spaces-over-fields-lemma-codim-1-point-in-schematic-locus}, on voit
que $W_\eta$ est un schéma. Puisque $\Spec(K)$ est la limite des
sous-schémas ouverts affines non vides de $W'$, nous concluons que nous
pouvons supposer que $W$ est un schéma d'après Limites d'espaces, lemme
\ref{spaces-limits-lemma-limit-is-scheme}. Enfin, le cas des schémas
(Homologie de Chow, lemme
\ref{chow-lemma-proper-pushforward-rational-equivalence}) donne $n = 0$.

\medskip\noindent
Le cas $\dim_\delta(W') = k + 1$. Dans ce cas, le lemme
\ref{lemma-proper-pushforward-alteration} s'applique, et on voit qu'en effet
$p'_*\text{div}(f) = \text{div}(g)$ pour un certain $g \in R(W')^*$, comme
voulu.
\end{proof}












\section{Le diviseur associé à un faisceau inversible}
\label{section-divisor-invertible-sheaf}

\noindent
Cette section est l'analogue de l'Homologie de Chow, section
\ref{chow-section-divisor-invertible-sheaf}. La définition suivante est
l'analogue d'Espaces sur les corps, définition
\ref{spaces-over-fields-definition-divisor-invertible-sheaf}, dans notre cadre.

\begin{definition}
\label{definition-divisor-invertible-sheaf}
Dans la situation \ref{situation-setup}, soit $X/B$ bon. Supposons que
$X$ soit intègre et que $n = \dim_\delta(X)$. Soit $\mathcal{L}$ un
$\mathcal{O}_X$-module inversible.
\begin{enumerate}
\item Pour toute section méromorphe non nulle $s$ de $\mathcal{L}$, nous
définissons le {\it diviseur de Weil associé à $s$} comme le $(n - 1)$-cycle
$$
\text{div}_\mathcal{L}(s) =
\sum \text{ord}_{Z, \mathcal{L}}(s) [Z]
$$
défini dans Espaces sur les corps, définition
\ref{spaces-over-fields-definition-divisor-invertible-sheaf}. Cela a un sens
car les diviseurs de Weil sont de $\delta$-dimension $n - 1$ par le lemme
\ref{lemma-divisor-delta-dimension}.
\item Nous définissons le {\it diviseur de Weil associé à $\mathcal{L}$} comme
$$
c_1(\mathcal{L}) \cap [X] =
\text{class of }\text{div}_\mathcal{L}(s) \in \CH_{n - 1}(X)
$$
où $s$ est une section méromorphe non nulle quelconque de $\mathcal{L}$ sur
$X$. Cette définition ne dépend pas du choix, d'après Espaces sur les corps, lemme
\ref{spaces-over-fields-lemma-divisor-meromorphic-well-defined}.
\end{enumerate}
\end{definition}

\noindent
Le schéma des zéros d'une section non nulle est un diviseur de Cartier effectif
dont la classe de Weil donne le diviseur de Weil associé au module inversible.

\begin{lemma}
\label{lemma-compute-c1}
Dans la situation \ref{situation-setup}, soit $X/B$ bon. Supposons que
$X$ soit intègre et que $n = \dim_\delta(X)$. Soit $\mathcal{L}$ un
$\mathcal{O}_X$-module inversible. Soit $s \in \Gamma(X, \mathcal{L})$ une
section globale non nulle. Alors
$$
\text{div}_\mathcal{L}(s) = [Z(s)]_{n - 1}
$$
dans $Z_{n - 1}(X)$ et
$$
c_1(\mathcal{L}) \cap [X] = [Z(s)]_{n - 1}
$$
dans $\CH_{n - 1}(X)$.
\end{lemma}

\begin{proof}
Soit $Z \subset X$ un sous-espace fermé intègre de $\delta$-dimension $n -
1$. Soit $\xi \in |Z|$ son point générique. Pour prouver la première égalité,
comparons les coefficients de $Z$ dans les deux membres. Choisissons un
voisinage \'etale élémentaire $(U, u) \to (X, \xi)$ ; voir Espaces décents, section
\ref{decent-spaces-section-residue-fields-henselian-local-rings}, et rappelons
que $\mathcal{O}_{X, \xi}^h = \mathcal{O}_{U, u}^h$ dans ce cas. Après avoir
remplacé $U$ par un voisinage ouvert de $u$, nous pouvons supposer qu'il existe
une section banalisante $s_U$ de $\mathcal{L}|_U$. Écrivons $s|_U = f s_U$ pour
un certain $f \in \Gamma(U, \mathcal{O}_U)$. Alors $Z \times_X U$ est égal à
$V(f)$ comme sous-schéma fermé de $U$ ; voir Diviseurs sur les espaces,
définition \ref{spaces-divisors-definition-zero-scheme-s}. Comme dans Espaces
sur les corps, section \ref{spaces-over-fields-section-c1}, notons
$\mathcal{L}_\xi$ l'image réciproque de $\mathcal{L}$ par le morphisme canonique
$c_\xi : \Spec(\mathcal{O}_{X, \xi}^h) \to X$. Notons $s_\xi$ l'image réciproque de
$s_U$ ; c'est une banalisation de $\mathcal{L}_\xi$. Alors
$c_\xi^*(s) = fs_\xi$. Le coefficient de $Z$ dans $[Z(s)]_{n - 1}$ est par
définition
$$
\text{length}_{\mathcal{O}_{U, u}}(\mathcal{O}_{U, u}/f\mathcal{O}_{U, u})
$$
Puisque $\mathcal{O}_{U, u} \to \mathcal{O}_{X, \xi}^h$ est plat et identifie
les corps résiduels, cela est égal à
$$
\text{length}_{\mathcal{O}_{X, \xi}^h}
(\mathcal{O}_{X, \xi}^h/f\mathcal{O}_{X, \xi}^h)
$$
par Algèbre, lemme \ref{algebra-lemma-pullback-module}. Cette dernière quantité
est égale à $\text{ord}_{Z, \mathcal{L}}(s)$ par Espaces sur les corps,
définition \ref{spaces-over-fields-definition-order-vanishing-meromorphic},
c'est-à-dire au coefficient de $Z$ dans $\text{div}_\mathcal{L}(s)$ comme
voulu.
\end{proof}

\begin{lemma}
\label{lemma-Gm-torsor}
Dans la situation \ref{situation-setup}, soit $X/B$ bon. Soit
$\mathcal{L}$ un $\mathcal{O}_X$-module inversible. Le morphisme
$$
q :
T = \underline{\Spec}\left(
\bigoplus\nolimits_{n \in \mathbf{Z}} \mathcal{L}^{\otimes n}\right)
\longrightarrow
X
$$
a les propriétés suivantes :
\begin{enumerate}
\item $q$ est surjectif, lisse, affine, de dimension relative $1$,
\item il existe un isomorphisme $\alpha : q^*\mathcal{L} \cong \mathcal{O}_T$,
\item la formation de $(q : T \to X, \alpha)$ commute aux changements de
base,
\item $q^* : Z_k(X) \to Z_{k + 1}(T)$ est injectif,
\item si $Z \subset X$ est un sous-espace fermé intègre, alors $q^{-1}(Z)
\subset T$ est un sous-espace fermé intègre,
\item si $Z \subset X$ est un sous-espace fermé de $X$ de $\delta$-dimension
$\leq k$, alors $q^{-1}(Z)$ est un sous-espace fermé de $T$ de
$\delta$-dimension $\leq k + 1$ et $q^*[Z]_k = [q^{-1}(Z)]_{k + 1}$,
\item si $\xi' \in |T|$ est le point générique de la fibre de $|T| \to |X|$
sur $\xi$, alors l'homomorphisme d'anneaux $\mathcal{O}_{X, \xi}^h \to
\mathcal{O}_{T, \xi'}^h$ est plat, nous avons $\mathfrak m_{\xi'}^h =
\mathfrak m_\xi^h \mathcal{O}_{T, \xi'}^h$, et l'extension de corps résiduels
est purement transcendante, de degré de transcendance $1$, et
\item ajouter ici d'autres propriétés au besoin.
\end{enumerate}
\end{lemma}

\begin{proof}
Soit $U \to X$ un morphisme \'etale tel que $\mathcal{L}|_U$ soit trivial.
Alors $T \times_X U \to U$ est isomorphe au morphisme de projection
$\mathbf{G}_m \times U \to U$, où $\mathbf{G}_m$ est le schéma en groupes
multiplicatif ; voir Groupoïdes, exemple
\ref{groupoids-example-multiplicative-group}. Ainsi (1) est clair.

\medskip\noindent
Pour (2), remarquons que $q_*q^*\mathcal{L} = \bigoplus_{n \in \mathbf{Z}}
\mathcal{L}^{\otimes n + 1}$. Il existe donc un isomorphisme évident
$q_*q^*\mathcal{L} \to q_*\mathcal{O}_T$ de $q_*\mathcal{O}_T$-modules. D'après
Morphismes d'espaces, lemme
\ref{spaces-morphisms-lemma-affine-equivalence-modules}, cela détermine un
isomorphisme $q^*\mathcal{L} \to \mathcal{O}_T$.

\medskip\noindent
La partie (3) résulte du fait que la formation du spectre relatif commute à
tout changement de base ; il en va clairement de même pour
l'isomorphisme $\alpha$.

\medskip\noindent
La partie (4) découle immédiatement de (1) et des définitions.

\medskip\noindent
La partie (5) découle du fait que si $Z$ est un espace algébrique intègre,
alors $\mathbf{G}_m \times Z$ est un espace algébrique intègre.

\medskip\noindent
La partie (6) découle de la conservation des longueurs : si $(A,
\mathfrak m)$ est un anneau local, si $B = A[x]_{\mathfrak m A[x]}$ et si $M$
est un $A$-module, alors $\text{length}_A(M) = \text{length}_B(M \otimes_A
B)$. Cela implique que, si $\mathcal{F}$ est un $\mathcal{O}_X$-module cohérent,
si $\xi \in |X|$ et si $\xi' \in |T|$ est le point générique de la fibre
au-dessus de $\xi$, alors la longueur de $\mathcal{F}$ en $\xi$ est égale à celle de
$q^*\mathcal{F}$ en $\xi'$. En remontant les définitions, on obtient (6), et
même davantage.

\medskip\noindent
L'homomorphisme de la partie (7) provient d'Espaces décents, remarque
\ref{decent-spaces-remark-functoriality-henselian-local-ring}. Cependant, dans
notre cas, nous avons
$$
\Spec(\mathcal{O}_{X, \xi}^h) \times_X T =
\mathbf{G}_m \times \Spec(\mathcal{O}_{X, \xi}^h) =
\Spec(\mathcal{O}_{X, \xi}^h[t, t^{-1}])
$$
et $\xi'$ correspond au point générique de la fibre spéciale de ce morphisme au-dessus de
$\Spec(\mathcal{O}_{X, \xi}^h)$. Ainsi $\mathcal{O}_{T, \xi'}^h$ est
l'hensélisation de la localisation de $\mathcal{O}_{X, \xi}^h[t, t^{-1}]$ en
l'idéal premier correspondant. La partie (7) en découle par quelques résultats
d'algèbre commutative ; détails omis.
\end{proof}

\begin{lemma}
\label{lemma-Gm-torsor-divisor-meromorphic-section}
Dans la situation \ref{situation-setup}, soit $X/B$ bon. Soit
$\mathcal{L}$ un $\mathcal{O}_X$-module inversible. Supposons que $X$ soit
intègre. Soit $s$ une section méromorphe non nulle de $\mathcal{L}$. Soit $q
: T \to X$ le morphisme du lemme \ref{lemma-Gm-torsor}. Alors
$$
q^*\text{div}_\mathcal{L}(s) = \text{div}_T(q^*(s))
$$
où nous considérons l'image réciproque $q^*(s)$ comme une fonction méromorphe non
nulle sur $T$ au moyen de l'isomorphisme $q^*\mathcal{L} \to \mathcal{O}_T$.
\end{lemma}

\begin{proof}
Remarquons que $\text{div}_T(q^*(s)) = \text{div}_{\mathcal{O}_T}(q^*(s))$ par
la compatibilité entre les constructions données dans Espaces sur les corps,
sections \ref{spaces-over-fields-section-Weil-divisors} et
\ref{spaces-over-fields-section-c1}. Dans cette preuve, nous établirons
l'égalité avec $\text{div}_{\mathcal{O}_T}(q^*(s))$. Nous utiliserons
toutes les propriétés de $q : T \to X$ indiquées dans le lemme
\ref{lemma-Gm-torsor} sans autre mention. Soit $Z \subset T$ un diviseur
premier. Ou bien $Z \to X$ est dominant, ou bien $Z = q^{-1}(Z')$ pour un
diviseur premier $Z' \subset X$. Si $Z \to X$ est dominant, alors le
coefficient de $Z$ dans chacun des deux membres de l'égalité du lemme est nul. On peut
donc supposer $Z = q^{-1}(Z')$ où $Z' \subset X$ est un diviseur premier. Soit
$\xi' \in |Z'|$ et $\xi \in |Z|$ les points génériques. On obtient alors un
diagramme commutatif
$$
\xymatrix{
\Spec(\mathcal{O}_{T, \xi}^h) \ar[r]_-{c_\xi} \ar[d]_h & T \ar[d]^q \\
\Spec(\mathcal{O}_{X, \xi'}^h) \ar[r]^-{c_{\xi'}} & X
}
$$
voir Espaces décents, remarque
\ref{decent-spaces-remark-functoriality-henselian-local-ring}. Choisissons une
banalisation $s_{\xi'}$ de $\mathcal{L}_{\xi'} = c_{\xi'}^*\mathcal{L}$.
Nous pouvons alors utiliser l'image réciproque $s_\xi$ de $s_{\xi'}$ par $h$
comme banalisation de $\mathcal{L}_\xi = c_\xi^* q^*\mathcal{L}$. Écrivons
$s/s_{\xi'} = a/b$ avec $a, b \in \mathcal{O}_{X, \xi'}$ non-diviseurs de
zéro. Par définition, le coefficient de $Z'$ dans $\text{div}_\mathcal{L}(s)$
est
$$
\text{length}_{\mathcal{O}_{X, \xi'}^h}(
\mathcal{O}_{X, \xi'}^h/a \mathcal{O}_{X, \xi'}^h)
-
\text{length}_{\mathcal{O}_{X, \xi'}^h}(
\mathcal{O}_{X, \xi'}^h/b \mathcal{O}_{X, \xi'}^h)
$$
Puisque $Z = q^{-1}(Z')$, c'est aussi le coefficient de $Z$ dans
$q^*\text{div}_\mathcal{L}(s)$. Étant donné que $\mathcal{O}_{X, \xi'}^h \to
\mathcal{O}_{T, \xi}^h$ est plat, les éléments $a, b$ sont envoyés sur des
non-diviseurs de zéro dans $\mathcal{O}_{T, \xi}^h$. Ainsi $q^*(s)/s_\xi = a/b$
dans l'anneau total des fractions de $\mathcal{O}_{T, \xi}^h$. Par définition, le
coefficient de $Z$ dans $\text{div}_T(q^*(s))$ est
$$
\text{length}_{\mathcal{O}_{T, \xi}^h}(
\mathcal{O}_{T, \xi}^h/a \mathcal{O}_{T, \xi}^h)
-
\text{length}_{\mathcal{O}_{T, \xi}^h}(
\mathcal{O}_{T, \xi}^h/b \mathcal{O}_{T, \xi}^h)
$$
La preuve est achevée, car ces longueurs sont les mêmes que précédemment
d'après Algèbre, lemme \ref{algebra-lemma-pullback-module}, et l'égalité
$\mathfrak m_\xi^h = \mathfrak m_{\xi'}^h\mathcal{O}_{T, \xi}^h$ établie dans
le lemme \ref{lemma-Gm-torsor}.
\end{proof}









\section{Intersection avec un faisceau inversible}
\label{section-intersecting-with-divisors}

\noindent
Cette section est l'analogue de l'Homologie de Chow, section
\ref{chow-section-intersecting-with-divisors}. Dans cette section, nous
étudions la construction suivante.

\begin{definition}
\label{definition-cap-c1}
Dans la situation \ref{situation-setup}, soit $X/B$ bon. Soit
$\mathcal{L}$ un $\mathcal{O}_X$-module inversible. On définit, pour tout
entier $k$, une opération
$$
c_1(\mathcal{L}) \cap - :
Z_{k + 1}(X) \to \CH_k(X)
$$
appelée {\it intersection avec la première classe de Chern de $\mathcal{L}$}.
\begin{enumerate}
\item Étant donné un sous-espace fermé intègre $i : W \to X$ avec
$\dim_\delta(W) = k + 1$, nous définissons
$$
c_1(\mathcal{L}) \cap [W] = i_*(c_1({i^*\mathcal{L}}) \cap [W])
$$
où le membre de droite est défini dans la définition
\ref{definition-divisor-invertible-sheaf}.
\item Pour un $(k + 1)$-cycle général $\alpha = \sum n_i [W_i]$, nous
définissons
$$
c_1(\mathcal{L}) \cap \alpha = \sum n_i c_1(\mathcal{L}) \cap [W_i]
$$
\end{enumerate}
\end{definition}

\noindent
Écrivons chaque $c_1(\mathcal{L}) \cap W_i = \sum_j n_{i, j} [Z_{i, j}]$, où
$\{Z_{i, j}\}_j$ est une famille localement finie de sous-espaces fermés intègres
de $W_i$. Puisque $\{W_i\}$ est une famille localement finie de
sous-espaces fermés intègres de $X$, il s'ensuit facilement que $\{Z_{i,
j}\}_{i, j}$ est une famille localement finie de sous-espaces fermés de
$X$. $c_1(\mathcal{L}) \cap \alpha = \sum n_in_{i, j}[Z_{i, j}]$ est donc un
cycle. Une autre interprétation, souvent plus commode, consiste à observer
que le morphisme $\coprod W_i \to X$ est propre. Par conséquent,
$c_1(\mathcal{L}) \cap \alpha$ peut être considéré comme l’image directe d’une
classe dans $\CH_k(\coprod W_i) = \prod \CH_k(W_i)$. Cela explique aussi
pourquoi le résultat est bien défini modulo l'équivalence rationnelle sur $X$.

\medskip\noindent
L'objectif principal des prochaines sections est de montrer que l'intersection
par $c_1(\mathcal{L})$ passe au quotient par l'équivalence rationnelle. Cela
n'est pas immédiat.

\begin{lemma}
\label{lemma-c1-cap-additive}
Dans la situation \ref{situation-setup}, soit $X/B$ bon. Soient
$\mathcal{L}$, $\mathcal{N}$ des faisceaux inversibles sur $X$. Alors
$$
c_1(\mathcal{L}) \cap \alpha  + c_1(\mathcal{N}) \cap \alpha =
c_1(\mathcal{L} \otimes_{\mathcal{O}_X} \mathcal{N}) \cap \alpha
$$
dans $\CH_k(X)$ pour tout $\alpha \in Z_{k - 1}(X)$. De plus,
$c_1(\mathcal{O}_X) \cap \alpha = 0$ pour tout $\alpha$.
\end{lemma}

\begin{proof}
L'additivité découle directement d'Espaces sur les corps, lemme
\ref{spaces-over-fields-lemma-c1-additive} et des définitions. Pour voir que
$c_1(\mathcal{O}_X) \cap \alpha = 0$, considérons la section $1 \in \Gamma(X,
\mathcal{O}_X)$. Celle-ci se restreint en une section partout non nulle sur tout
sous-espace fermé intègre $W \subset X$. D'où $c_1(\mathcal{O}_X) \cap [W] =
0$, comme voulu.
\end{proof}

\noindent
Rappelons que $Z(s) \subset X$ désigne le schéma des zéros d'une section globale $s$
d'un faisceau inversible sur un espace algébrique $X$, voir Diviseurs sur les
espaces, définition \ref{spaces-divisors-definition-zero-scheme-s}.

\begin{lemma}
\label{lemma-prepare-geometric-cap}
Dans la situation \ref{situation-setup}, soit $Y/B$ bon. Soit
$\mathcal{L}$ un $\mathcal{O}_Y$-module inversible. Soit $s \in \Gamma(Y,
\mathcal{L})$ une section régulière et supposons $\dim_\delta(Y) \leq k + 1$.
Écrivons $[Y]_{k + 1} = \sum n_i[Y_i]$ où $Y_i \subset Y$ sont les composantes
irréductibles de $Y$ de $\delta$-dimension $k + 1$. Posons $s_i = s|_{Y_i}
\in \Gamma(Y_i, \mathcal{L}|_{Y_i})$. Alors
\begin{equation}
\label{equation-equal-as-cycles}
[Z(s)]_k =  \sum n_i[Z(s_i)]_k
\end{equation}
comme $k$-cycles sur $Y$.
\end{lemma}

\begin{proof}
Soit $\varphi : V \to Y$ un morphisme \'etale surjectif, où $V$ est un schéma.
Il suffit de démontrer l'égalité après image réciproque par $\varphi$ ; voir le lemme
\ref{lemma-etale-pullback}. Ce même lemme nous dit que $\varphi^*[Y_i] =
[\varphi^{-1}(Y_i)] = \sum [V_{i, j}]$ où $V_{i, j}$ sont les composantes
irréductibles de $V$ situées au-dessus de $Y_i$. Par conséquent, si nous
appliquons d'abord le cas des schémas (Homologie de Chow, lemme
\ref{chow-lemma-prepare-geometric-cap}) à $\varphi^*s_i$ sur $Y_i \times_Y V$,
nous trouvons que $\varphi^*[Z(s_i)]_k = [Z(\varphi^*s_i)] = \sum [Z(s_{i,
j})]_k$ où $s_{i, j}$ est l'image réciproque de $s$ sur $V_{i, j}$. En
appliquant le cas des schémas à $\varphi^*s$, on obtient
$$
\varphi^*[Z(s)]_k =
[Z(\varphi^*s)]_k =
\sum n_i[Z(s_{i, j})]_k
$$
par notre remarque sur les multiplicités ci-dessus. En combinant tout ce qui
précède, la preuve est complète.
\end{proof}

\noindent
Le lemme suivant permet de calculer le produit d'intersection de la classe
$c_1$ d'un faisceau inversible avec le cycle associé à un sous-schéma fermé.
Rappelons que $Z(s) \subset X$ désigne le schéma des zéros d'une section globale $s$
d'un faisceau inversible sur un schéma $X$ ; voir Diviseurs, définition
\ref{divisors-definition-zero-scheme-s}.

\begin{lemma}
\label{lemma-geometric-cap}
Dans la situation \ref{situation-setup}, soit $X/B$ bon. Soit
$\mathcal{L}$ un $\mathcal{O}_X$-module inversible. Soit $Y \subset X$ un
sous-schéma fermé avec $\dim_\delta(Y) \leq k + 1$, et soit $s \in \Gamma(Y,
\mathcal{L}|_Y)$ une section régulière. Alors
$$
c_1(\mathcal{L}) \cap [Y]_{k + 1} = [Z(s)]_k
$$
dans $\CH_k(X)$.
\end{lemma}

\begin{proof}
Écrivons
$$
[Y]_{k + 1} = \sum n_i[Y_i]
$$
où $Y_i \subset Y$ sont les composantes irréductibles de $Y$ de $\delta$-dimension $k + 1$ et $n_i > 0$. Par hypothèse, la restriction $s_i = s|_{Y_i}
\in \Gamma(Y_i, \mathcal{L}|_{Y_i})$ n'est pas nulle et constitue donc une
section régulière. D'après le lemme \ref{lemma-compute-c1},
$[Z(s_i)]_k$ représente $c_1(\mathcal{L}|_{Y_i})$. Par définition,
$$
c_1(\mathcal{L}) \cap [Y]_{k + 1} = \sum n_i[Z(s_i)]_k
$$
Le résultat découle donc du lemme \ref{lemma-prepare-geometric-cap}.
\end{proof}









\section{Intersection avec un faisceau inversible, image directe et image réciproque}
\label{section-intersecting-with-divisors-push-pull}

\noindent
Cette section est l'analogue de l'Homologie de Chow, section
\ref{chow-section-intersecting-with-divisors-push-pull}. Dans cette section,
nous prouvons que l'opération $c_1(\mathcal{L}) \cap -$ commute avec
l'image réciproque plate et l'image directe propre.

\begin{lemma}
\label{lemma-prepare-flat-pullback-cap-c1}
Dans la situation \ref{situation-setup}, soient $X, Y/B$ bons. Soit $f :
X \to Y$ un morphisme plat de dimension relative $r$. Soit $\mathcal{L}$ un
faisceau inversible sur $Y$. Supposons $Y$ intègre et posons $n =
\dim_\delta(Y)$. Soit $s$ une section méromorphe non nulle de $\mathcal{L}$.
Alors
$$
f^*\text{div}_\mathcal{L}(s) = \sum n_i\text{div}_{f^*\mathcal{L}|_{X_i}}(s_i)
$$
dans $Z_{n + r - 1}(X)$. Ici, la somme porte sur les composantes irréductibles
$X_i \subset X$ de $\delta$-dimension $n + r$, la section $s_i =
f|_{X_i}^*(s)$ est l'image réciproque de $s$ et $n_i = m_{X_i, X}$ est la multiplicité
de $X_i$ dans $X$.
\end{lemma}

\begin{proof}
Par une astuce, nous allons déduire cet énoncé du lemme
\ref{lemma-prepare-flat-pullback-rational-equivalence}. (Une alternative est
de refaire la preuve de ce lemme dans le cadre de sections méromorphes de
modules inversibles.) Plus précisément, soit $q : T \to Y$ le morphisme du lemme
\ref{lemma-Gm-torsor} construit en utilisant $\mathcal{L}$ sur $Y$. Nous
utiliserons toutes les propriétés de $T$ énoncées dans ce lemme. Considérons le
diagramme cartésien
$$
\xymatrix{
T' \ar[r]_{q'} \ar[d]_h & X \ar[d]^f \\
T \ar[r]^q & Y
}
$$
Alors $q' : T' \to X$ est le morphisme construit en utilisant $f^*\mathcal{L}$
sur $X$. Il suffit alors de démontrer
$$
(q')^*f^*\text{div}_\mathcal{L}(s) =
\sum n_i (q')^*\text{div}_{f^*\mathcal{L}|_{X_i}}(s_i)
$$
Remarquons que $T'_i = q^{-1}(X_i)$ sont les composantes irréductibles de $T'$
et que $n_i$ est la multiplicité de $T'_i$ dans $T'$. Le membre de gauche est
égal à
$$
h^*q^*\text{div}_\mathcal{L}(s) = h^*\text{div}_T(q^*(s))
$$
par le lemme \ref{lemma-Gm-torsor-divisor-meromorphic-section} (et le lemme
\ref{lemma-compose-flat-pullback}). D'autre part, si $q'_i : T'_i \to
X_i$ désigne la restriction de $q'$, nous constatons que le lemme
\ref{lemma-Gm-torsor-divisor-meromorphic-section} nous dit également que le
membre de droite est égal à
$$
\sum n_i \text{div}_{T_i}((q'_i)^*(s_i))
$$
Dans ces deux formules, les expressions $q^*(s)$ et $(q'_i)^*(s_i)$
représentent les fonctions rationnelles correspondant aux sections méromorphes
de $q^*\mathcal{L}$ et de $(q'_i)^*f^*\mathcal{L}|_{X_i}$ obtenues par image
réciproque ; l'identification se fait au moyen de l'isomorphisme
$\alpha : q^*\mathcal{L} \to \mathcal{O}_T$ et de ses images réciproques sur les
espaces au-dessus de $T$. Avec cette convention, il est clair que
$(q'_i)^*(s_i)$ est la composition de la fonction rationnelle $q^*(s)$ sur $T$
et du morphisme $h|_{T'_i} : T'_i \to T$. Ainsi le lemme
\ref{lemma-prepare-flat-pullback-rational-equivalence} dit exactement que
$$
h^*\text{div}_T(q^*(s)) = \sum n_i \text{div}_{T_i}((q'_i)^*(s_i))
$$
comme voulu.
\end{proof}

\begin{lemma}
\label{lemma-flat-pullback-cap-c1}
Dans la situation \ref{situation-setup}, soient $X, Y/B$ bons. Soit $f :
X \to Y$ un morphisme plat de dimension relative $r$. Soit $\mathcal{L}$ un
faisceau inversible sur $Y$. Soit $\alpha$ un $k$-cycle sur $Y$. Alors
$$
f^*(c_1(\mathcal{L}) \cap \alpha) = c_1(f^*\mathcal{L}) \cap f^*\alpha
$$
dans $\CH_{k + r - 1}(X)$.
\end{lemma}

\begin{proof}
Écrivons $\alpha = \sum n_i[W_i]$. Nous allons montrer que
$$
f^*(c_1(\mathcal{L}) \cap [W_i]) = c_1(f^*\mathcal{L}) \cap f^*[W_i]
$$
dans $\CH_{k + r - 1}(X)$ en produisant une équivalence rationnelle sur le
sous-espace fermé $f^{-1}(W_i)$ de $X$. D'après la discussion de la remarque
\ref{remark-infinite-sums-rational-equivalences}, cela prouvera que l'égalité
du lemme est vraie.

\medskip\noindent
Soit $W \subset Y$ un sous-espace fermé intègre de $\delta$-dimension $k$.
Considérons le sous-espace fermé $W' = f^{-1}(W) = W \times_Y X$, ce qui donne
le diagramme cartésien
$$
\xymatrix{
W' \ar[r] \ar[d]_h & X \ar[d]^f \\
W \ar[r] & Y
}
$$
Nous devons montrer que $f^*(c_1(\mathcal{L}) \cap [W]) = c_1(f^*\mathcal{L})
\cap f^*[W]$. Choisissons une section méromorphe non nulle $s$ de
$\mathcal{L}|_W$. Soient $W'_i \subset W'$ les composantes irréductibles de
$\delta$-dimension $k + r$. Écrivons $[W']_{k + r} = \sum n_i[W'_i]$ avec $n_i$
la multiplicité de $W'_i$ dans $W'$ selon la définition. Donc $f^*[W] = \sum
n_i[W'_i]$ dans $Z_{k + r}(X)$. Puisque chaque $W'_i \to W$ est dominant, nous
voyons que $s_i = s|_{W'_i}$ est une section méromorphe non nulle pour tout
$i$. Par le lemme \ref{lemma-prepare-flat-pullback-cap-c1}, nous avons
l'égalité de cycles suivante
$$
h^*\text{div}_{\mathcal{L}|_W}(s) =
\sum n_i\text{div}_{f^*\mathcal{L}|_{W'_i}}(s_i)
$$
dans $Z_{k + r - 1}(W')$. Ceci achève la preuve, puisque le membre de gauche est un
cycle sur $W'$ dont l'image directe est $f^*(c_1(\mathcal{L}) \cap [W])$ dans $\CH_{k +
r - 1}(X)$ et le membre de droite est un cycle sur $W'$ dont l'image directe est
$c_1(f^*\mathcal{L}) \cap f^*[W]$ dans $\CH_{k + r - 1}(X)$.
\end{proof}

\begin{lemma}
\label{lemma-equal-c1-as-cycles}
Dans la situation \ref{situation-setup}, soient $X, Y/B$ bons. Soit $f :
X \to Y$ un morphisme propre. Soit $\mathcal{L}$ un faisceau inversible sur $Y$.
Supposons $X$, $Y$ intègres, $f$ dominant et $\dim_\delta(X) =
\dim_\delta(Y)$. Soit $s$ une section méromorphe non nulle, la notation $s$ désignant une section de
$\mathcal{L}$ sur $Y$. Alors
$$
f_*\left(\text{div}_{f^*\mathcal{L}}(f^*s)\right) =
[R(X) : R(Y)]\text{div}_\mathcal{L}(s).
$$
Il s'agit d'une égalité de cycles sur $Y$. En particulier,
$$
f_*(c_1(f^*\mathcal{L}) \cap [X]) = c_1(\mathcal{L}) \cap f_*[Y].
$$
\end{lemma}

\begin{proof}
La dernière égalité découle de la première puisque $f_*[X] = [R(X) :
R(Y)][Y]$ par définition. Démontrons la première égalité. Soit $q : T \to Y$ le
morphisme du lemme \ref{lemma-Gm-torsor} construit en utilisant $\mathcal{L}$
sur $Y$. Nous utiliserons toutes les propriétés de $T$ énoncées dans ce lemme.
Considérons le diagramme cartésien
$$
\xymatrix{
T' \ar[r]_{q'} \ar[d]_h & X \ar[d]^f \\
T \ar[r]^q & Y
}
$$
Alors $q' : T' \to X$ est le morphisme construit en utilisant $f^*\mathcal{L}$
sur $X$. Il suffit de démontrer l'égalité après image réciproque sur $T'$. L'image
réciproque du membre de gauche vaut
\begin{align*}
q^*f_*\left(\text{div}_{f^*\mathcal{L}}(f^*s)\right)
& =
h_*(q')^*\text{div}_{f^*\mathcal{L}}(f^*s) \\
& =
h_*\text{div}_{(q')^*f^*\mathcal{L}}((q')^*f^*s) \\
& =
h_*\text{div}_{h^*q^*\mathcal{L}}(h^*q^*s)
\end{align*}
La première égalité découle du lemme \ref{lemma-flat-pullback-proper-pushforward}.
La deuxième découle du lemme \ref{lemma-prepare-flat-pullback-cap-c1}, puisque
$T'$ est intègre. La troisième résulte de la commutativité du diagramme affiché.
L'image réciproque du membre de droite est
$$
[R(X) : R(Y)]q^*\text{div}_\mathcal{L}(s) =
[R(T') : R(T)]\text{div}_{q^*\mathcal{L}}(q^*s)
$$
Cela découle du lemme \ref{lemma-prepare-flat-pullback-cap-c1}, du fait que
$T$ est intègre, et de l'égalité $[R(T') : R(T)] = [R(X) : R(Y)]$ dont nous
omettons la preuve (cela découle du lemme
\ref{lemma-flat-pullback-proper-pushforward}, mais ce serait une manière
inutilement détournée de démontrer l'égalité). Il suffit donc de démontrer le lemme pour $h : T' \to T$,
le module inversible $q^\mathcal{L}$ et la section $q^*s$. Puisque
$q^*\mathcal{L} = \mathcal{O}_T$, nous nous ramenons au cas où $\mathcal{L} \cong
\mathcal{O}$ discuté dans le paragraphe suivant.

\medskip\noindent
Supposons que $\mathcal{L} = \mathcal{O}_Y$. Dans ce cas $s$ correspond à une
fonction rationnelle $g \in R(Y)$. En utilisant le plongement $R(Y) \subset
R(X)$, nous pouvons considérer $g$ comme une fonction rationnelle sur $X$ et il
s'agit simplement de démontrer
$$
f_*\left(\text{div}_X(g)\right) = [R(X) : R(Y)]\text{div}_Y(g).
$$
En comparant avec le résultat du lemme
\ref{lemma-proper-pushforward-alteration}, nous voyons que cela est vrai
puisque $\text{Nm}_{R(X)/R(Y)}(g) = g^{[R(X) : R(Y)]}$ pour $g \in R(Y)^*$.
\end{proof}

\begin{lemma}
\label{lemma-pushforward-cap-c1}
Dans la situation \ref{situation-setup}, soient $X, Y/B$ bons. Soit $p :
X \to Y$ un morphisme propre. Soit $\alpha \in Z_{k + 1}(X)$. Soit
$\mathcal{L}$ un faisceau inversible sur $Y$. Alors
$$
p_*(c_1(p^*\mathcal{L}) \cap \alpha) = c_1(\mathcal{L}) \cap p_*\alpha
$$
dans $\CH_k(Y)$.
\end{lemma}

\begin{proof}
Supposons que $p$ ait la propriété que, pour tout sous-espace fermé intègre
$W \subset X$, l'application $p|_W : W \to Y$ est une immersion fermée.
Alors, par définition de l'intersection avec $c_1(\mathcal{L})$, le lemme est
valable.

\medskip\noindent
Nous allons utiliser cette remarque pour nous ramener à un cas particulier. Plus précisément,
écrivons $\alpha = \sum n_i[W_i]$ avec $n_i \not = 0$ et $W_i$ distincts deux à
deux. Soit $W'_i \subset Y$ l'« image » de $W_i$ comme dans le lemme
\ref{lemma-proper-image}. Considérons le diagramme
$$
\xymatrix{
X' = \coprod W_i \ar[r]_-q \ar[d]_{p'} & X \ar[d]^p \\
Y' = \coprod W'_i \ar[r]^-{q'} & Y.
}
$$
Puisque $\{W_i\}$ est localement fini sur $X$ et que $p$ est propre, on voit
que $\{W'_i\}$ est localement fini sur $Y$ et que $q, q', p'$ sont eux aussi
des morphismes propres. Nous pouvons également considérer $\sum n_i[W_i]$
comme un $k$-cycle $\alpha' \in Z_k(X')$. Clairement, $q_*\alpha' = \alpha$.
Nous avons $q_*(c_1(q^*p^*\mathcal{L}) \cap \alpha') = c_1(p^*\mathcal{L})
\cap q_*\alpha'$ et $(q')_*(c_1((q')^*\mathcal{L}) \cap p'_*\alpha') =
c_1(\mathcal{L}) \cap q'_*p'_*\alpha'$ par la remarque initiale de la preuve.
Il suffit donc de démontrer le lemme pour le morphisme $p'$ et le cycle $\sum
n_i[W_i]$. Cela signifie que nous pouvons supposer $X$, $Y$
intègres, $f : X \to Y$ dominant et $\alpha = [X]$. Dans ce cas, le résultat
découle du lemme \ref{lemma-equal-c1-as-cycles}.
\end{proof}













\section{La formule clé}
\label{section-key}

\noindent
Cette section est l'analogue de l'Homologie de Chow, section
\ref{chow-section-key}. Nous recommandons vivement au lecteur de lire d'abord
la preuve dans ce cas.

\medskip\noindent
Dans la situation \ref{situation-setup}, soit $X/B$ bon. Supposons
$X$ intègre et $\dim_\delta(X) = n$. Soient $\mathcal{L}$ et
$\mathcal{N}$ des $\mathcal{O}_X$-modules inversibles. Soit $s$ une section
méromorphe non nulle de $\mathcal{L}$ et soit $t$ une section méromorphe non
nulle de $\mathcal{N}$. Soit $Z \subset X$ un diviseur premier ayant pour point
générique $\xi \in |Z|$. Considérons le morphisme
$$
c_\xi : \Spec(\mathcal{O}_{X, \xi}^h) \longrightarrow X
$$
utilisé dans Espaces sur les corps, section \ref{spaces-over-fields-section-c1}.
Nous notons $\mathcal{L}_\xi$ et $\mathcal{N}_\xi$ les images réciproques de
$\mathcal{L}$ et $\mathcal{N}$ par $c_\xi$ ; nous considérons souvent
$\mathcal{L}_\xi$ et $\mathcal{N}_\xi$ comme
les modules libres de rang $1$ sur $\mathcal{O}_{X, \xi}^h$ qu'ils définissent.
Notons que l'image réciproque de $s$, resp.\ $t$, est une section méromorphe régulière de $\mathcal{L}_\xi$,
resp.\ $\mathcal{N}_\xi$.

\medskip\noindent
Soit $Z_i \subset X$, $i \in I$, une famille localement finie de diviseurs
premiers ayant la propriété suivante : si $Z \not \in \{Z_i\}$, alors $s$ est
un générateur de $\mathcal{L}_\xi$ et $t$ est un générateur de
$\mathcal{N}_\xi$. Une telle famille existe d'après le lemme d'Espaces sur les corps
\ref{spaces-over-fields-lemma-divisor-meromorphic-locally-finite}. Alors
$$
\text{div}_\mathcal{L}(s) = \sum \text{ord}_{Z_i, \mathcal{L}}(s) [Z_i]
$$
et de même
$$
\text{div}_\mathcal{N}(t) = \sum \text{ord}_{Z_i, \mathcal{N}}(t) [Z_i]
$$
En explicitant davantage les définitions, choisissons pour tout $i$ des générateurs
$s_i \in \mathcal{L}_{\xi_i}$ et $t_i \in \mathcal{N}_{\xi_i}$,
où $\xi_i$ est le point générique de $Z_i$. Nous pouvons alors écrire
$$
s = f_i s_i
\quad\text{et}\quad
t = g_i t_i
$$
où $f_i, g_i$ sont des éléments inversibles de l'anneau total des fractions
$Q(\mathcal{O}_{X, \xi_i}^h)$. Posons, pour abréger, $B_i = \mathcal{O}_{X,
\xi_i}^h$. Notons
$$
\text{ord}_{B_i} : Q(B_i)^* \longrightarrow \mathbf{Z},\quad
a/b \longmapsto
\text{length}_{B_i}(B_i/aB_i) - \text{length}_{B_i}(B_i/bB_i)
$$
Autrement dit, nous étendons provisoirement la définition
\ref{algebra-definition-ord} d'Algèbre à ces anneaux locaux noethériens réduits
de dimension $1$. Alors, par définition,
$$
\text{ord}_{Z_i, \mathcal{L}}(s) = \text{ord}_{B_i}(f_i)
\quad\text{et}\quad
\text{ord}_{Z_i, \mathcal{N}}(t) = \text{ord}_{B_i}(g_i)
$$
Puisque $\xi_i$ est le point générique de $Z_i$, nous voyons que le corps
résiduel $\kappa(\xi_i)$ est le corps des fonctions de $Z_i$. De plus,
$\kappa(\xi_i)$ est le corps résiduel de $B_i$ ; voir le lemme d'Espaces décents
\ref{decent-spaces-lemma-residue-field-henselian-local-ring}. Puisque $t_i$
est un générateur de $\mathcal{N}_{\xi_i}$, on voit que son image dans la fibre
$\mathcal{N}_{\xi_i} \otimes_{B_i} \kappa(\xi_i)$ est une section méromorphe
non nulle de $\mathcal{N}|_{Z_i}$. Nous noterons cette image $t_i|_{Z_i}$. Il
résulte de nos définitions que
$$
c_1(\mathcal{N}) \cap \text{div}_\mathcal{L}(s) =
\sum \text{ord}_{B_i}(f_i)
(Z_i \to X)_*\text{div}_{\mathcal{N}|_{Z_i}}(t_i|_{Z_i})
$$
et de même
$$
c_1(\mathcal{L}) \cap \text{div}_\mathcal{N}(t) =
\sum \text{ord}_{B_i}(g_i)
(Z_i \to X)_*\text{div}_{\mathcal{L}|_{Z_i}}(s_i|_{Z_i})
$$
dans $\CH_{n - 2}(X)$. Nous allons construire une équivalence rationnelle entre
ces deux cycles. Pour cela, considérons le symbole modéré
$$
\partial_{B_i}(f_i, g_i) \in \kappa(\xi_i)^* = R(Z_i)^*
$$
voir la section \ref{chow-section-tame-symbol} de l'Homologie de Chow.

\begin{lemma}[Formule clé]
\label{lemma-key-formula}
Dans la situation ci-dessus, le cycle
$$
\sum
(Z_i \to X)_*\left(
\text{ord}_{B_i}(f_i) \text{div}_{\mathcal{N}|_{Z_i}}(t_i|_{Z_i}) -
\text{ord}_{B_i}(g_i) \text{div}_{\mathcal{L}|_{Z_i}}(s_i|_{Z_i}) \right)
$$
est égal au cycle
$$
\sum (Z_i \to X)_*\text{div}(\partial_{B_i}(f_i, g_i))
$$
\end{lemma}

\begin{proof}
Voici la stratégie de la preuve : nous nous ramenons d'abord au cas où $\mathcal{L}$ et
$\mathcal{N}$ sont des modules inversibles triviaux, puis nous modifions nos choix
de banalisations locales et, enfin, nous utilisons la localisation \'etale pour
nous ramener au cas des schémas\footnote{Il est possible qu'une preuve plus courte
s'obtienne en appliquant immédiatement la localisation \'etale.}.

\medskip\noindent
Première étape. Soit $q : T \to X$ le morphisme construit dans le lemme
\ref{lemma-Gm-torsor}. Nous utiliserons toutes les propriétés énoncées dans ce
lemme sans plus ample mention. En particulier, il suffit de montrer que les cycles
deviennent égaux après passage aux images réciproques par $q$. Désignons par $s'$ et $t'$ les images réciproques de $s$ et
$t$, vues comme sections méromorphes de $q^*\mathcal{L}$ et $q^*\mathcal{N}$.
Notons $Z'_i = q^{-1}(Z_i)$, notons $\xi'_i \in |Z'_i|$ le point générique,
notons $B'_i = \mathcal{O}_{T, \xi'_i}^h$, notons $\mathcal{L}_{\xi'_i}$ et
$\mathcal{N}_{\xi'_i}$ les images réciproques de $\mathcal{L}$ et $\mathcal{N}$ sur
$\Spec(B'_i)$. Rappelons que nous avons des diagrammes commutatifs
$$
\xymatrix{
\Spec(B'_i) \ar[r]_-{c_{\xi'_i}} \ar[d] & T \ar[d]^q \\
\Spec(B_i) \ar[r]^-{c_{\xi_i}} & X
}
$$
voir la remarque d'Espaces décents
\ref{decent-spaces-remark-functoriality-henselian-local-ring}. Notons $s'_i$
et $t'_i$ les images réciproques de $s_i$ et $t_i$ ; ce sont des générateurs de
$\mathcal{L}_{\xi'_i}$ et $\mathcal{N}_{\xi'_i}$. Alors
$$
s' = f'_i s'_i
\quad\text{et}\quad
t' = g'_i t'_i
$$
où $f'_i$ et $g'_i$ sont les images de $f_i, g_i$ par l'application $Q(B_i) \to
Q(B'_i)$ induite par $B_i \to B'_i$. D'après le lemme d'Algèbre
\ref{algebra-lemma-pullback-module}, nous avons
$$
\text{ord}_{B_i}(f_i) = \text{ord}_{B'_i}(f'_i)
\quad\text{et}\quad
\text{ord}_{B_i}(g_i) = \text{ord}_{B'_i}(g'_i)
$$
Par le lemme \ref{lemma-prepare-flat-pullback-cap-c1} appliqué au morphisme $q : Z'_i \to
Z_i$, nous avons
$$
q^*\text{div}_{\mathcal{N}|_{Z_i}}(t_i|_{Z_i}) =
\text{div}_{q^*\mathcal{N}|_{Z'_i}}(t'_i|_{Z'_i})
\quad\text{et}\quad
q^*\text{div}_{\mathcal{L}|_{Z_i}}(s_i|_{Z_i}) =
\text{div}_{q^*\mathcal{L}|_{Z'_i}}(s'_i|_{Z'_i})
$$
Cela montre déjà que l'image réciproque du premier cycle de l'énoncé du lemme
est le cycle correspondant aux données $s', t', Z'_i, s'_i, t'_i$. Pour établir
la même assertion pour le second, remarquons que le lemme de l'Homologie de Chow
\ref{chow-lemma-tame-symbol-formally-smooth} donne
$$
\partial_{B_i}(f_i, g_i) \mapsto
\partial_{B'_i}(f'_i, g'_i)
\quad\text{via}\quad
\kappa(\xi_i) \to \kappa(\xi'_i)
$$
Le même lemme qu'auparavant montre donc que
$$
q^*\text{div}(\partial_{B_i}(f_i, g_i)) =
\text{div}(\partial_{B'_i}(f'_i, g'_i))
$$
Puisque $q^*\mathcal{L} \cong \mathcal{O}_T$, nous constatons qu'il suffit de
démontrer l'égalité dans le cas où $\mathcal{L}$ est trivial. En échangeant les
rôles de $\mathcal{L}$ et $\mathcal{N}$, nous voyons que nous pouvons de même
supposer que $\mathcal{N}$ est trivial. Ceci achève la preuve de la
première étape.

\medskip\noindent
Deuxième étape. Supposons $\mathcal{L} = \mathcal{O}_X$ et $\mathcal{N} =
\mathcal{O}_X$. Désignons par $1$ la section qui banalise $\mathcal{L}$. Alors
$s_i = u \cdot 1$ pour une certaine unité $u \in B_i$. Examinons ce qui se passe
si nous remplaçons $s_i$ par $1$. Alors, $f_i$ est remplacé par $u f_i$.
Ainsi la première partie de la première expression du lemme est inchangée et
au second terme nous ajoutons
$$
\text{ord}_{B_i}(g_i)\text{div}(u|_{Z_i})
$$
où $u|_{Z_i}$ est l'image de $u$ dans le corps résiduel, d'après le lemme d'Espaces sur les corps \ref{spaces-over-fields-lemma-divisor-meromorphic-well-defined},
et à la deuxième expression nous ajoutons
$$
\text{div}(\partial_{B_i}(u, g_i))
$$
par bilinéarité du symbole modéré. Ces termes sont égaux en vertu de la
propriété du symbole modéré donnée par l'équation de l'Homologie de Chow
(\ref{chow-item-normalization}).

\medskip\noindent
Soit $Y \subset X$ un sous-espace fermé intègre avec $\dim_\delta(Y) = n -
2$. Pour montrer que les coefficients de $Y$ dans les deux cycles du lemme sont
égaux, nous pouvons remplacer $s_i$ par $1$ comme au paragraphe
précédent. De même, nous pouvons remplacer $t_i$ par $1$. Puisqu'il n'y a qu'un nombre fini de $Z_i$ tels
que $Y \subset Z_i$, nous pouvons supposer $s_i = 1$ et $t_i = 1$ pour tous
ces $Z_i$.

\medskip\noindent
Troisième étape. Nous montrons ici que les coefficients de $Y$ dans les cycles
du lemme coïncident pour un sous-espace fermé intègre $Y$ avec
$\dim_\delta(Y) = n - 2$, dans la situation où, de plus, $\mathcal{L} = \mathcal{O}_X$ et
$\mathcal{N} = \mathcal{O}_X$, et où $s_i = 1$ et $t_i = 1$ pour tous les $Z_i$ tels
que $Y \subset Z_i$. Quitte à remplacer $X$ par un sous-espace ouvert plus
petit, nous pouvons supposer que $s_i$ et $t_i$ sont égaux à $1$ pour
tout $i$. Dans ce cas, le premier cycle est nul. Il reste à montrer
que le coefficient de $Y$ dans le deuxième cycle est également nul.

\medskip\noindent
Tout d'abord, puisque $\mathcal{L} = \mathcal{O}_X$ et $\mathcal{N} =
\mathcal{O}_X$, nous identifions $s, t$ à des
fonctions rationnelles $f, g$ sur $X$. Puisque $s_i$ et $t_i$ sont égaux à
$1$, nous voyons que $f_i$, resp.\ $g_i$, est l'image de $f$, resp.\ $g$, dans
$Q(B_i)$ pour tous les $i$. Soit $\zeta \in |Y|$ le point générique.
Choisissons un voisinage \'etale
$$
(U, u) \longrightarrow (X, \zeta)
$$
et notons $Y' = \overline{\{u\}} \subset U$. Puisqu'un morphisme \'etale
est plat, nous pouvons prendre les images réciproques de $f$ et $g$, qui sont des fonctions méromorphes
régulières sur $U$ que nous noterons encore $f$ et $g$. Pour tout
diviseur premier $Y \subset Z \subset X$, le schéma $Z \times_X U$ est une
réunion de diviseurs premiers de $U$. Réciproquement, étant donné un diviseur
premier $Y' \subset Z' \subset U$, il existe un diviseur premier $Y \subset Z
\subset X$ tel que $Z'$ soit une composante de $Z \times_X U$. Pour une
telle paire $(Z, Z')$, l'homomorphisme d'anneaux
$$
\mathcal{O}_{X, \xi}^h \to \mathcal{O}_{U, \xi'}^h
$$
est \'etale (il est même fini \'etale). On obtient donc
$$
\partial_{\mathcal{O}_{X, \xi}^h}(f, g) \mapsto
\partial_{\mathcal{O}_{U, \xi'}^h}(f, g)
\quad\text{via}\quad
\kappa(\xi) \to \kappa(\xi')
$$
en vertu du lemme de l'Homologie de Chow \ref{chow-lemma-tame-symbol-formally-smooth}. Ainsi
le lemme \ref{lemma-etale-pullback-principal-divisor} s'applique et montre
$$
(Z \times_X U \to Z)^*\text{div}_Z(\partial_{\mathcal{O}_{X, \xi}^h}(f, g))
=
\sum\nolimits_{Z' \subset Z \times_X U}
\text{div}_{Z'}(\partial_{\mathcal{O}_{U, \xi'}^h}(f, g))
$$
Puisque l'image réciproque plate commute à l'image directe par les immersions
fermées (lemme \ref{lemma-flat-pullback-proper-pushforward}), nous voyons
qu'il suffit de démontrer que le coefficient de $Y'$ dans
$$
\sum\nolimits_{Z' \subset U} 
(Z' \to U)_*\text{div}_{Z'}(\partial_{\mathcal{O}_{U, \xi'}^h}(f, g))
$$
est nul.

\medskip\noindent
Soit $A = \mathcal{O}_{U, u}$. Alors $f, g \in Q(A)^*$. Écrivons
$f = a/b$ et $g = c/d$, où $a, b, c, d \in A$ sont des non-diviseurs de zéro.
Le coefficient de $Y'$ dans l'expression ci-dessus est
$$
\sum\nolimits_{\mathfrak q \subset A\text{ height }1}
\text{ord}_{A/\mathfrak q}(\partial_{A_\mathfrak q}(f, g))
$$
Par bilinéarité de $\partial_A$, il suffit de démontrer que
$$
\sum\nolimits_{\mathfrak q \subset A\text{ height }1}
\text{ord}_{A/\mathfrak q}(\partial_{A_\mathfrak q}(a, c))
$$
est nul. Il en est de même pour les autres paires $(a, d)$, $(b, c)$ et $(b, d)$. Cela
résulte du lemme de l'Homologie de Chow \ref{chow-lemma-key-nonzerodivisors}.
\end{proof}











\section{Intersection avec un faisceau inversible et équivalence rationnelle}
\label{section-commutativity}

\noindent
Cette section est l'analogue de l'Homologie de Chow, section
\ref{chow-section-commutativity}. En appliquant le lemme clé, nous obtenons
les propriétés fondamentales de l'intersection avec des faisceaux inversibles.
En particulier, nous verrons que $c_1(\mathcal{L}) \cap -$ passe au quotient par
l'équivalence rationnelle et que les opérations associées à différents faisceaux inversibles
commutent.

\begin{lemma}
\label{lemma-commutativity-on-integral}
Dans la situation \ref{situation-setup}, soit $X/B$ bon. Supposons
$X$ intègre et $\dim_\delta(X) = n$. Soient $\mathcal{L}$ et $\mathcal{N}$ des faisceaux
inversibles sur $X$. Choisissons une section méromorphe non nulle $s$ de
$\mathcal{L}$ et une section méromorphe non nulle $t$ de $\mathcal{N}$. Posons
$\alpha = \text{div}_\mathcal{L}(s)$ et $\beta = \text{div}_\mathcal{N}(t)$.
Alors
$$
c_1(\mathcal{N}) \cap \alpha
=
c_1(\mathcal{L}) \cap \beta
$$
dans $\CH_{n - 2}(X)$.
\end{lemma}

\begin{proof}
Cela résulte immédiatement du lemme clé \ref{lemma-key-formula} et de la discussion qui le
précède.
\end{proof}

\begin{lemma}
\label{lemma-factors}
Dans la situation \ref{situation-setup}, soit $X/B$ bon. Soit
$\mathcal{L}$ un faisceau inversible sur $X$. L'opération $\alpha \mapsto c_1(\mathcal{L})
\cap \alpha$ passe au quotient par l'équivalence rationnelle pour donner une opération
$$
c_1(\mathcal{L}) \cap - : \CH_{k + 1}(X) \to \CH_k(X)
$$
\end{lemma}

\begin{proof}
Soit $\alpha \in Z_{k + 1}(X)$ et $\alpha \sim_{rat} 0$. Nous devons montrer
que $c_1(\mathcal{L}) \cap \alpha$, tel qu'il est défini dans la définition
\ref{definition-cap-c1}, est nul. D'après la définition
\ref{definition-rational-equivalence}, il existe une famille localement finie
$\{W_j\}$ de sous-espaces fermés intègres avec $\dim_\delta(W_j) = k + 2$ et
des fonctions rationnelles $f_j \in R(W_j)^*$ telles que
$$
\alpha = \sum (i_j)_*\text{div}_{W_j}(f_j)
$$
Notons que $p : \coprod W_j \to X$ est un morphisme propre, et donc $\alpha =
p_*\alpha'$ où $\alpha' \in Z_{k + 1}(\coprod W_j)$ est la somme des diviseurs
principaux $\text{div}_{W_j}(f_j)$. Par le lemme
\ref{lemma-pushforward-cap-c1}, nous avons $c_1(\mathcal{L}) \cap \alpha =
p_*(c_1(p^*\mathcal{L}) \cap \alpha')$. Il suffit donc de montrer que chacun des cycles
$c_1(\mathcal{L}|_{W_j}) \cap \text{div}_{W_j}(f_j)$ est nul. En d'autres
termes, nous pouvons supposer que $X$ est intègre et $\alpha =
\text{div}_X(f)$ pour un certain $f \in R(X)^*$.

\medskip\noindent
Supposons $X$ intègre et $\alpha = \text{div}_X(f)$ pour
un certain $f \in R(X)^*$. Nous pouvons considérer $f$ comme une section
méromorphe régulière du faisceau inversible $\mathcal{N} = \mathcal{O}_X$.
Choisissons une section méromorphe $s$ de $\mathcal{L}$ et posons $\beta =
\text{div}_\mathcal{L}(s)$. Par le lemme \ref{lemma-commutativity-on-integral}
nous concluons que
$$
c_1(\mathcal{L}) \cap \alpha = c_1(\mathcal{O}_X) \cap \beta.
$$
Cependant, le lemme \ref{lemma-c1-cap-additive} montre que le membre
de droite est nul dans $\CH_k(X)$, comme voulu.
\end{proof}

\noindent
Dans la situation \ref{situation-setup}, soit $X/B$ bon. Soit
$\mathcal{L}$ un faisceau inversible sur $X$. Nous noterons
$$
c_1(\mathcal{L})^s \cap - : \CH_{k + s}(X) \to \CH_k(X)
$$
l'opération $c_1(\mathcal{L}) \cap - $. Cette notation a un sens d'après le lemme
\ref{lemma-factors}. Nous noterons $c_1(\mathcal{L}^s \cap -$ l'itérée
$s$-ième de cette opération pour tout $s \geq 0$.

\begin{lemma}
\label{lemma-cap-commutative}
Dans la situation \ref{situation-setup}, soit $X/B$ bon. Soient
$\mathcal{L}$ et $\mathcal{N}$ des faisceaux inversibles sur $X$. Pour tout $\alpha \in \CH_{k
+ 2}(X)$, nous avons
$$
c_1(\mathcal{L}) \cap c_1(\mathcal{N}) \cap \alpha
=
c_1(\mathcal{N}) \cap c_1(\mathcal{L}) \cap \alpha
$$
comme éléments de $\CH_k(X)$.
\end{lemma}

\begin{proof}
Écrivons $\alpha = \sum m_j[Z_j]$ pour une famille localement finie de
sous-espaces fermés intègres $Z_j \subset X$ avec $\dim_\delta(Z_j) = k + 2$.
Considérons le morphisme propre $p : \coprod Z_j \to X$. Posons
$\alpha' = \sum m_j[Z_j]$ comme un $(k + 2)$-cycle sur $\coprod Z_j$. Après
plusieurs applications du lemme \ref{lemma-pushforward-cap-c1}, nous voyons que
$c_1(\mathcal{L}) \cap c_1(\mathcal{N}) \cap \alpha = p_*(c_1(p^*\mathcal{L})
\cap c_1(p^*\mathcal{N}) \cap \alpha')$ et $c_1(\mathcal{N}) \cap
c_1(\mathcal{L}) \cap \alpha = p_*(c_1(p^*\mathcal{N}) \cap
c_1(p^*\mathcal{L}) \cap \alpha')$. Il suffit donc de démontrer la formule dans
le cas où $X$ est intègre et $\alpha = [X]$. Dans ce cas, le résultat
découle du lemme \ref{lemma-commutativity-on-integral} et des définitions.
\end{proof}

















\section{Intersection avec les diviseurs de Cartier effectifs}
\label{section-intersecting-effective-Cartier}

\noindent
Cette section est l'analogue de l'Homologie de Chow, section
\ref{chow-section-intersecting-effective-Cartier}. Nous renvoyons à
l'introduction de cette section pour les motivations.

\medskip\noindent
Rappelons que les diviseurs de Cartier effectifs correspondent de façon $1$-à-$1$ aux
classes d'isomorphisme de couples $(\mathcal{L}, s)$ où $\mathcal{L}$ est un
faisceau inversible et $s$ une section globale ; voir le lemme de Diviseurs sur les
espaces \ref{spaces-divisors-lemma-characterize-OD}. Si $D$ correspond
à $(\mathcal{L}, s)$, alors $\mathcal{L} = \mathcal{O}_X(D)$. Gardons
ce fait à l'esprit au cours de cette section.

\begin{definition}
\label{definition-gysin-homomorphism}
Dans la situation \ref{situation-setup}, soit $X/B$ bon. Soit
$(\mathcal{L}, s)$ un couple constitué d'un faisceau inversible et d'une section
globale $s \in \Gamma(X, \mathcal{L})$. Soit $D = Z(s)$ le schéma des zéros
de $s$, et notons $i : D \to X$ l'immersion fermée. Nous définissons, pour
tout entier $k$, un {\it homomorphisme de Gysin raffiné}
$$
i^* : Z_{k + 1}(X) \to \CH_k(D).
$$
suivant les règles ci-dessous :
\begin{enumerate}
\item Étant donné un sous-espace fermé intègre $W \subset X$ avec
$\dim_\delta(W) = k + 1$, nous posons
\begin{enumerate}
\item si $W \not \subset D$, alors $i^*[W] = [D \cap W]_k$ comme
$k$-cycle sur $D$, et
\item si $W \subset D$, alors $i^*[W] = i'_*(c_1(\mathcal{L}|_W) \cap [W])$,
où $i' : W \to D$ est l'immersion fermée induite.
\end{enumerate}
\item Pour un $(k + 1)$-cycle général $\alpha = \sum n_j[W_j]$, nous
posons
$$
i^*\alpha = \sum n_j i^*[W_j]
$$
\item Si $D$ est un diviseur de Cartier effectif, nous notons $D \cdot
\alpha = i_*i^*\alpha$ l'image directe de la classe, considérée comme une classe sur $X$.
\end{enumerate}
\end{definition}

\noindent
En fait, comme nous le verrons plus tard, cet homomorphisme de Gysin $i^*$
peut être considéré comme un exemple d'image réciproque non plate. Ainsi, nous
appellerons parfois de manière informelle la classe $i^*\alpha$ l'{\it
image réciproque} de la classe $\alpha$.

\begin{remark}
\label{remark-on-cycles}
Soient $S$, $B$, $X$, $\mathcal{L}$, $s$, $i : D \to X$ les données de la définition
\ref{definition-gysin-homomorphism} et supposons que $\mathcal{L}|_D \cong
\mathcal{O}_D$. Dans ce cas, nous pouvons définir une application canonique $i^* :
Z_{k + 1}(X) \to Z_k(D)$ au niveau des cycles, en imposant $i^*[W] = 0$ lorsque
$W \subset D$. Cette possibilité sera utile plus loin.
\end{remark}

\begin{remark}
\label{remark-pullback-pairs}
Soit $f : X' \to X$ un morphisme de bons espaces algébriques au-dessus de $B$ comme
dans la situation \ref{situation-setup}. Soit $(\mathcal{L}, s, i : D \to X)$
un triplet comme dans la définition \ref{definition-gysin-homomorphism}.
Posons $\mathcal{L}' = f^*\mathcal{L}$, $s' = f^*s$ et
$D' = X' \times_X D = Z(s')$. Nous obtenons un diagramme commutatif
$$
\xymatrix{
D' \ar[d]_g \ar[r]_{i'} & X' \ar[d]^f \\
D \ar[r]^i & X
}
$$
et nous pouvons demander diverses compatibilités entre $i^*$ et $(i')^*$.
\end{remark}

\begin{lemma}
\label{lemma-support-cap-effective-Cartier}
Dans la situation \ref{situation-setup}, soit $X/B$ bon. Soit
$(\mathcal{L}, s, i : D \to X)$ comme dans la définition
\ref{definition-gysin-homomorphism}. Soit $\alpha$ un $(k + 1)$-cycle sur $X$.
Alors $i_*i^*\alpha = c_1(\mathcal{L}) \cap \alpha$ dans $\CH_k(X)$. En
particulier, si $D$ est un diviseur de Cartier effectif, alors $D \cdot \alpha =
c_1(\mathcal{O}_X(D)) \cap \alpha$.
\end{lemma}

\begin{proof}
Écrivons $\alpha = \sum n_j[W_j]$, où $i_j : W_j \to X$ sont les immersions de sous-espaces
fermés intègres avec $\dim_\delta(W_j) = k$. Puisque $D$ est le schéma des zéros
de $s$, nous voyons que $D \cap W_j$ est le schéma des zéros de
la restriction $s|_{W_j}$. Par conséquent, pour tout $j$ tel que $W_j \not
\subset D$, nous avons $c_1(\mathcal{L}) \cap [W_j] = [D \cap W_j]_k$ par le
lemme \ref{lemma-geometric-cap}. Nous avons donc
$$
c_1(\mathcal{L}) \cap \alpha
=
\sum\nolimits_{W_j \not \subset D} n_j[D \cap W_j]_k
+
\sum\nolimits_{W_j \subset D}
n_j i_{j, *}(c_1(\mathcal{L})|_{W_j}) \cap [W_j])
$$
dans $\CH_k(X)$ d'après la définition \ref{definition-cap-c1}. Le membre de droite
coïncide, terme à terme, avec l'image directe de la classe $i^*\alpha$ sur $D$
définie dans la définition \ref{definition-gysin-homomorphism}. Le résultat en découle.
\end{proof}

\begin{lemma}
\label{lemma-closed-in-X-gysin}
Dans la situation \ref{situation-setup}, soit $f : X' \to X$ un morphisme
propre entre de bons espaces algébriques au-dessus de $B$. Soit $(\mathcal{L}, s, i : D \to
X)$ comme dans la définition \ref{definition-gysin-homomorphism}. Formons le
diagramme
$$
\xymatrix{
D' \ar[d]_g \ar[r]_{i'} & X' \ar[d]^f \\
D \ar[r]^i & X
}
$$
comme dans la remarque \ref{remark-pullback-pairs}. Pour tout $(k + 1)$-cycle
$\alpha'$ sur $X'$, nous avons $i^*f_*\alpha' = g_*(i')^*\alpha'$ dans
$\CH_k(D)$ (ce qui a un sens puisque $f_*$ est défini au niveau des cycles).
\end{lemma}

\begin{proof}
Supposons $\alpha = [W']$ pour un sous-espace fermé intègre $W' \subset X'$.
Soit $W \subset X$ l'« image » de $W'$ comme dans le lemme
\ref{lemma-proper-image}. Si $W' \not \subset D'$, alors $W \not
\subset D$ et nous voyons que
$$
[W' \cap D']_k = \text{div}_{\mathcal{L}'|_{W'}}({s'|_{W'}})
\quad\text{et}\quad
[W \cap D]_k = \text{div}_{\mathcal{L}|_W}(s|_W)
$$
L'image par $f_*$ du premier cycle est donc égale au second, d'après le lemme
\ref{lemma-equal-c1-as-cycles}.
L'égalité vaut ainsi au niveau des cycles.
Si $W' \subset D'$, alors $W \subset D$ et
$f_*(c_1(\mathcal{L}|_{W'}) \cap [W'])$ est égal à $c_1(\mathcal{L}|_W) \cap
[W]$ dans $\CH_k(W)$ par la deuxième assertion du lemme
\ref{lemma-equal-c1-as-cycles}. Par la remarque
\ref{remark-infinite-sums-rational-equivalences}, le résultat en découle pour
$\alpha'$ quelconque.
\end{proof}

\begin{lemma}
\label{lemma-gysin-flat-pullback}
Dans la situation \ref{situation-setup}, soit $f : X' \to X$ un morphisme plat
de dimension relative $r$ entre de bons espaces algébriques au-dessus de $B$. Soit
$(\mathcal{L}, s, i : D \to X)$ comme dans la définition
\ref{definition-gysin-homomorphism}. Considérons le diagramme
$$
\xymatrix{
D' \ar[d]_g \ar[r]_{i'} & X' \ar[d]^f \\
D \ar[r]^i & X
}
$$
comme dans la remarque \ref{remark-pullback-pairs}. Pour tout $(k + 1)$-cycle
$\alpha$ sur $X$, nous avons $(i')^*f^*\alpha = g^*i^*\alpha'$ dans $\CH_{k +
r}(D)$ (ce qui a un sens puisque $f^*$ est défini au niveau des cycles).
\end{lemma}

\begin{proof}
Supposons $\alpha = [W]$ pour un sous-espace fermé intègre $W \subset X$.
Soit $W' = f^{-1}(W) \subset X'$. Si $W \not \subset D$, alors $W'
\not \subset D'$ et nous voyons que
$$
W' \cap D' = g^{-1}(W \cap D)
$$
comme sous-espaces fermés de $D'$. L'égalité vaut donc au niveau des cycles ;
voir le lemme \ref{lemma-pullback-coherent}. Si $W
\subset D$, alors $W' \subset D'$ et $W' = g^{-1}(W)$, avec $[W']_{k + 1 + r} =
g^*[W]$, et l'égalité vaut dans $\CH_{k + r}(D')$ par le lemme
\ref{lemma-flat-pullback-cap-c1}. Par la remarque
\ref{remark-infinite-sums-rational-equivalences}, le résultat en découle pour
$\alpha'$ quelconque.
\end{proof}

\begin{lemma}
\label{lemma-easy-gysin}
Dans la situation \ref{situation-setup}, soit $X/B$ bon. Soit
$(\mathcal{L}, s, i : D \to X)$ comme dans la définition
\ref{definition-gysin-homomorphism}. Soit $Z \subset X$ un sous-schéma fermé
tel que $\dim_\delta(Z) \leq k + 1$ et tel que $D \cap Z$ soit un diviseur de
Cartier effectif sur $Z$. Alors $i^*([Z]_{k + 1}) = [D \cap Z]_k$.
\end{lemma}

\begin{proof}
L'hypothèse signifie que $s|_Z$ est une section régulière de $\mathcal{L}|_Z$.
Ainsi, $D \cap Z = Z(s)$ et nous obtenons
$$
[D \cap Z]_k = \sum n_i [Z(s_i)]_k
$$
comme égalité de cycles, où $s_i = s|_{Z_i}$, les $Z_i$ sont les composantes irréductibles
de $\delta$-dimension $k + 1$ et $[Z]_{k + 1} = \sum n_i[Z_i]$ ; voir le lemme
\ref{lemma-prepare-geometric-cap}. Nous avons $D \cap Z_i = Z(s_i)$. La
comparaison avec la définition de l'application de Gysin achève la preuve.
\end{proof}













\section{Homomorphismes de Gysin}
\label{section-gysin}

\noindent
Cette section est l'analogue de l'Homologie de Chow, section
\ref{chow-section-gysin}. Dans cette section, nous utilisons la formule clé
pour montrer que l'homomorphisme de Gysin passe au quotient par l'équivalence rationnelle.

\begin{lemma}
\label{lemma-gysin-factors-general}
Dans la situation \ref{situation-setup}, soit $X/B$ bon. Supposons
$X$ intègre et $n = \dim_\delta(X)$. Soit $i : D \to X$ un
diviseur de Cartier effectif. Soit $\mathcal{N}$ un $\mathcal{O}_X$-module
inversible et soit $t$ une section méromorphe non nulle de $\mathcal{N}$. Alors
$i^*\text{div}_\mathcal{N}(t) = c_1(\mathcal{N}) \cap [D]_{n - 1}$ dans
$\CH_{n - 2}(D)$.
\end{lemma}

\begin{proof}
Écrivons $\text{div}_\mathcal{N}(t) = \sum \text{ord}_{Z_i,
\mathcal{N}}(t)[Z_i]$ pour des sous-espaces fermés intègres $Z_i \subset
X$ de $\delta$-dimension $n - 1$. Nous pouvons supposer que la famille $\{Z_i\}$
est localement finie, que $t \in \Gamma(U, \mathcal{N}|_U)$ est un générateur
où $U = X \setminus \bigcup Z_i$, et que chaque composante irréductible de $D$
est l'un des $Z_i$ ; voir les lemmes d'Espaces sur les corps
\ref{spaces-over-fields-lemma-components-locally-finite},
\ref{spaces-over-fields-lemma-divisor-locally-finite} et
\ref{spaces-over-fields-lemma-divisor-meromorphic-locally-finite}.

\medskip\noindent
Posons $\mathcal{L} = \mathcal{O}_X(D)$. Notons $s \in \Gamma(X,
\mathcal{O}_X(D)) = \Gamma(X, \mathcal{L})$ la section canonique. Nous
appliquerons la discussion de la section \ref{section-key} à notre situation
actuelle. Pour tout $i$, soit $\xi_i \in |Z_i|$ son point générique. Soit
$B_i = \mathcal{O}_{X, \xi_i}^h$. Pour tout $i$, choisissons des
générateurs $s_i$ de $\mathcal{L}_{\xi_i}$ et $t_i$ de $\mathcal{N}_{\xi_i}$
sur $B_i$, en imposant $s_i = s$ si
$Z_i \not \subset D$. Écrivons $s = f_i s_i$ et $t = g_i t_i$ avec $f_i, g_i
\in B_i$. Alors $\text{ord}_{Z_i, \mathcal{N}}(t) = \text{ord}_{B_i}(g_i)$.
D'autre part, nous avons $f_i \in B_i$ et
$$
[D]_{n - 1} = \sum \text{ord}_{B_i}(f_i)[Z_i]
$$
grâce à notre choix des $s_i$. Nous affirmons que
$$
i^*\text{div}_\mathcal{N}(t) =
\sum \text{ord}_{B_i}(g_i) \text{div}_{\mathcal{L}|_{Z_i}}(s_i|_{Z_i})
$$
comme égalité de cycles. Plus précisément, le membre de droite est un cycle
représentant le membre de gauche. En effet, cela résulte de notre formule
pour $\text{div}_\mathcal{N}(t)$ et le fait que
$\text{div}_{\mathcal{L}|_{Z_i}}(s_i|_{Z_i}) = [Z(s_i|_{Z_i})]_{n - 2} = [Z_i
\cap D]_{n - 2}$ lorsque $Z_i \not \subset D$ car dans ce cas $s_i|_{Z_i} =
s|_{Z_i}$ est une section régulière ; voir le lemme \ref{lemma-compute-c1}. De
même,
$$
c_1(\mathcal{N}) \cap [D]_{n - 1} =
\sum \text{ord}_{B_i}(f_i) \text{div}_{\mathcal{N}|_{Z_i}}(t_i|_{Z_i})
$$
La formule clé (lemme \ref{lemma-key-formula}) donne l'égalité
$$
\sum \left(
\text{ord}_{B_i}(f_i) \text{div}_{\mathcal{N}|_{Z_i}}(t_i|_{Z_i}) -
\text{ord}_{B_i}(g_i) \text{div}_{\mathcal{L}|_{Z_i}}(s_i|_{Z_i}) \right) =
\sum \text{div}_{Z_i}(\partial_{B_i}(f_i, g_i))
$$
de cycles. Si $Z_i \not \subset D$, alors $f_i = 1$ et donc
$\text{div}_{Z_i}(\partial_{B_i}(f_i, g_i)) = 0$. Nous obtenons ainsi une
équivalence rationnelle entre nos cycles spécifiques représentant
$i^*\text{div}_\mathcal{N}(t)$ et $c_1(\mathcal{N}) \cap [D]_{n - 1}$ sur $D$.
Ceci achève la preuve.
\end{proof}

\begin{lemma}
\label{lemma-gysin-factors}
Dans la situation \ref{situation-setup}, soit $X/B$ bon. Soit
$(\mathcal{L}, s, i : D \to X)$ comme dans la définition
\ref{definition-gysin-homomorphism}. L'homomorphisme de Gysin passe au quotient par
l'équivalence rationnelle et donne une application $i^* : \CH_{k + 1}(X) \to
\CH_k(D)$.
\end{lemma}

\begin{proof}
Soit $\alpha \in Z_{k + 1}(X)$ et supposons que $\alpha \sim_{rat} 0$. Cela
signifie qu'il existe une famille localement finie de sous-espaces fermés
intègres $W_j \subset X$ de $\delta$-dimension $k + 2$ et des fonctions
$f_j \in R(W_j)^*$ telles que $\alpha = \sum i_{j, *}\text{div}_{W_j}(f_j)$. Posons $X' =
\coprod W_i$ et considérons le diagramme
$$
\xymatrix{
D' \ar[d]_q \ar[r]_{i'} & X' \ar[d]^p \\
D \ar[r]^i & X
}
$$
de la remarque \ref{remark-pullback-pairs}. Puisque $X' \to X$ est propre,
nous voyons que $i^*p_* = q_*(i')^*$ par le lemme
\ref{lemma-closed-in-X-gysin}. Comme nous savons que $q_*$ passe au quotient par
l'équivalence rationnelle (lemme
\ref{lemma-proper-pushforward-rational-equivalence}), il suffit de démontrer le
résultat pour $\alpha' = \sum \text{div}_{W_j}(f_j)$ sur $X'$. Cela
nous ramène au cas où $X$ est intègre et $\alpha = \text{div}(f)$ pour
un certain $f \in R(X)^*$.

\medskip\noindent
Supposons $X$ intègre et $\alpha = \text{div}(f)$ pour
un certain $f \in R(X)^*$. Si $X = D$, alors nous voyons que $i^*\alpha$ est
égal à $c_1(\mathcal{L}) \cap \alpha$. Ceci est rationnellement équivalent à
zéro par le lemme \ref{lemma-factors}. Si $D \not = X$, alors nous voyons que
$i^*\text{div}_X(f)$ est égal à $c_1(\mathcal{O}_D) \cap [D]_{n - 1}$ dans
$\CH_k(D)$ par le lemme \ref{lemma-gysin-factors-general}. Bien entendu,
l'intersection avec $c_1(\mathcal{O}_D)$ est l'application nulle.
\end{proof}

\begin{lemma}
\label{lemma-gysin-commutes-cap-c1}
Dans la situation \ref{situation-setup}, soit $X/B$ bon. Soit
$(\mathcal{L}, s, i : D \to X)$ un triplet comme dans la définition
\ref{definition-gysin-homomorphism}. Soit $\mathcal{N}$ un
$\mathcal{O}_X$-module inversible. Alors $i^*(c_1(\mathcal{N}) \cap \alpha) =
c_1(i^*\mathcal{N}) \cap i^*\alpha$ dans $\CH_{k - 2}(D)$ pour tout
$\alpha \in \CH_k(Z)$.
\end{lemma}

\begin{proof}
Avec exactement la même preuve que dans le lemme \ref{lemma-gysin-factors},
cela découle des lemmes \ref{lemma-pushforward-cap-c1},
\ref{lemma-cap-commutative} et \ref{lemma-gysin-factors-general}.
\end{proof}

\begin{lemma}
\label{lemma-gysin-commutes-gysin}
Dans la situation \ref{situation-setup}, soit $X/B$ bon. Soient
$(\mathcal{L}, s, i : D \to X)$ et $(\mathcal{L}', s', i' : D' \to X)$ deux
triplets comme dans la définition \ref{definition-gysin-homomorphism}. Alors le
diagramme
$$
\xymatrix{
\CH_k(X) \ar[r]_{i^*} \ar[d]_{(i')^*} & \CH_{k - 1}(D) \ar[d] \\
\CH_{k - 1}(D') \ar[r] & \CH_{k - 2}(D \cap D')
}
$$
est commutatif, chaque flèche étant un homomorphisme de Gysin.
\end{lemma}

\begin{proof}
Désignons par $j : D \cap D' \to D$ et $j' : D \cap D' \to D'$ les immersions fermées
correspondant à $(\mathcal{L}|_{D'}, s|_{D'}$ et $(\mathcal{L}'_D, s|_D)$.
Nous devons montrer que $(j')^*i^*\alpha = j^* (i')^*\alpha$ pour tout
$\alpha \in \CH_k(X)$. Soit $W \subset X$ un sous-schéma fermé intègre de
dimension $k$. Nous établirons l'égalité dans le cas $\alpha = [W]$. Le cas
général découlera alors de l'observation de la remarque
\ref{remark-infinite-sums-rational-equivalences} (et de la forme particulière de
notre équivalence rationnelle produite ci-dessous). Nous déduirons l'égalité
pour $\alpha = [W]$ de la formule clé.

\medskip\noindent
Prenons pour $\sigma$ une section méromorphe non nulle de
$\mathcal{L}|_W$, égale à $s|_W$ si $W \not
\subset D$. Prenons pour $\sigma'$ une section méromorphe non nulle de
$\mathcal{L}'|_W$, égale à $s'|_W$ si $W \not
\subset D'$. Écrivons
$$
\text{div}_{\mathcal{L}|_W}(\sigma) =
\sum \text{ord}_{Z_i, \mathcal{L}|_W}(\sigma)[Z_i] = \sum n_i[Z_i]
$$
et de même
$$
\text{div}_{\mathcal{L}'|_W}(\sigma') =
\sum \text{ord}_{Z_i, \mathcal{L}'|_W}(\sigma')[Z_i] = \sum n'_i[Z_i]
$$
comme dans la discussion de la section \ref{section-key}. Nous voyons alors
que $Z_i \subset D$ si $n_i \not = 0$ et $Z'_i \subset D'$ si $n'_i \not = 0$.
Pour tout $i$, soit $\xi_i \in |Z_i|$ le point générique. Comme dans la
section \ref{section-key}, choisissons pour tout $i$ un élément $\sigma_i \in
\mathcal{L}_{\xi_i}$, resp.\ $\sigma'_i \in \mathcal{L}'_{\xi_i}$, qui soit un générateur
sur $B_i = \mathcal{O}_{W, \xi_i}^h$ et qui soit égal à l'image de $s$, resp.\
$s'$, si $Z_i \not \subset D$, resp.\ $Z_i \not \subset D'$. Écrivons $\sigma =
f_i \sigma_i$ et $\sigma' = f'_i\sigma'_i$ de sorte que $n_i =
\text{ord}_{B_i}(f_i)$ et $n'_i = \text{ord}_{B_i}(f'_i)$. Il résulte de nos définitions que
$$
(j')^*i^*[W] =
\sum \text{ord}_{B_i}(f_i) \text{div}_{\mathcal{L}'|_{Z_i}}(\sigma'_i|_{Z_i})
$$
comme égalité de cycles, et
$$
j^*(i')^*[W] =
\sum \text{ord}_{B_i}(f'_i) \text{div}_{\mathcal{L}|_{Z_i}}(\sigma_i|_{Z_i})
$$
La formule clé (lemme \ref{lemma-key-formula}) donne alors l'égalité
$$
\sum \left(
\text{ord}_{B_i}(f_i) \text{div}_{\mathcal{L}'|_{Z_i}}(\sigma'_i|_{Z_i}) -
\text{ord}_{B_i}(f'_i) \text{div}_{\mathcal{L}|_{Z_i}}(\sigma_i|_{Z_i})
\right) =
\sum \text{div}_{Z_i}(\partial_{B_i}(f_i, f'_i))
$$
de cycles. Notons que $\text{div}_{Z_i}(\partial_{B_i}(f_i, f'_i)) = 0$ si
$Z_i \not \subset D \cap D'$, car dans ce cas soit $f_i = 1$, soit $f'_i = 1$. Nous
obtenons ainsi une équivalence rationnelle entre nos cycles spécifiques
représentant $(j')^*i^*[W]$ et $j^*(i')^*[W]$ sur $D \cap D' \cap W$.
\end{proof}










\section{Diviseurs de Cartier effectifs relatifs}
\label{section-relative-effective-cartier}

\noindent
Cette section est l'analogue de l'Homologie de Chow, section
\ref{chow-section-relative-effective-cartier}. Les diviseurs de Cartier effectifs
relatifs sont définis dans la section de Diviseurs sur les espaces
\ref{spaces-divisors-section-effective-Cartier-morphisms}. Pour développer les
résultats de base sur les classes de Chern de fibrés vectoriels, nous n'avons
besoin que du cas où le schéma ambiant et le diviseur de Cartier effectif sont
plats sur la base.

\begin{lemma}
\label{lemma-relative-effective-cartier}
Dans la situation \ref{situation-setup}, soient $X, Y/B$ bons. Soit $p : X
\to Y$ un morphisme plat de dimension relative $r$. Soit $i : D \to X$ un
diviseur de Cartier effectif relatif (Diviseurs sur les espaces, définition
\ref{spaces-divisors-definition-relative-effective-Cartier-divisor}). Soit
$\mathcal{L} = \mathcal{O}_X(D)$. Pour tout $\alpha \in \CH_{k + 1}(Y)$, nous
avons
$$
i^*p^*\alpha = (p|_D)^*\alpha
$$
dans $\CH_{k + r}(D)$ et
$$
c_1(\mathcal{L}) \cap p^*\alpha = i_* ((p|_D)^*\alpha)
$$
dans $\CH_{k + r}(X)$.
\end{lemma}

\begin{proof}
Soit $W \subset Y$ un sous-espace fermé intègre de $\delta$-dimension $k +
1$. Le lemme de Diviseurs sur les espaces
\ref{spaces-divisors-lemma-relative-Cartier} montre que $D \cap p^{-1}W$
est un diviseur de Cartier effectif sur $p^{-1}W$. Par le lemme
\ref{lemma-easy-gysin}, nous obtenons la première égalité dans
$$
i^*[p^{-1}W]_{k + r + 1} =
[D \cap p^{-1}W]_{k + r} =
[(p|_D)^{-1}(W)]_{k + r}.
$$
et la seconde vient de ce que $D \cap p^{-1}(W) = (p|_D)^{-1}(W)$ comme espaces
algébriques. Puisque, par définition, $p^*[W] = [p^{-1}W]_{k + r + 1}$, nous
voyons que $i^*p^*[W] = (p|_D)^*[W]$ comme égalité de cycles. Si $\alpha = \sum
m_j[W_j]$ est un $k + 1$-cycle général, alors nous obtenons $i^*\alpha = \sum
m_j i^*p^*[W_j] = \sum m_j(p|_D)^*[W_j]$ comme égalité de cycles. Cela démontre
la première égalité. Pour déduire la seconde de la première, appliquons le lemme
\ref{lemma-support-cap-effective-Cartier}.
\end{proof}












\section{Fibrés affines}
\label{section-affine-vector}

\noindent
Cette section est l'analogue de l'Homologie de Chow, section
\ref{chow-section-affine-vector}. Pour un fibré affine, l'application d'image réciproque
est surjective sur les groupes de Chow.

\begin{lemma}
\label{lemma-pullback-affine-fibres-surjective}
Dans la situation \ref{situation-setup}, soient $X, Y/B$ bons. Soit $f :
X \to Y$ un morphisme plat quasi-compact au-dessus de $B$, de dimension relative $r$.
Supposons que, pour tout $y \in Y$, nous ayons $X_y \cong
\mathbf{A}^r_{\kappa(y)}$. Alors $f^* : \CH_k(Y) \to \CH_{k + r}(X)$ est
surjectif pour tout $k \in \mathbf{Z}$.
\end{lemma}

\begin{proof}
Soit $\alpha \in \CH_{k + r}(X)$. Écrivons $\alpha = \sum m_j[W_j]$ avec
$m_j \not = 0$ et $W_j$ des sous-espaces fermés intègres distincts deux à deux
de $\delta$-dimension $k + r$. Alors la famille $\{W_j\}$ est localement
finie sur $X$. Soit $Z_j \subset Y$ le sous-espace fermé intègre tel que
le morphisme obtenu $W_j \to Z_j$ soit dominant, comme dans le lemme
\ref{lemma-proper-image}. Pour tout $V \subset Y$ ouvert quasi-compact, nous
voyons que $f^{-1}(V) \cap W_j$ n'est non vide que pour un nombre fini de $j$.
Par conséquent, les $Z_j$, fermetures des images, forment une famille
localement finie de sous-espaces fermés intègres de $Y$.

\medskip\noindent
Considérons les diagrammes cartésiens
$$
\xymatrix{
f^{-1}(Z_j) \ar[r] \ar[d]_{f_j} & X \ar[d]^f \\
Z_j \ar[r] & Y
}
$$
Supposons que $[W_j] \in Z_{k + r}(f^{-1}(Z_j))$ soit rationnellement
équivalent à $f_j^*\beta_j$ pour un $k$-cycle $\beta_j \in \CH_k(Z_j)$.
Alors $\beta = \sum m_j \beta_j$ sera un $k$-cycle sur $Y$ et $f^*\beta = \sum
m_j f_j^*\beta_j$ sera rationnellement équivalent à $\alpha$ (voir la remarque
\ref{remark-infinite-sums-rational-equivalences}). Cela nous ramène au cas où $Y$ est
intègre et où $\alpha = [W]$ pour un sous-schéma fermé intègre de $X$ dominant
$Y$. En particulier, nous pouvons supposer que $d = \dim_\delta(Y) < \infty$.

\medskip\noindent
Nous pouvons donc raisonner par récurrence sur $d = \dim_\delta(Y)$. Si $d < k$,
alors $\CH_{k + r}(X) = 0$ et le lemme est valable ; c'est le cas initial de
la récurrence. Soit $V \subset Y$ un ouvert non vide. Supposons que nous
puissions montrer que $\alpha|_{f^{-1}(V)} = f^*\beta$ pour un certain $\beta
\in Z_k(V)$. Le lemme \ref{lemma-exact-sequence-open} montre que
$\beta = \beta'|_V$ pour un certain $\beta' \in Z_k(Y)$. La suite exacte
$\CH_k(f^{-1}(Y \setminus V)) \to \CH_k(X) \to \CH_k(f^{-1}(V))$ du lemme
\ref{lemma-restrict-to-open} montre que $\alpha - f^*\beta'$ provient d'un
cycle $\alpha' \in \CH_{k + r}(f^{-1}(Y \setminus V))$. Puisque $\dim_\delta(Y
\setminus V) < d$, nous concluons par récurrence sur $d$.

\medskip\noindent
En particulier, quitte à remplacer $Y$ par un ouvert convenable, nous pouvons
supposer que $Y$ est un schéma ayant pour point générique $\eta$. L'isomorphisme
$Y_\eta \cong \mathbf{A}^r_\eta$ s'étend à un isomorphisme sur un $V \subset
Y$ ouvert non vide ; voir le lemme de Limites d'espaces
\ref{spaces-limits-lemma-descend-finite-presentation}. Cela nous ramène au cas
des schémas, traité dans le lemme de l'Homologie de Chow
\ref{chow-lemma-pullback-affine-fibres-surjective}.
\end{proof}

\begin{lemma}
\label{lemma-linebundle}
Dans la situation \ref{situation-setup}, soit $X/B$ bon. Soit
$\mathcal{L}$ un $\mathcal{O}_X$-module inversible. Soit
$$
p :
L = \underline{\Spec}(\text{Sym}^*(\mathcal{L}))
\longrightarrow
X
$$
le fibré vectoriel associé au-dessus de $X$. Alors $p^* : \CH_k(X) \to \CH_{k +
1}(L)$ est un isomorphisme pour tout $k$.
\end{lemma}

\begin{proof}
Pour la surjectivité, voir le lemme
\ref{lemma-pullback-affine-fibres-surjective}. Soit $o : X \to L$ la section
nulle de $L \to X$, c'est-à-dire le morphisme correspondant à la surjection
$\text{Sym}^*(\mathcal{L}) \to \mathcal{O}_X$ qui envoie $\mathcal{L}^{\otimes
n}$ sur zéro pour tout $n > 0$. Alors $p \circ o = \text{id}_X$ et $o(X)$
est un diviseur de Cartier effectif sur $L$. Ainsi, d'après le lemme
\ref{lemma-relative-effective-cartier}, nous avons $o^* \circ p^* =
\text{id}$ et nous concluons que $p^*$ est également injectif.
\end{proof}























\section{Théorie bivariante de l'intersection}
\label{section-bivariant}

\noindent
Cette section est l'analogue de l'Homologie de Chow, section
\ref{chow-section-bivariant}. Afin de traiter convenablement les classes de Chern
supérieures des fibrés vectoriels, nous introduisons la notion suivante, à la
suite de \cite{FM}. Il résulte de \cite[théorème 17.1]{F} que notre définition
coïncide avec celle de \cite{F}, sous réserve de la différence entre les
contextes considérés.

\begin{definition}
\label{definition-bivariant-class}
\begin{reference}
Comparer avec \cite[Definition 17.1]{F}
\end{reference}
Dans la situation \ref{situation-setup}, soit $f : X \to Y$ un morphisme de
bons espaces algébriques au-dessus de $B$. Soit $p \in \mathbf{Z}$. Une
{\it classe bivariante $c$ de degré $p$ pour $f$} est la donnée d'une règle qui
associe à tout morphisme $Y' \to Y$ de bons espaces algébriques au-dessus de $B$ et à
tout $k$ une application
$$
c \cap - : \CH_k(Y') \longrightarrow \CH_{k - p}(X')
$$
où $X' = Y' \times_Y X$, et qui satisfait aux conditions suivantes :
\begin{enumerate}
\item si $Y'' \to Y'$ est un morphisme propre, alors $c \cap (Y'' \to
Y')_*\alpha'' = (X'' \to X')_*(c \cap \alpha'')$ pour tout $\alpha''$ sur
$Y''$,
\item si $Y'' \to Y'$ est un morphisme de bons espaces algébriques au-dessus de $B$ qui est
plat de dimension relative $r$, alors $c \cap (Y'' \to Y')^*\alpha' = (X'' \to
X')^*(c \cap \alpha')$ pour tout $\alpha'$ sur $Y'$,
\item si $(\mathcal{L}', s', i' : D' \to Y')$ est comme dans la définition
\ref{definition-gysin-homomorphism}, dont l'image réciproque est $(\mathcal{N}', t', j' : E'
\to X')$ sur $X'$, alors nous avons $c \cap (i')^*\alpha' = (j')^*(c \cap
\alpha')$ pour tout $\alpha'$ sur $Y'$.
\end{enumerate}
L'ensemble de toutes les classes bivariantes de degré $p$ pour $f$ est noté
$A^p(X \to Y)$.
\end{definition}

\noindent
Dans la situation \ref{situation-setup}, soient $X \to Y$ et $Y \to Z$ des
morphismes de bons espaces algébriques au-dessus de $B$. Soit $p \in \mathbf{Z}$. Il
est clair que $A^p(X \to Y)$ est un groupe abélien. De plus, il est clair que
nous avons une composition bilinéaire
$$
A^p(X \to Y) \times A^q(Y \to Z) \to A^{p + q}(X \to Z)
$$
qui est associative. Nous nous intéresserons surtout à $A^p(X) = A^p(X \to X)$,
notation qui désignera toujours les classes de cohomologie bivariantes pour
$\text{id}_X$. C'est en effet là que vivront les classes de Chern.

\begin{definition}
\label{definition-chow-cohomology}
Dans la situation \ref{situation-setup}, soit $X/B$ bon. La
{\it cohomologie de Chow} de $X$ est la $\mathbf{Z}$-algèbre graduée $A^*(X)$ dont
la composante de degré $p$ est $A^p(X \to X)$.
\end{definition}

\noindent
Attention : il n'est pas clair que la structure de $\mathbf{Z}$-algèbre sur
$A^*(X)$ soit commutative, mais nous verrons que les classes de Chern vivent
en son centre.

\begin{remark}
\label{remark-pullback-cohomology}
Dans la situation \ref{situation-setup}, soit $f : X \to Y$ un morphisme de
bons espaces algébriques au-dessus de $B$. Il existe alors un homomorphisme canonique de
$\mathbf{Z}$-algèbres $A^*(Y) \to A^*(X)$. Plus précisément, étant donné $c \in
A^p(Y)$ et $X' \to X$, nous pouvons définir $f^*c$ par
l'application $c \cap - : \CH_k(X') \to \CH_{k - p}(X')$ qui est donnée en
considérant $X'$ comme un espace algébrique au-dessus de $Y$.
\end{remark}

\begin{lemma}
\label{lemma-cap-c1-bivariant}
Dans la situation \ref{situation-setup}, soit $X/B$ bon. Soit
$\mathcal{L}$ un $\mathcal{O}_X$-module inversible. Alors la règle qui
associe à $f : X' \to X$ $c_1(f^*\mathcal{L}) \cap - : \CH_k(X') \to \CH_{k -
1}(X')$ est une classe bivariante de degré $1$.
\end{lemma}

\begin{proof}
Cela découle des lemmes \ref{lemma-factors}, \ref{lemma-pushforward-cap-c1},
\ref{lemma-flat-pullback-cap-c1} et \ref{lemma-gysin-commutes-cap-c1}.
\end{proof}

\begin{lemma}
\label{lemma-flat-pullback-bivariant}
Dans la situation \ref{situation-setup}, soit $f : X \to Y$ un morphisme de
bons espaces algébriques au-dessus de $B$ qui est plat de dimension relative $r$.
Alors la règle qui associe à $Y' \to Y$ $(f')^* : \CH_k(Y') \to \CH_{k +
r}(X')$ où $X' = X \times_Y Y'$ est une classe bivariante de degré $-r$.
\end{lemma}

\begin{proof}
Cela découle des lemmes \ref{lemma-flat-pullback-rational-equivalence},
\ref{lemma-compose-flat-pullback},
\ref{lemma-flat-pullback-proper-pushforward} et
\ref{lemma-gysin-flat-pullback}.
\end{proof}

\begin{lemma}
\label{lemma-gysin-bivariant}
Dans la situation \ref{situation-setup}, soit $X/B$ bon. Soit
$(\mathcal{L}, s, i : D \to X)$ un triplet comme dans la définition
\ref{definition-gysin-homomorphism}. Alors la règle qui associe à $f : X'
\to X$ $(i')^* : \CH_k(X') \to \CH_{k - 1}(D')$ où $D' = D \times_X X'$ est
une classe bivariante de degré $1$.
\end{lemma}

\begin{proof}
Cela découle des lemmes \ref{lemma-gysin-factors},
\ref{lemma-closed-in-X-gysin}, \ref{lemma-gysin-flat-pullback} et
\ref{lemma-gysin-commutes-gysin}.
\end{proof}

\begin{lemma}
\label{lemma-push-proper-bivariant}
Dans la situation \ref{situation-setup}, soient $f : X \to Y$ et $g : Y \to Z$
des morphismes de bons espaces algébriques au-dessus de $B$. Soit $c \in A^p(X \to Z)$
et supposons $f$ propre. Alors la règle qui associe à $X' \to X$
$\alpha \longmapsto f_*(c \cap \alpha)$ est une classe bivariante de degré
$p$.
\end{lemma}

\begin{proof}
Cela découle des lemmes \ref{lemma-compose-pushforward},
\ref{lemma-flat-pullback-proper-pushforward} et \ref{lemma-closed-in-X-gysin}.
\end{proof}

\noindent
Ici, nous voyons que $c_1(\mathcal{L})$ est au centre de $A^*(X)$.

\begin{lemma}
\label{lemma-c1-center}
Dans la situation \ref{situation-setup}, soit $X/B$ bon. Soit
$\mathcal{L}$ un $\mathcal{O}_X$-module inversible. Alors $c_1(\mathcal{L})
\in A^1(X)$ commute avec tout élément $c \in A^p(X)$.
\end{lemma}

\begin{proof}
Soit $p : L \to X$ comme dans le lemme \ref{lemma-linebundle} et soit $o : X
\to L$ la section nulle. Remarquons que $p^*\mathcal{L}^{\otimes -1}$ possède une
section canonique dont le schéma des zéros est exactement le diviseur de
Cartier effectif $o(X)$. Soit $\alpha \in \CH_k(X)$. Nous avons alors
$$
p^*(c_1(\mathcal{L}^{\otimes -1}) \cap \alpha) =
c_1(p^*\mathcal{L}^{\otimes -1}) \cap p^*\alpha =
o_* o^* p^*\alpha
$$
par les lemmes \ref{lemma-flat-pullback-cap-c1} et
\ref{lemma-relative-effective-cartier}. Puisque $c$ est une classe bivariante,
nous avons
\begin{align*}
p^*(c \cap c_1(\mathcal{L}^{\otimes -1}) \cap \alpha)
& =
c \cap p^*(c_1(\mathcal{L}^{\otimes -1}) \cap \alpha) \\
& =
c \cap o_* o^* p^*\alpha \\
& =
o_* o^* p^*(c \cap \alpha) \\
& =
p^*(c_1(\mathcal{L}^{\otimes -1}) \cap c \cap \alpha)
\end{align*}
(la dernière égalité résulte de ce qui précède, appliqué à $c \cap \alpha$). Puisque
$p^*$ est injectif d'après un lemme cité plus haut, nous obtenons que
$c_1(\mathcal{L}^{\otimes -1})$ est au centre de $A^*(X)$. Cela prouve le
lemme.
\end{proof}

\noindent
Voici un critère pour savoir quand une classe bivariante est nulle.

\begin{lemma}
\label{lemma-bivariant-zero}
Dans la situation \ref{situation-setup}, soit $X/B$ bon. Soit $c
\in A^p(X)$. Alors $c$ est nul si et seulement si $c \cap [Y] = 0$ dans
$\CH_*(Y)$ pour tout espace algébrique intègre $Y$ localement de type fini
sur $X$.
\end{lemma}

\begin{proof}
Le sens direct est clair. Réciproquement, supposons que $c \cap [Y] = 0$ dans
$\CH_*(Y)$ pour tout espace algébrique intègre $Y$ localement de type fini
sur $X$. Soit $X' \to X$ localement de type fini. Soit $\alpha \in
\CH_k(X')$. Écrivons $\alpha = \sum n_i [Y_i]$, les $Y_i \subset X'$ formant une
famille localement finie de sous-schémas fermés intègres de $\delta$-dimension $k$. Alors $\alpha$ est l'image directe du cycle
$\alpha' = \sum n_i[Y_i]$ sur $X'' = \coprod Y_i$ par le morphisme propre
$X'' \to X'$. Par les propriétés des classes bivariantes, il suffit de prouver
que $c \cap \alpha' = 0$ dans $\CH_{k - p}(X'')$. Nous avons $\CH_{k - p}(X'')
= \prod \CH_{k - p}(Y_i)$, comme il résulte immédiatement des définitions. Les applications
de projection $\CH_{k - p}(X'') \to \CH_{k - p}(Y_i)$ sont données par image réciproque
plate. Puisque l'intersection avec $c$ commute avec l'image réciproque plate,
nous voyons qu'il suffit de montrer que $c \cap [Y_i]$ est nul dans $\CH_{k -
p}(Y_i)$, ce qui est vrai par hypothèse.
\end{proof}















\section{Formule du fibré projectif}
\label{section-projective-space-bundle-formula}

\noindent
Dans la situation \ref{situation-setup}, soit $X/B$ bon. Considérons un
$\mathcal{O}_X$-module localement libre de type fini $\mathcal{E}$, de rang
$r$. Par convention, le {\it fibré projectif associé à $\mathcal{E}$} est le
morphisme
$$
\xymatrix{
\mathbf{P}(\mathcal{E}) =
\underline{\text{Proj}}_X(\text{Sym}^*(\mathcal{E}))
\ar[r]^-\pi
& X
}
$$
au-dessus de $X$, muni de $\mathcal{O}_{\mathbf{P}(\mathcal{E})}(1)$ normalisé
par l'égalité $\pi_*(\mathcal{O}_{\mathbf{P}(\mathcal{E})}(1)) = \mathcal{E}$.
En particulier, il existe une surjection $\pi^*\mathcal{E} \to
\mathcal{O}_{\mathbf{P}(\mathcal{E})}(1)$. Nous dirons informellement « soit
$(\pi : P \to X, \mathcal{O}_P(1))$ le fibré projectif associé à
$\mathcal{E}$ » pour désigner la situation où $P = \mathbf{P}(\mathcal{E})$ et
$\mathcal{O}_P(1) = \mathcal{O}_{\mathbf{P}(\mathcal{E})}(1)$.

\begin{lemma}
\label{lemma-cap-projective-bundle}
Dans la situation \ref{situation-setup}, soit $X/B$ bon. Soit $\mathcal{E}$ un
$\mathcal{O}_X$-module localement libre de type fini ; ce module
$\mathcal{E}$ est de rang $r$. Soit $(\pi : P \to X, \mathcal{O}_P(1))$ le
fibré projectif associé à $\mathcal{E}$. Pour tout $\alpha \in \CH_k(X)$,
l'élément
$$
\pi_*\left(
c_1(\mathcal{O}_P(1))^s \cap \pi^*\alpha
\right)
\in
\CH_{k + r - 1 - s}(X)
$$
est $0$ si $s < r - 1$ et égal à $\alpha$ si $s = r - 1$.
\end{lemma}

\begin{proof}
Soit $Z \subset X$ un sous-espace fermé intègre de $\delta$-dimension $k$.
Nous démontrerons le lemme pour $\alpha = [Z]$. Nous omettons l'argument qui
déduit le cas général de ce cas particulier ; indication : procéder comme dans
la remarque \ref{remark-infinite-sums-rational-equivalences}.

\medskip\noindent
Soit $P_Z = P \times_X Z$ le changement de base ; bien entendu, $\pi_Z : P_Z
\to Z$ est le fibré projectif associé à $\mathcal{E}|_Z$, et l'image réciproque
de $\mathcal{O}_P(1)$ est le module inversible correspondant sur $P_Z$. Puisque
$c_1(\mathcal{O}_P(1) \cap -$ et $\pi^*$ sont des classes bivariantes par les
lemmes \ref{lemma-cap-c1-bivariant} et \ref{lemma-flat-pullback-bivariant},
il vient
$$
\pi_*\left(
c_1(\mathcal{O}_P(1))^s \cap \pi^*[Z]
\right)
=
(Z \to X)_*\pi_{Z, *}\left(
c_1(\mathcal{O}_{P_Z}(1))^s \cap \pi_Z^*[Z]
\right)
$$
Il suffit donc de prouver le lemme dans le cas où $X$ est intègre et $\alpha
= [X]$.

\medskip\noindent
Supposons $X$ intègre, $\dim_\delta(X) = k$ et $\alpha = [X]$. Remarquons que
$\pi^*[X] = [P]$, puisque $P$ est intègre de $\delta$-dimension $r - 1$. Si
$s < r - 1$, alors, par construction, $c_1(\mathcal{O}_P(1))^s \cap [P]$ est
un $(k + r - 1 - s)$-cycle. Son image directe est donc nulle pour des raisons
de dimension.

\medskip\noindent
Soit $s = r - 1$. L'argument précédent donne
$\pi_*(c_1(\mathcal{O}_P(1))^s \cap [P]) = n [X]$ pour un certain $n \in
\mathbf{Z}$. Nous voulons montrer que $n = 1$. Pour les mêmes raisons de
dimension, il suffit de le faire après remplacement de $X$ par un ouvert
dense. Nous pouvons donc supposer que $X$ est un schéma, et le résultat découle
alors du lemme \ref{chow-lemma-cap-projective-bundle} de l'Homologie de Chow.
\end{proof}

\begin{lemma}[Formule du fibré projectif]
\label{lemma-chow-ring-projective-bundle}
Soit $(S, \delta)$ comme dans la situation \ref{situation-setup}. Soit $X$
localement de type fini sur $S$. Soit $\mathcal{E}$ un $\mathcal{O}_X$-module
localement libre de type fini ; ce module $\mathcal{E}$ est de rang $r$. Soit
$(\pi : P \to X, \mathcal{O}_P(1))$ le fibré projectif associé à
$\mathcal{E}$. L'application
$$
\bigoplus\nolimits_{i = 0}^{r - 1}
\CH_{k + i}(X)
\longrightarrow
\CH_{k + r - 1}(P),
$$
$$
(\alpha_0, \ldots, \alpha_{r-1})
\longmapsto
\pi^*\alpha_0 +
c_1(\mathcal{O}_P(1)) \cap \pi^*\alpha_1
+ \ldots +
c_1(\mathcal{O}_P(1))^{r - 1} \cap \pi^*\alpha_{r-1}
$$
est un isomorphisme.
\end{lemma}

\begin{proof}
Fixons $k \in \mathbf{Z}$. Montrons d'abord que l'application est injective.
Supposons que $(\alpha_0, \ldots, \alpha_{r - 1})$ soit un élément du membre de
gauche dont l'image est nulle. Le lemme \ref{lemma-cap-projective-bundle} donne
$$
0 = \pi_*(\pi^*\alpha_0 +
c_1(\mathcal{O}_P(1)) \cap \pi^*\alpha_1
+ \ldots +
c_1(\mathcal{O}_P(1))^{r - 1} \cap \pi^*\alpha_{r-1})
= \alpha_{r - 1}
$$
Ensuite,
$$
0 = \pi_*(c_1(\mathcal{O}_P(1)) \cap (\pi^*\alpha_0 +
c_1(\mathcal{O}_P(1)) \cap \pi^*\alpha_1
+ \ldots +
c_1(\mathcal{O}_P(1))^{r - 2} \cap \pi^*\alpha_{r - 2}))
= \alpha_{r - 2}
$$
et ainsi de suite. L'application est donc injective.

\medskip\noindent
Pour démontrer la surjectivité, nous raisonnerons exactement comme dans la
preuve du lemme \ref{lemma-pullback-affine-fibres-surjective}, afin de nous
ramener au cas des schémas. Le lecteur est invité à omettre cette preuve.

\medskip\noindent
Soit $\beta \in \CH_{k + r - 1}(P)$. Écrivons $\beta = \sum m_j[W_j]$ avec
$m_j \not = 0$ et $W_j$ des sous-espaces fermés intègres distincts par paires
et de $\delta$-dimension $k + r$. La famille $\{W_j\}$ est alors localement
finie dans $P$. Soit $Z_j \subset X$ l'« image » de $W_j$ comme dans le lemme
\ref{lemma-proper-image}. Pour toute partie ouverte quasi-compacte $U \subset
X$, l'intersection $\pi^{-1}(U) \cap W_j$ n'est non vide que pour un nombre
fini de $j$. Par conséquent, les images $Z_j$ forment une famille localement
finie de sous-espaces fermés intègres de $X$.

\medskip\noindent
Considérons les diagrammes cartésiens
$$
\xymatrix{
P_j \ar[r] \ar[d]_{\pi_j} & P \ar[d]^\pi \\
Z_j \ar[r] & X
}
$$
Supposons que $[W_j] \in Z_{k + r - 1}(P_j)$ soit rationnellement équivalent à
$$
\pi_j^*\alpha_{j, 0} +
c_1(\mathcal{O}(1)) \cap \pi_j^*\alpha_{j, 1} +
\ldots +
c_1(\mathcal{O}(1))^{r - 1} \cap \pi_j^*\alpha_{j, r - 1}
$$
pour certains $(k + i)$-cycles $\alpha_{j, i} \in \CH_{k + i}(Z_j)$. Alors
$\alpha_i = \sum m_j \beta_{j, i}$ sera un $(k + i)$-cycle sur $X$, et
$$
\pi^*\alpha_0 +
c_1(\mathcal{O}(1)) \cap \pi^*\alpha_1 +
\ldots +
c_1(\mathcal{O}(1))^{r - 1} \cap \pi^*\alpha_{r - 1}
$$
sera rationnellement équivalent à $\beta$ (voir la remarque
\ref{remark-infinite-sums-rational-equivalences}). Cela nous réduit au cas $X$
intègre et $\alpha = [W]$ pour un sous-schéma fermé intègre de $P$ dominant
$X$. En particulier, nous pouvons supposer que $d = \dim_\delta(X) < \infty$.

\medskip\noindent
Nous pouvons donc raisonner par récurrence sur $d = \dim_\delta(X)$. Si $d <
k$, alors $\CH_{k + r - 1}(X) = 0$ et le lemme est établi ; c'est le cas
initial de la récurrence. Considérons une partie ouverte non vide $U \subset
X$. Supposons que
l'on puisse montrer que
$$
\beta|_{\pi^{-1}(U)} =
\pi^*\alpha_0 +
c_1(\mathcal{O}(1)) \cap \pi^*\alpha_1 +
\ldots +
c_1(\mathcal{O}(1))^{r - 1} \cap \pi^*\alpha_{r - 1}
$$
pour certains $\alpha_i \in Z_{k + i}(U)$. Le lemme
\ref{lemma-exact-sequence-open} donne $\alpha_i = \alpha'_i|_U$ pour certains
$\alpha'_i \in Z_{k + i}(X)$. Par les suites exactes $\CH_{k +
i}(\pi^{-1}(X \setminus U)) \to \CH_{k + i}(P) \to \CH_{k + i}(\pi^{-1}(U))$
du lemme \ref{lemma-restrict-to-open}, il vient que
$$
\beta -
\left(\pi^*\alpha'_0 +
c_1(\mathcal{O}(1)) \cap \pi^*\alpha'_1 +
\ldots +
c_1(\mathcal{O}(1))^{r - 1} \cap \pi^*\alpha'_{r - 1}\right)
$$
provient d'un cycle $\beta' \in \CH_{k + r}(\pi^{-1}(X \setminus U))$. Puisque
$\dim_\delta(X \setminus U) < d$, on conclut par récurrence sur $d$.

\medskip\noindent
En particulier, après remplacement de $X$ par un ouvert convenable, nous
pouvons supposer que $X$ est un schéma ; le problème est alors ramené au lemme
\ref{chow-lemma-chow-ring-projective-bundle} de l'Homologie de Chow.
\end{proof}

\begin{lemma}
\label{lemma-vectorbundle}
Dans la situation \ref{situation-setup}, soit $X/B$ bon. Soit $\mathcal{E}$ un
faisceau localement libre de type fini et de rang $r$ sur $X$. Soit
$$
p :
E = \underline{\Spec}(\text{Sym}^*(\mathcal{E}))
\longrightarrow
X
$$
le fibré vectoriel associé au-dessus de $X$. Alors $p^* : \CH_k(X) \to
\CH_{k + r}(E)$ est un isomorphisme pour tout $k$.
\end{lemma}

\begin{proof}
(Pour le cas des fibrés en droites, voir le lemme \ref{lemma-linebundle}.) La
surjectivité résulte du lemme \ref{lemma-pullback-affine-fibres-surjective}.
Soit $(\pi  : P \to X, \mathcal{O}_P(1))$ le fibré en espaces projectifs
associé au faisceau localement libre de type fini $\mathcal{E} \oplus
\mathcal{O}_X$. Soit $s \in \Gamma(P, \mathcal{O}_P(1))$ la section globale
correspondant à $(0, 1) \in \Gamma(X, \mathcal{E} \oplus \mathcal{O}_X)$.
Soit $D = Z(s) \subset P$. Remarquons que $(\pi|_D : D \to X ,
\mathcal{O}_P(1)|_D)$ est le fibré en espaces projectifs associé à
$\mathcal{E}$. Posons $\pi_D = \pi|_D$ et $\mathcal{O}_D(1) =
\mathcal{O}_P(1)|_D$. De plus, $D$ est un diviseur de Cartier effectif sur $P$.
Ainsi, $\mathcal{O}_P(D) = \mathcal{O}_P(1)$ (voir Diviseurs sur les espaces,
lemme \ref{spaces-divisors-lemma-characterize-OD}). Il existe aussi un
isomorphisme $E \cong P \setminus D$. Notons $j : E \to P$ l'immersion ouverte
correspondante. Pour l'injectivité, utilisons le fait que les éléments du noyau
de
$$
j^* :
\CH_{k + r}(P)
\longrightarrow
\CH_{k + r}(E)
$$
sont les cycles à support dans le diviseur de Cartier effectif $D$ ; voir le
lemme \ref{lemma-restrict-to-open}. Ainsi, si $p^*\alpha = 0$, alors
$\pi^*\alpha = i_*\beta$ pour un certain $\beta \in \CH_{k + r}(D)$. Le lemme
\ref{lemma-chow-ring-projective-bundle} permet d'écrire
$$
\beta = \pi_D^*\beta_0 +
\ldots + c_1(\mathcal{O}_D(1))^{r - 1} \cap \pi_D^* \beta_{r - 1}.
$$
pour certains $\beta_i \in \CH_{k + i}(X)$. Les lemmes
\ref{lemma-relative-effective-cartier} et \ref{lemma-pushforward-cap-c1}
donnent alors
$$
\pi^*\alpha = i_*\beta =
c_1(\mathcal{O}_P(1)) \cap \pi^*\beta_0 +
\ldots +
c_1(\mathcal{O}_D(1))^r \cap \pi^*\beta_{r - 1}.
$$
Puisque le rang de $\mathcal{E} \oplus \mathcal{O}_X$ est $r + 1$, cela
contredit le lemme \ref{lemma-pushforward-cap-c1}, à moins que tous les
$\alpha$ et tous les $\beta_i$ ne soient nuls.
\end{proof}








\section{Les classes de Chern d'un fibré vectoriel}
\label{section-chern-classes-vector-bundles}

\noindent
Cette section est l'analogue des sections
\ref{chow-section-chern-classes-vector-bundles} et
\ref{chow-section-intersecting-chern-classes} de l'Homologie de Chow.
Cependant, contrairement à ce qui y est fait, nous définissons directement les
classes de Chern d'un fibré vectoriel comme des classes bivariantes. Cela
économise un travail considérable.

\begin{lemma}
\label{lemma-segre-classes}
Dans la situation \ref{situation-setup}, soit $X/B$ bon. Soit $\mathcal{E}$ un
faisceau localement libre de type fini et de rang $r$ sur $X$. Soit
$(\pi : P \to X, \mathcal{O}_P(1))$ le fibré en espaces projectifs associé à
$\mathcal{E}$. Pour tout morphisme $X' \to X$ de bons espaces algébriques
au-dessus de $B$, il existe des applications uniques
$$
c_i(\mathcal{E}) \cap - : \CH_k(X') \longrightarrow \CH_{k - i}(X'),\quad
i = 0, \ldots, r
$$
telles que, pour $\alpha \in \CH_k(X')$, nous ayons $c_0(\mathcal{E}) \cap
\alpha = \alpha$ et
$$
\sum\nolimits_{i = 0, \ldots, r}
(-1)^i c_1(\mathcal{O}_{P'}(1))^i \cap
(\pi')^*\left(c_{r - i}(\mathcal{E}) \cap \alpha\right) = 0
$$
où $\pi' : P' \to X'$ est le changement de base de $\pi$. De plus, ces applications
définissent une classe bivariante $c_i(\mathcal{E})$ de degré $i$ sur $X$.
\end{lemma}

\begin{proof}
L'existence et l'unicité des applications $c_i(\mathcal{E}) \cap -$ résultent
immédiatement du lemme \ref{lemma-chow-ring-projective-bundle} et de la
description donnée de $c_0(\mathcal{E})$. Pour tout $i \in \mathbf{Z}$, la
règle qui associe à tout morphisme $X' \to X$ de bons espaces algébriques
au-dessus de $B$ l'application
$$
t_i(\mathcal{E}) \cap - :
\CH_k(X') \longrightarrow \CH_{k - i}(X'),\quad
\alpha \longmapsto
\pi'_*(c_1(\mathcal{O}_{P'}(1))^{r - 1 + i} \cap (\pi')^*\alpha)
$$
est une classe bivariante\footnote{À des signes près, ce sont les classes de
Segre de $\mathcal{E}$.}, par les lemmes \ref{lemma-cap-c1-bivariant},
\ref{lemma-flat-pullback-bivariant} et \ref{lemma-push-proper-bivariant}. Par le
lemme \ref{lemma-cap-projective-bundle}, nous avons $t_i(\mathcal{E}) = 0$ pour
$i < 0$ et $t_0(\mathcal{E}) = 1$. En appliquant l'image directe à l'équation
de l'énoncé du lemme, le lemme \ref{lemma-cap-projective-bundle} donne
$$
(-1)^r t_1(\mathcal{E}) + (-1)^{r - 1}c_1(\mathcal{E}) = 0
$$
En particulier, $c_1(\mathcal{E})$ est une classe bivariante. Si nous
multiplions l'équation de l'énoncé du lemme par
$c_1(\mathcal{O}_{P'}(1))$ et prenons l'image directe du résultat sur $X'$, il
vient
$$
(-1)^r t_2(\mathcal{E}) +
(-1)^{r - 1} t_1(\mathcal{E}) \cap c_1(\mathcal{E}) +
(-1)^{r - 2} c_2(\mathcal{E}) = 0
$$
Comme précédemment, nous concluons que $c_2(\mathcal{E})$ est une classe
bivariante. Et ainsi de suite.
\end{proof}

\begin{definition}
\label{definition-chern-classes}
Dans la situation \ref{situation-setup}, soit $X/B$ bon. Soit $\mathcal{E}$ un
faisceau localement libre de type fini et de rang $r$ sur $X$. Pour $i = 0,
\ldots, r$, la {\it $i$-ième classe de Chern de $\mathcal{E}$} est la classe
bivariante $c_i(\mathcal{E}) \in A^i(X)$ de degré $i$ construite dans le lemme
\ref{lemma-segre-classes}. La {\it classe totale de Chern de $\mathcal{E}$} est
la somme formelle
$$
c(\mathcal{E}) =
c_0(\mathcal{E}) + c_1(\mathcal{E}) + \ldots + c_r(\mathcal{E})
$$
qui est considérée comme une classe bivariante non homogène sur $X$.
\end{definition}

\noindent
Par commodité, nous posons souvent $c_i(\mathcal{E}) = 0$ pour
$i > r$ et $i < 0$. Par définition, nous avons $c_0(\mathcal{E}) = 1 \in
A^0(X)$. Voici une vérification élémentaire.

\begin{lemma}
\label{lemma-first-chern-class}
Dans la situation \ref{situation-setup}, soit $X/B$ bon. Soit $\mathcal{L}$ un
$\mathcal{O}_X$-module inversible. La première classe de Chern de
$\mathcal{L}$ sur $X$, définie dans la définition
\ref{definition-chern-classes}, est égale à la classe bivariante du lemme
\ref{lemma-cap-c1-bivariant}.
\end{lemma}

\begin{proof}
En effet, dans ce cas $P = \mathbf{P}(\mathcal{L}) = X$ et $\mathcal{O}_P(1) =
\mathcal{L}$ par notre normalisation du fibré projectif ; voir la section
\ref{section-projective-space-bundle-formula}. L'équation du lemme
\ref{lemma-segre-classes} s'écrit donc
$$
(-1)^0 c_1(\mathcal{L})^0 \cap c^{new}_1(\mathcal{L}) \cap \alpha +
(-1)^1 c_1(\mathcal{L})^1 \cap c^{new}_0(\mathcal{L}) \cap \alpha = 0
$$
où $c_i^{new}(\mathcal{L})$ est défini comme dans la définition
\ref{definition-chern-classes}. Puisque $c_0^{new}(\mathcal{L}) = 1$ et
$c_1(\mathcal{L})^0 = 1$, la conclusion suit.
\end{proof}

\noindent
Nous voyons ensuite que les classes de Chern appartiennent au centre de
l'anneau de cohomologie de Chow bivariante $A^*(X)$.

\begin{lemma}
\label{lemma-cap-commutative-chern}
Dans la situation \ref{situation-setup}, soit $X/B$ bon. Soit $\mathcal{E}$ un
$\mathcal{O}_X$-module localement libre de rang $r$. Alors
$c_j(\mathcal{L}) \in A^j(X)$ commute avec tout élément $c \in A^p(X)$. En
particulier, si $\mathcal{F}$ est un second $\mathcal{O}_X$-module localement
libre sur $X$, de rang $s$, alors
$$
c_i(\mathcal{E}) \cap c_j(\mathcal{F}) \cap \alpha
=
c_j(\mathcal{F}) \cap c_i(\mathcal{E}) \cap \alpha
$$
comme éléments de $\CH_{k - i - j}(X)$ pour tout $\alpha \in \CH_k(X)$.
\end{lemma}

\begin{proof}
Soit $X' \to X$ un morphisme de bons espaces algébriques au-dessus de $B$.
Soit $\alpha \in \CH_k(X')$. Écrivons $\alpha_j = c_j(\mathcal{E}) \cap
\alpha$, de sorte que $\alpha_0 = \alpha$. Le lemme
\ref{lemma-segre-classes} donne
$$
\sum\nolimits_{i = 0}^r
(-1)^i c_1(\mathcal{O}_{P'}(1))^i \cap
(\pi')^*(\alpha_{r - i}) = 0
$$
dans le groupe de Chow du fibré projectif $(\pi' : P' \to X',
\mathcal{O}_{P'}(1))$ associé à $(X' \to X)^*\mathcal{E}$. En appliquant $c
\cap -$, puis en utilisant le lemme \ref{lemma-c1-center} et les propriétés des
classes bivariantes, nous obtenons
$$
\sum\nolimits_{i = 0}^r
(-1)^i c_1(\mathcal{O}_{P'}(1))^i \cap
\pi^*(c \cap \alpha_{r - i}) = 0
$$
dans le groupe de Chow de $P'$. L'unicité affirmée dans le lemme
\ref{lemma-segre-classes} montre donc que $c \cap \alpha_j$ est égal à
$c_j(\mathcal{E}) \cap (c \cap \alpha)$. Cela prouve le lemme.
\end{proof}

\begin{remark}
\label{remark-extend-to-finite-locally-free}
Dans la situation \ref{situation-setup}, soit $X/B$ bon. Soit $\mathcal{E}$ un
$\mathcal{O}_X$-module localement libre de type fini. Si le rang de
$\mathcal{E}$ n'est pas constant, nous pouvons encore définir les classes de
Chern de $\mathcal{E}$. En effet, dans ce cas, on peut écrire
$$
X = X_0 \amalg X_1 \amalg X_2 \amalg \ldots
$$
où $X_r \subset X$ est le sous-espace ouvert et fermé sur lequel le rang de
$\mathcal{E}$ est $r$. Si $X' \to X$ est un morphisme de bons espaces
algébriques au-dessus de $B$, l'image réciproque fournit une décomposition
correspondante de $X'$, et les définitions donnent
$$
\CH_*(X') = \prod\nolimits_{r \geq 0} \CH_*(X'_r)
$$
Nous définissons alors $c_i(\mathcal{E})$ comme la classe bivariante qui
préserve ces décompositions en produits directs et agit sur les facteurs par
les opérations déjà définies $c_i(\mathcal{E}|_{X_r}) \cap -$. Remarquons que,
dans ce cadre, $c_i(\mathcal{E})$ peut être non nul pour une infinité de $i$.
\end{remark}












\section{Relations polynomiales entre les classes de Chern}
\label{section-relations-chern-classes}

\noindent
Dans la situation \ref{situation-setup}, soit $X/B$ bon. Soient
$\mathcal{E}_i$ les membres d'une famille finie de $\mathcal{O}_X$-modules
localement libres de type fini. Le lemme \ref{lemma-cap-commutative-chern}
montre que les classes de Chern
$$
c_j(\mathcal{E}_i) \in A^*(X)
$$
engendrent une sous-$\mathbf{Z}$-algèbre commutative, et même centrale, de la
cohomologie de Chow $A^*(X)$. Il est donc bien défini qu'un polynôme en ces
classes de Chern soit nul, ou que deux tels polynômes soient égaux. À titre
d'exemple, dire que $c_1(\mathcal{E}_1)^5 +
c_2(\mathcal{E}_2)c_3(\mathcal{E}_3) = 0$ signifie que les opérations
$$
\CH_k(Y) \longrightarrow \CH_{k - 5}(Y), \quad
\alpha \longmapsto
c_1(\mathcal{E}_1)^5 \cap \alpha +
c_2(\mathcal{E}_2) \cap c_3(\mathcal{E}_3) \cap \alpha
$$
sont nulles pour tout morphisme $f : Y \to X$ de bons espaces algébriques
au-dessus de $B$. Par le lemme \ref{lemma-bivariant-zero}, cette condition
équivaut à demander que, pour tout morphisme $f : Y \to X$, où $Y$ est un
espace algébrique intègre localement de type fini sur $X$, le cycle
$$
c_1(\mathcal{E}_1)^5 \cap [Y] +
c_2(\mathcal{E}_2) \cap c_3(\mathcal{E}_3) \cap [Y]
$$
est nul dans $\CH_{\dim(Y) - 5}(Y)$.

\medskip\noindent
Un exemple concret est la relation
$$
c_1(\mathcal{L} \otimes_{\mathcal{O}_X} \mathcal{N})
=
c_1(\mathcal{L}) + c_1(\mathcal{N})
$$
établie dans le lemme \ref{lemma-c1-cap-additive}. Plus généralement, voici ce
qui se passe lorsqu'on tensorise un faisceau localement libre quelconque par un
faisceau inversible.

\begin{lemma}
\label{lemma-chern-classes-E-tensor-L}
Dans la situation \ref{situation-setup}, soit $X/B$ bon. Soit $\mathcal{E}$ un
faisceau localement libre de type fini et de rang $r$ sur $X$. Soit
$\mathcal{L}$ un faisceau inversible sur $X$. Alors
\begin{equation}
\label{equation-twist}
c_i({\mathcal E} \otimes {\mathcal L})
=
\sum\nolimits_{j = 0}^i
\binom{r - i + j}{j} c_{i - j}({\mathcal E}) c_1({\mathcal L})^j
\end{equation}
dans $A^*(X)$.
\end{lemma}

\begin{proof}
La preuve est identique à celle du lemme
\ref{chow-lemma-chern-classes-E-tensor-L} de l'Homologie de Chow, les lemmes
\ref{lemma-bivariant-zero} et \ref{lemma-segre-classes} remplaçant ceux qui y
sont utilisés.
\end{proof}












\section{Additivité des classes de Chern}
\label{section-additivity-chern-classes}

\noindent
Cette section est l'analogue de la section
\ref{chow-section-additivity-chern-classes} de l'Homologie de Chow.

\begin{lemma}
\label{lemma-get-rid-of-trivial-subbundle}
Dans la situation \ref{situation-setup}, soit $X/B$ bon. Soient
$\mathcal{E}$, $\mathcal{F}$ des faisceaux localement libres de type fini sur
$X$, de rangs respectifs $r$, $r - 1$, qui s'insèrent dans une suite exacte
courte
$$
0 \to \mathcal{O}_X \to \mathcal{E} \to \mathcal{F} \to 0
$$
Alors
$$
c_r(\mathcal{E}) = 0, \quad
c_j(\mathcal{E}) = c_j(\mathcal{F}), \quad j = 0, \ldots, r - 1
$$
dans $A^*(X)$.
\end{lemma}

\begin{proof}
La preuve est identique à celle du lemme
\ref{chow-lemma-get-rid-of-trivial-subbundle} de l'Homologie de Chow, les
lemmes \ref{lemma-bivariant-zero},
\ref{lemma-relative-effective-cartier}, \ref{lemma-pushforward-cap-c1} et
\ref{lemma-segre-classes} remplaçant ceux qui y sont utilisés.
\end{proof}

\begin{lemma}
\label{lemma-additivity-invertible-subsheaf}
Dans la situation \ref{situation-setup}, soit $X/B$ bon. Soient
$\mathcal{E}$, $\mathcal{F}$ des faisceaux localement libres de type fini sur
$X$, de rangs respectifs $r$, $r - 1$, qui s'insèrent dans une suite exacte
courte
$$
0 \to \mathcal{L} \to \mathcal{E} \to \mathcal{F} \to 0
$$
où $\mathcal{L}$ est un faisceau inversible. Alors
$$
c(\mathcal{E}) = c(\mathcal{L}) c(\mathcal{F})
$$
dans $A^*(X)$.
\end{lemma}

\begin{proof}
La preuve est identique à celle du lemme
\ref{chow-lemma-additivity-invertible-subsheaf} de l'Homologie de Chow, les
lemmes \ref{lemma-get-rid-of-trivial-subbundle} et
\ref{lemma-chern-classes-E-tensor-L} remplaçant ceux qui y sont utilisés.
\end{proof}

\begin{lemma}
\label{lemma-additivity-chern-classes}
Dans la situation \ref{situation-setup}, soit $X/B$ bon. Supposons que
$\mathcal{E}$ s'insère dans une suite exacte
$$
0
\to
\mathcal{E}_1
\to
\mathcal{E}
\to
\mathcal{E}_2
\to
0
$$
de faisceaux localement libres de type fini $\mathcal{E}_i$, de rang $r_i$.
Les classes totales de Chern vérifient
$$
c(\mathcal{E}) = c(\mathcal{E}_1) c(\mathcal{E}_2)
$$
dans $A^*(X)$.
\end{lemma}

\begin{proof}
La preuve est identique à celle du lemme
\ref{chow-lemma-additivity-chern-classes} de l'Homologie de Chow, les lemmes
\ref{lemma-bivariant-zero}, \ref{lemma-additivity-invertible-subsheaf} et
\ref{lemma-segre-classes} remplaçant ceux qui y sont utilisés.
\end{proof}

\begin{lemma}
\label{lemma-chern-filter-by-linebundles}
Dans la situation \ref{situation-setup}, soit $X/B$ bon. Soient
${\mathcal L}_i$, $i = 1, \ldots, r$, des $\mathcal{O}_X$-modules inversibles.
Soit $\mathcal{E}$ un $\mathcal{O}_X$-module localement libre, muni d'une
filtration
$$
0 = \mathcal{E}_0 \subset \mathcal{E}_1 \subset \mathcal{E}_2
\subset \ldots \subset \mathcal{E}_r = \mathcal{E}
$$
tel que $\mathcal{E}_i/\mathcal{E}_{i - 1} \cong \mathcal{L}_i$. Posons
$c_1({\mathcal L}_i) = x_i$. Alors
$$
c(\mathcal{E})
=
\prod\nolimits_{i = 1}^r (1 + x_i)
$$
dans $A^*(X)$.
\end{lemma}

\begin{proof}
Appliquons le lemme \ref{lemma-additivity-invertible-subsheaf} et raisonnons
par récurrence.
\end{proof}



















\section{Le principe de scindage}
\label{section-splitting-principle}

\noindent
Cette section est l'analogue de la section
\ref{chow-section-additivity-chern-classes} de l'Homologie de Chow.

\begin{lemma}
\label{lemma-splitting-principle}
Dans la situation \ref{situation-setup}, soit $X/B$ bon. Soient
$\mathcal{E}_i$ les membres d'une famille finie de $\mathcal{O}_X$-modules
localement libres, de rang $r_i$. Il existe un morphisme projectif et plat
$\pi : P \to X$, de dimension relative $d$, tel que
\begin{enumerate}
\item pour tout morphisme $f : Y \to X$ de bons espaces algébriques au-dessus de $B$,
l'application $\pi_Y^* : \CH_*(Y) \to \CH_{* + d}(Y \times_X P)$ est
injective, et
\item chaque $\pi^*\mathcal{E}_i$ possède une filtration dont les quotients
successifs $\mathcal{L}_{i, 1}, \ldots, \mathcal{L}_{i, r_i}$ sont des
${\mathcal O}_P$-modules inversibles.
\end{enumerate}
\end{lemma}

\begin{proof}
Raisonnons par récurrence sur l'entier $r = \sum r_i$. Si $r = 0$, nous
pouvons prendre $\pi = \text{id}_X$. Si $r_i = 1$ pour tout $i$, nous pouvons
encore prendre $\pi = \text{id}_X$. Supposons que $r_{i_0} > 1$ pour un certain
$i_0$. Soit $(\pi : P \to X, \mathcal{O}_P(1))$ le fibré projectif associé à
$\mathcal{E}_{i_0}$. L'homomorphisme canonique $\pi^*\mathcal{E}_{i_0} \to
\mathcal{O}_P(1)$ est surjectif ; son noyau $\mathcal{E}'_{i_0}$ est donc
localement libre de type fini et de rang $r_{i_0} - 1$. Remarquons que
$\pi_Y^*$ est injectif pour tout morphisme $f : Y \to X$ de bons espaces
algébriques au-dessus de $B$ ; voir le lemme
\ref{lemma-chow-ring-projective-bundle}. Il suffit donc de démontrer le lemme
pour $P$ et les faisceaux localement libres $\pi^*\mathcal{E}_i$. Cependant,
puisque nous disposons du sous-fibré $\mathcal{E}_{i_0} \subset
\pi^*\mathcal{E}_{i_0}$ à quotient inversible, il suffit maintenant de traiter
la famille $\{\mathcal{E}_i\}_{i \not = i_0} \cup
\{\mathcal{E}'_{i_0}\}$. Cela diminue $r$ de $1$, et l'hypothèse de récurrence
permet de conclure.
\end{proof}

\noindent
Plutôt que d'expliquer ce qu'affirme le principe de scindage, utilisons-le
dans la preuve de quelques lemmes.

\begin{lemma}
\label{lemma-chern-classes-dual}
Dans la situation \ref{situation-setup}, soit $X/B$ bon. Soit $\mathcal{E}$ un
$\mathcal{O}_X$-module localement libre de type fini, dont le dual est
$\mathcal{E}^\vee$. Alors
$$
c_i(\mathcal{E}^\vee) = (-1)^i c_i(\mathcal{E})
$$
dans $A^i(X)$.
\end{lemma}

\begin{proof}
Choisissons un morphisme $\pi : P \to X$ comme dans le lemme
\ref{lemma-splitting-principle}. Puisque $\pi^*$ est injectif après tout
changement de base, il suffit d'établir la relation entre les classes de Chern
de $\mathcal{E}$ et de $\mathcal{E}^\vee$ après image réciproque sur $P$. Nous
pouvons donc supposer qu'il existe des $\mathcal{O}_X$-modules inversibles
${\mathcal L}_i$, $i = 1, \ldots, r$, et une filtration
$$
0 = \mathcal{E}_0 \subset \mathcal{E}_1 \subset \mathcal{E}_2
\subset \ldots \subset \mathcal{E}_r = \mathcal{E}
$$
telle que $\mathcal{E}_i/\mathcal{E}_{i - 1} \cong \mathcal{L}_i$. On obtient
alors la filtration duale
$$
0 = \mathcal{E}_r^\perp \subset \mathcal{E}_1^\perp \subset \mathcal{E}_2^\perp
\subset \ldots \subset \mathcal{E}_0^\perp = \mathcal{E}^\vee
$$
telle que $\mathcal{E}_{i - 1}^\perp/\mathcal{E}_i^\perp \cong
\mathcal{L}_i^{\otimes -1}$. Posons $x_i = c_1(\mathcal{L}_i)$. Alors
$c_1(\mathcal{L}_i^{\otimes -1}) = - x_i$ par le lemme
\ref{lemma-c1-cap-additive}. Par le lemme
\ref{lemma-chern-filter-by-linebundles}, nous avons
$$
c(\mathcal{E}) = \prod\nolimits_{i = 1}^r (1 + x_i)
\quad\text{et}\quad
c(\mathcal{E}^\vee) = \prod\nolimits_{i = 1}^r (1 - x_i)
$$
dans $A^*(X)$. Le résultat découle d'un calcul formel que nous omettons.
\end{proof}

\begin{lemma}
\label{lemma-chern-classes-tensor-product}
Dans la situation \ref{situation-setup}, soit $X/B$ bon. Soient
$\mathcal{E}$ et $\mathcal{F}$ des $\mathcal{O}_X$-modules localement libres de
type fini, de rangs respectifs $r$ et $s$. Alors
$$
c_1(\mathcal{E} \otimes \mathcal{F})
=
r c_1(\mathcal{F}) + s c_1(\mathcal{E})
$$
$$
c_2(\mathcal{E} \otimes \mathcal{F})
=
r^2 c_2(\mathcal{F}) +
rs c_1(\mathcal{F})c_1(\mathcal{E}) +
s^2 c_2(\mathcal{E})
$$
et ainsi de suite (voir preuve).
\end{lemma}

\begin{proof}
En raisonnant exactement comme dans la preuve du lemme
\ref{lemma-chern-classes-dual}, nous pouvons supposer donnés des
$\mathcal{O}_X$-modules inversibles ${\mathcal L}_i$, $i = 1, \ldots, r$, des
${\mathcal N}_i$, $i = 1, \ldots, s$, et des filtrations
$$
0 = \mathcal{E}_0 \subset \mathcal{E}_1 \subset \mathcal{E}_2
\subset \ldots \subset \mathcal{E}_r = \mathcal{E}
\quad\text{et}\quad
0 = \mathcal{F}_0 \subset \mathcal{F}_1 \subset \mathcal{F}_2
\subset \ldots \subset \mathcal{F}_s = \mathcal{F}
$$
telles que $\mathcal{E}_i/\mathcal{E}_{i - 1} \cong \mathcal{L}_i$ et que
$\mathcal{F}_j/\mathcal{F}_{j - 1} \cong \mathcal{N}_j$. En ordonnant
lexicographiquement les couples $(i, j)$, on obtient une filtration
$$
0 \subset \ldots \subset
\mathcal{E}_i \otimes \mathcal{F}_j
+
\mathcal{E}_{i - 1} \otimes \mathcal{F}
\subset \ldots \subset \mathcal{E} \otimes \mathcal{F}
$$
dont les quotients successifs sont
$$
\mathcal{L}_1 \otimes \mathcal{N}_1,
\mathcal{L}_1 \otimes \mathcal{N}_2,
\ldots,
\mathcal{L}_1 \otimes \mathcal{N}_s,
\mathcal{L}_2 \otimes \mathcal{N}_1,
\ldots,
\mathcal{L}_r \otimes \mathcal{N}_s
$$
Par le lemme \ref{lemma-chern-filter-by-linebundles}, nous avons
$$
c(\mathcal{E}) = \prod (1 + x_i),
\quad
c(\mathcal{F}) = \prod (1 + y_j),
\quad\text{et}\quad
c(\mathcal{F}) = \prod (1 + x_i + y_j),
$$
dans $A^*(X)$. Le résultat découle d'un calcul formel que nous omettons.
\end{proof}











\section{Degrés des zéro-cycles}
\label{section-degree-zero-cycles}

\noindent
Cette section est l'analogue de la section
\ref{chow-section-degree-zero-cycles} de l'Homologie de Chow. Commençons par
définir le degré d'un zéro-cycle sur un espace algébrique propre au-dessus d'un
corps.

\begin{definition}
\label{definition-degree-zero-cycle}
Soit $k$ un corps. Soit $p : X \to \Spec(k)$ un morphisme propre d'espaces
algébriques. Le {\it degré d'un zéro-cycle} sur $X$ est donné par l'image
directe propre
$$
p_* : \CH_0(X) \longrightarrow \CH_0(\Spec(k)) \longrightarrow \mathbf{Z}
$$
(lemme \ref{lemma-proper-pushforward-rational-equivalence}), composée avec
l'isomorphisme naturel $\CH_0(\Spec(k)) \to \mathbf{Z}$ qui envoie $[\Spec(k)]$
à $1$. Notation : $\deg(\alpha)$.
\end{definition}

\noindent
Expliquons cela plus en détail.

\begin{lemma}
\label{lemma-spell-out-degree-zero-cycle}
Soit $k$ un corps. Soit $X$ un espace algébrique propre sur $k$. Soit $\alpha
= \sum n_i[Z_i]$ dans $Z_0(X)$. Alors
$$
\deg(\alpha) = \sum n_i\deg(Z_i)
$$
où $\deg(Z_i)$ est le degré de $Z_i \to \Spec(k)$, c'est-à-dire $\deg(Z_i) =
\dim_k \Gamma(Z_i, \mathcal{O}_{Z_i})$.
\end{lemma}

\begin{proof}
C'est la définition de l'image directe propre (définition
\ref{definition-proper-pushforward}).
\end{proof}

\begin{lemma}
\label{lemma-degrees-and-numerical-intersections}
Soit $k$ un corps. Soit $X$ un espace algébrique propre sur $k$. Soit $Z
\subset X$ un sous-espace fermé de dimension $d$. Soient $\mathcal{L}_1,
\ldots, \mathcal{L}_d$ des $\mathcal{O}_X$-modules inversibles. Alors
$$
(\mathcal{L}_1 \cdots \mathcal{L}_d \cdot Z) =
\deg(
c_1(\mathcal{L}_1) \cap \ldots \cap c_1(\mathcal{L}_1) \cap [Z]_d)
$$
où le membre de gauche est défini dans Espaces sur les corps, définition
\ref{spaces-over-fields-definition-intersection-number}.
\end{lemma}

\begin{proof}
Soient $Z_i \subset Z$, $i = 1, \ldots, t$, les composantes irréductibles de
dimension $d$. Soit $m_i$ la multiplicité de $Z_i$ dans $Z$. Alors $[Z]_d =
\sum m_i[Z_i]$, et $c_1(\mathcal{L}_1) \cap \ldots \cap c_1(\mathcal{L}_d)
\cap [Z]_d$ est la somme des cycles $m_i c_1(\mathcal{L}_1) \cap \ldots \cap
c_1(\mathcal{L}_d) \cap [Z_i]$. Comme on dispose d'une décomposition analogue
pour $(\mathcal{L}_1 \cdots \mathcal{L}_d \cdot Z)$ d'après le lemme
\ref{spaces-over-fields-lemma-numerical-polynomial-leading-term} d'Espaces sur
les corps, il suffit de démontrer le lemme lorsque $Z = X$ est un espace
algébrique intègre propre sur $k$.

\medskip\noindent
D'après le lemme de Chow, il existe un morphisme propre $f : X' \to X$ qui est
un isomorphisme au-dessus d'une partie ouverte dense $U \subset X$ et tel que
$X'$ soit un schéma ; voir le lemme
\ref{spaces-more-morphisms-lemma-chow-noetherian-separated} de Compléments sur
les morphismes d'espaces. Alors $X'$ est un schéma propre sur $k$. Après avoir
remplacé $X'$ par l'adhérence schématique de $f^{-1}(U)$, nous pouvons supposer
que $X'$ est intègre. Alors
$$
(f^*\mathcal{L}_1 \cdots f^*\mathcal{L}_d \cdot X') =
(\mathcal{L}_1 \cdots \mathcal{L}_d \cdot X)
$$
par le lemme \ref{spaces-over-fields-lemma-intersection-number-and-pullback}
d'Espaces sur les corps, et nous avons
$$
f_*(c_1(f^*\mathcal{L}_1) \cap \ldots \cap c_1(f^*\mathcal{L}_d) \cap [Y]) =
c_1(\mathcal{L}_1) \cap \ldots \cap c_1(\mathcal{L}_d) \cap [X]
$$
par le lemme \ref{lemma-pushforward-cap-c1}. Nous pouvons donc remplacer $X$
par $X'$ et supposer que $X$ est un schéma propre sur $k$. Ce cas est établi
dans le lemme \ref{chow-lemma-degrees-and-numerical-intersections} de
l'Homologie de Chow.
\end{proof}















\begin{multicols}{2}[\section{Autres chapitres}]
\noindent
Préliminaires
\begin{enumerate}
\item \hyperref[introduction-section-phantom]{Introduction}
\item \hyperref[conventions-section-phantom]{Conventions}
\item \hyperref[sets-section-phantom]{Théorie des ensembles}
\item \hyperref[categories-section-phantom]{Catégories}
\item \hyperref[topology-section-phantom]{Topologie}
\item \hyperref[sheaves-section-phantom]{Faisceaux sur les espaces}
\item \hyperref[sites-section-phantom]{Sites et faisceaux}
\item \hyperref[stacks-section-phantom]{Champs}
\item \hyperref[fields-section-phantom]{Corps}
\item \hyperref[algebra-section-phantom]{Algèbre commutative}
\item \hyperref[brauer-section-phantom]{Groupes de Brauer}
\item \hyperref[homology-section-phantom]{Algèbre homologique}
\item \hyperref[derived-section-phantom]{Catégories dérivées}
\item \hyperref[simplicial-section-phantom]{Méthodes simpliciales}
\item \hyperref[more-algebra-section-phantom]{Compléments d'algèbre}
\item \hyperref[smoothing-section-phantom]{Lissage des morphismes d'anneaux}
\item \hyperref[modules-section-phantom]{Faisceaux de modules}
\item \hyperref[sites-modules-section-phantom]{Modules sur les sites}
\item \hyperref[injectives-section-phantom]{Objets injectifs}
\item \hyperref[cohomology-section-phantom]{Cohomologie des faisceaux}
\item \hyperref[sites-cohomology-section-phantom]{Cohomologie sur les sites}
\item \hyperref[dga-section-phantom]{Algèbre différentielle graduée}
\item \hyperref[dpa-section-phantom]{Algèbre à puissances divisées}
\item \hyperref[sdga-section-phantom]{Faisceaux différentiels gradués}
\item \hyperref[hypercovering-section-phantom]{Hyperrecouvrements}
\end{enumerate}
Schémas
\begin{enumerate}
\setcounter{enumi}{25}
\item \hyperref[schemes-section-phantom]{Schémas}
\item \hyperref[constructions-section-phantom]{Constructions de schémas}
\item \hyperref[properties-section-phantom]{Propriétés des schémas}
\item \hyperref[morphisms-section-phantom]{Morphismes de schémas}
\item \hyperref[coherent-section-phantom]{Cohomologie des schémas}
\item \hyperref[divisors-section-phantom]{Diviseurs}
\item \hyperref[limits-section-phantom]{Limites de schémas}
\item \hyperref[varieties-section-phantom]{Variétés}
\item \hyperref[topologies-section-phantom]{Topologies sur les schémas}
\item \hyperref[descent-section-phantom]{Descente}
\item \hyperref[perfect-section-phantom]{Catégories dérivées de schémas}
\item \hyperref[more-morphisms-section-phantom]{Compléments sur les morphismes}
\item \hyperref[flat-section-phantom]{Compléments sur la platitude}
\item \hyperref[groupoids-section-phantom]{Schémas en groupoïdes}
\item \hyperref[more-groupoids-section-phantom]{Compléments sur les schémas en groupoïdes}
\item \hyperref[etale-section-phantom]{Morphismes étales de schémas}
\end{enumerate}
Thèmes de théorie des schémas
\begin{enumerate}
\setcounter{enumi}{41}
\item \hyperref[chow-section-phantom]{Homologie de Chow}
\item \hyperref[intersection-section-phantom]{Théorie de l'intersection}
\item \hyperref[pic-section-phantom]{Schémas de Picard des courbes}
\item \hyperref[weil-section-phantom]{Théories cohomologiques de Weil}
\item \hyperref[adequate-section-phantom]{Modules adéquats}
\item \hyperref[dualizing-section-phantom]{Complexes dualisants}
\item \hyperref[duality-section-phantom]{Dualité pour les schémas}
\item \hyperref[discriminant-section-phantom]{Discriminants et différentes}
\item \hyperref[derham-section-phantom]{Cohomologie de de Rham}
\item \hyperref[local-cohomology-section-phantom]{Cohomologie locale}
\item \hyperref[algebraization-section-phantom]{Géométrie algébrique et formelle}
\item \hyperref[curves-section-phantom]{Courbes algébriques}
\item \hyperref[resolve-section-phantom]{Résolution des surfaces}
\item \hyperref[models-section-phantom]{Réduction semi-stable}
\item \hyperref[functors-section-phantom]{Foncteurs et morphismes}
\item \hyperref[equiv-section-phantom]{Catégories dérivées de variétés}
\item \hyperref[pione-section-phantom]{Groupes fondamentaux de schémas}
\item \hyperref[etale-cohomology-section-phantom]{Cohomologie étale}
\item \hyperref[crystalline-section-phantom]{Cohomologie cristalline}
\item \hyperref[proetale-section-phantom]{Cohomologie pro-étale}
\item \hyperref[relative-cycles-section-phantom]{Cycles relatifs}
\item \hyperref[more-etale-section-phantom]{Compléments de cohomologie étale}
\item \hyperref[trace-section-phantom]{La formule des traces}
\end{enumerate}
Espaces algébriques
\begin{enumerate}
\setcounter{enumi}{64}
\item \hyperref[spaces-section-phantom]{Espaces algébriques}
\item \hyperref[spaces-properties-section-phantom]{Propriétés des espaces algébriques}
\item \hyperref[spaces-morphisms-section-phantom]{Morphismes d'espaces algébriques}
\item \hyperref[decent-spaces-section-phantom]{Espaces algébriques décents}
\item \hyperref[spaces-cohomology-section-phantom]{Cohomologie des espaces algébriques}
\item \hyperref[spaces-limits-section-phantom]{Limites d'espaces algébriques}
\item \hyperref[spaces-divisors-section-phantom]{Diviseurs sur les espaces algébriques}
\item \hyperref[spaces-over-fields-section-phantom]{Espaces algébriques sur des corps}
\item \hyperref[spaces-topologies-section-phantom]{Topologies sur les espaces algébriques}
\item \hyperref[spaces-descent-section-phantom]{Descente et espaces algébriques}
\item \hyperref[spaces-perfect-section-phantom]{Catégories dérivées d'espaces}
\item \hyperref[spaces-more-morphisms-section-phantom]{Compléments sur les morphismes d'espaces}
\item \hyperref[spaces-flat-section-phantom]{Platitude sur les espaces algébriques}
\item \hyperref[spaces-groupoids-section-phantom]{Groupoïdes dans les espaces algébriques}
\item \hyperref[spaces-more-groupoids-section-phantom]{Compléments sur les groupoïdes dans les espaces}
\item \hyperref[bootstrap-section-phantom]{Amorçage}
\item \hyperref[spaces-pushouts-section-phantom]{Sommes amalgamées d'espaces algébriques}
\end{enumerate}
Thèmes de géométrie
\begin{enumerate}
\setcounter{enumi}{81}
\item \hyperref[spaces-chow-section-phantom]{Groupes de Chow des espaces}
\item \hyperref[groupoids-quotients-section-phantom]{Quotients de groupoïdes}
\item \hyperref[spaces-more-cohomology-section-phantom]{Compléments sur la cohomologie des espaces}
\item \hyperref[spaces-simplicial-section-phantom]{Espaces simpliciaux}
\item \hyperref[spaces-duality-section-phantom]{Dualité pour les espaces}
\item \hyperref[formal-spaces-section-phantom]{Espaces algébriques formels}
\item \hyperref[restricted-section-phantom]{Algébrisation des espaces formels}
\item \hyperref[spaces-resolve-section-phantom]{Résolution des surfaces revisitée}
\end{enumerate}
Théorie des déformations
\begin{enumerate}
\setcounter{enumi}{89}
\item \hyperref[formal-defos-section-phantom]{Théorie formelle des déformations}
\item \hyperref[defos-section-phantom]{Théorie des déformations}
\item \hyperref[cotangent-section-phantom]{Le complexe cotangent}
\item \hyperref[examples-defos-section-phantom]{Problèmes de déformation}
\end{enumerate}
Champs algébriques
\begin{enumerate}
\setcounter{enumi}{93}
\item \hyperref[algebraic-section-phantom]{Champs algébriques}
\item \hyperref[examples-stacks-section-phantom]{Exemples de champs}
\item \hyperref[stacks-sheaves-section-phantom]{Faisceaux sur les champs algébriques}
\item \hyperref[criteria-section-phantom]{Critères de représentabilité}
\item \hyperref[artin-section-phantom]{Axiomes d'Artin}
\item \hyperref[quot-section-phantom]{Espaces de Quot et de Hilbert}
\item \hyperref[stacks-properties-section-phantom]{Propriétés des champs algébriques}
\item \hyperref[stacks-morphisms-section-phantom]{Morphismes de champs algébriques}
\item \hyperref[stacks-limits-section-phantom]{Limites de champs algébriques}
\item \hyperref[stacks-cohomology-section-phantom]{Cohomologie des champs algébriques}
\item \hyperref[stacks-perfect-section-phantom]{Catégories dérivées de champs}
\item \hyperref[stacks-introduction-section-phantom]{Introduction aux champs algébriques}
\item \hyperref[stacks-more-morphisms-section-phantom]{Compléments sur les morphismes de champs}
\item \hyperref[stacks-geometry-section-phantom]{La géométrie des champs}
\end{enumerate}
Thèmes de théorie des espaces de modules
\begin{enumerate}
\setcounter{enumi}{107}
\item \hyperref[moduli-section-phantom]{Champs de modules}
\item \hyperref[moduli-curves-section-phantom]{Espaces de modules de courbes}
\end{enumerate}
Divers
\begin{enumerate}
\setcounter{enumi}{109}
\item \hyperref[examples-section-phantom]{Exemples}
\item \hyperref[exercises-section-phantom]{Exercices}
\item \hyperref[guide-section-phantom]{Guide bibliographique}
\item \hyperref[desirables-section-phantom]{Desiderata}
\item \hyperref[coding-section-phantom]{Style de programmation}
\item \hyperref[obsolete-section-phantom]{Obsolète}
\item \hyperref[fdl-section-phantom]{Licence de documentation libre GNU}
\item \hyperref[index-section-phantom]{Index généré automatiquement}
\end{enumerate}
\end{multicols}


\bibliography{my}
\bibliographystyle{amsalpha}

\end{document}
